\documentclass[a4paper,reqno]{amsart}

\textheight 220mm
\textwidth 160mm
%\textheight 230mm
%\textwidth 150mm
\hoffset -16mm
%\voffset -16mm

\usepackage{amssymb}
\usepackage{amstext}
\usepackage{amsmath}
\usepackage{amscd}
\usepackage{amsthm}
\usepackage{amsfonts}
%\usepackage{txfonts}
\usepackage{enumerate}
\usepackage{graphicx}
\usepackage{latexsym}
%\usepackage{showkeys}
\usepackage[all]{xy}
\input xy
\xyoption{all}
\usepackage[usenames]{color}

\newtheorem{theorem}{Theorem}[section]
\newtheorem{acknowledgement}[theorem]{Acknowledgement}
\newtheorem{corollary}[theorem]{Corollary}
\newtheorem{lemma}[theorem]{Lemma}
\newtheorem{proposition}[theorem]{Proposition}

\theoremstyle{definition}
\newtheorem{definition}[theorem]{Definition}
\newtheorem{remark}[theorem]{Remark}
\newtheorem{example}[theorem]{Example}
\newtheorem{question}[theorem]{Question}

\newcommand{\Mod}[1]{{\rm{Mod}}\text{-}#1} %Mod-
\newcommand{\PP}{{\mathcal P}}
\newcommand{\CC}{{\mathcal C}}
\newcommand{\DD}{{\mathcal D}}
\newcommand{\SSS}{{\mathcal S}}
\newcommand{\UU}{{\mathcal U}}
\newcommand{\VV}{{\mathcal V}}
\newcommand{\TT}{{\mathcal T}}
\newcommand{\MM}{{\mathcal M}}
\newcommand{\NN}{{\mathcal N}}
\newcommand{\KK}{{\mathcal K}}
\newcommand{\LL}{{\mathcal L}}
\newcommand{\XX}{{\mathcal X}}
\newcommand{\YY}{{\mathcal Y}}
\newcommand{\Z}{\mathbb{Z}}
\newcommand{\ind}{\mathsf{ind}}
\newcommand{\add}{\mathsf{add}}
\newcommand{\Add}{\mathsf{Add}}
\newcommand{\summand}{\mathsf{smd}}
\newcommand{\thick}{\mathsf{thick}}
\newcommand{\Hom}{\operatorname{Hom}\nolimits}
\newcommand{\silt}{\operatorname{silt}\nolimits}

%%%%%%%%%%%%%%%%%%Command for Aihara Section%%%%%%%%%%%%%%%%%%%%%%%%%%%%%%%% 
\newcommand{\modu}[1]{{\rm{mod}}\text{-}#1} %mod-
\newcommand{\proj}[1]{{\rm{proj}}\text{-}#1} %proj-
\newcommand{\inj}[1]{{\rm{inj}}\text{-}#1} %inj-

\newcommand{\sta}[1]{{\underline{{\rm{mod}}}}\text{-}#1} %stable module cat
\newcommand{\ista}[1]{{\overline{{\rm{mod}}}}\text{-}#1} %injectively stable module cat

\newcommand{\Modu}[1]{{\rm{Mod}}\text{-}#1} %Mod- 
\newcommand{\Proj}[1]{{\rm{Proj}}\text{-}#1} %Proj-
\newcommand{\Inj}[1]{{\rm{Inj}}\text{-}#1} %Inj-

\newcommand{\homo}[3]{{\rm{Hom}}_{#1}(#2, #3)} %homo
\newcommand{\ext}[4]{{\rm{Ext}}_{#1}^{#2}(#3, #4)} %ext
\newcommand{\tor}[4]{{\rm{Tor}}^{#1}_{#2}(#3, #4)} %tor
\newcommand{\ehomo}[2]{{\rm{End}}_{#1}(#2)} %end
\newcommand{\shomo}[3]{\underline{{\rm{Hom}}}_{#1}(#2, #3)} %stable homo
\newcommand{\ishomo}[3]{{\overline{{\rm{Hom}}}}_{#1}(#2, #3)} % injectively stable homo

\newcommand{\syz}[2]{\Omega ^{#1}(#2)} %syzygy

\newcommand{\length}[1]{\mathsf{length}(#1)} %length 

\newcommand{\CCC}{\mathcal{C}} %C
\newcommand{\DDD}{\mathcal{D}} %D
\newcommand{\EEE}{\mathcal{E}} %E

\newcommand{\ann}[1]{\mathsf{ann}#1} %ann

\newcommand{\fl}[1]{{\rm{f.l}}\text{-}#1} %f.l-
%%%%%%%%%%%%%%%%%%%%%%%%%%%%%%%%%%%%%%%%%%%%%%%%%%%%%%%%%%%%%%%%%%%%%%%%%%%%%

%%%%%%%%%%%%%%%%%%Command for New Section%%%%%%%%%%%%%%%%%%%%%%%%%%%%%%%%%%%%
\newtheorem{assumption}[theorem]{Assumption}
\newcommand{\exc}{\operatorname{exp}\nolimits}
\newcommand{\K}{\mathsf{K}}
\newcommand{\D}{\mathsf{D}}
\newcommand{\End}{\operatorname{End}\nolimits}
\newcommand{\Ext}{\operatorname{Ext}\nolimits}
\renewcommand{\AA}{{\mathcal A}}
\renewcommand{\mod}{\operatorname{mod}\nolimits}
\newcommand{\RHom}{\mathbf{R}\operatorname{Hom}\nolimits}
%%%%%%%%%%%%%%%%%%%%%%%%%%%%%%%%%%%%%%%%%%%%%%%%%%%%%%%%%%%%%%%%%

%\usepackage{pstricks}
%\usepackage{lscape}
%\usepackage{comment}

%\newcommand{\new}[1]{{\blue #1}}
%\newcommand{\old}[1]{{\red #1}}
%\renewcommand{\comment}[1]{{\green #1}}

\begin{document}

\title{Silting mutation in triangulated categories}
\author{Takuma Aihara and Osamu Iyama}
%\date{\today}
\address{T. Aihara: Division of Mathematical Science and Physics, 
Graduate School of Science and Technology, Chiba University, Yayoi-cho, Chiba 263-8522, Japan}
\email{taihara@math.s.chiba-u.ac.jp}
\address{O. Iyama: Graduate School of Mathematics, Nagoya University, Chikusa-ku, Nagoya,
464-8602 Japan}
\email{iyama@math.nagoya-u.ac.jp}
\thanks{2010 {\em Mathematics Subject Classification.} 16E30, 18E30}

\begin{abstract}
In representation theory of algebras the notion of `mutation' often plays important roles,
and two cases are well known, i.e. `cluster tilting mutation' and `exceptional mutation'.
In this paper we focus on `tilting mutation', which has a disadvantage that it is often 
impossible, i.e. some of summands of a tilting object can not be replaced to get a new tilting 
object. The aim of this paper is to take away this disadvantage by introducing `silting mutation'
for silting objects as a generalization of `tilting mutation'. We shall develop a basic
theory of silting mutation.
In particular, we introduce a partial order on the set of silting objects and establish the relationship with `silting mutation' by generalizing the theory of
Riedtmann-Schofield and Happel-Unger. We show that iterated silting mutation act transitively on the set of silting objects for local, hereditary or canonical algebras.
Finally we give a bijection between silting subcategories and certain t-structures.
\end{abstract}
\maketitle
\tableofcontents

%%%%%%%%%%%%%%%%%%%%%%%%%%%%%%%%%%%%%%%%%%%%%%%%%%%%%%%%%%%%%%%%%%%%%%%%%%%%%%%%%%%%%%%%%%%%%%%%%%%%%%%%%%%
%%%%%%%%%%%%%%%%%%%%%%%%%%%%%%%%%%%%%%%%%%%%%%%%%%%%%%%%%%%%%%%%%%%%%%%%%%%%%%%%%%%%%%%%%%%%%%%%%%%%%%%%%%%

\section{Introduction}

In representation theory of algebras the notion of `mutation' often plays important roles. 
Mutation is an operation for a certain class of objects (e.g. cluster tilting objects)
in a fixed category to construct a new object from a given one by replacing
a summand.
Two important cases are well-known: One is `cluster tilting mutation' \cite{BMRRT,IY,IW}
for cluster tilting objects which is applied for categorification of
Fomin-Zelevinsky cluster algebras,
and the other is `exceptional mutation' for exceptional sequences 
which is used to study the structure of derived categories of algebraic varieties \cite{GR,C}.

From Morita theoretic viewpoint \cite{Ric1}, tilting objects are the most important
class of objects, and `tilting mutation' for tilting objects has
been studied by several authors. `Tilting mutation' has the origin in
BGP (=Bernstein-Gelfand-Ponomarev) reflection functor \cite{BGP} in quiver
representation theory, and reflection functors are understood as a special class
of APR (=Auslander-Platzeck-Reiten) tilting modules \cite{APR},
which are `tilting mutations' obtained by replacing a simple summand of
the tilting $A$-modules $A$. In 1991 general `tilting mutation' for tilting modules was
introduced by Riedtmann-Schofield \cite{RS} in their study of combinatorial aspects
of tilting theory. It has been shown by Happel-Unger \cite{HU} that `tilting mutation'
is closely related to the partial order of tilting modules given by the inclusion
relation of the associated t-structures, and this is a big advantage of `tilting mutation'
which `cluster tilting mutation' and `exceptional mutation' do not have.
Their beautiful theory of `tilting mutation' has a lot of important applications, especially in the study of
cluster categories and preprojective algebras \cite{BMRRT,IR,BIRS,BIKR,SY}.
Another important source of `tilting mutation' is modular representation theory, 
where `tilting mutation' for tilting complexes was introduced by 
Okuyama and Rickard \cite{O,Ric2,HK}, and has played an important role 
in the study of Brou\'e's abelian defect group conjecture.

It is remarkable that `tilting mutation' has a big disadvantage
that it is often impossible, i.e. some of summands of a tilting
object can not be replaced to get a new tilting object. So it is usual that we can not
get sufficiently many tilting objects by iterated `tilting mutation'.
For example the tilting $A$-module $DA$ usually can not be obtained by
iterated `tilting mutation' of $A$ even for the case of hereditary algebras $A$.
The aim of this paper is to take away this disadvantage by introducing 
`\emph{silting mutation}' for \emph{silting objects}. In this context
`tilting mutation' should be understood as a special case of `silting mutation'.
Actually classical notion of APR and BB tilting modules and Okuyama-Rickard complexes
can be understood as special cases of `silting mutation'  (Theorems \ref{prop:APR is APR} and \ref{prop:1_right_mu}).
Silting objects are generalization of tilting objects,
and sometimes appeared in representation theory mainly in the study of
t-structures, e.g. Keller-Vossieck \cite{KV}, 
Hoshino-Kato-Miyachi \cite{HKM}, Assem-Salorio-Trepode \cite{AST} and Wei \cite{W}.
A point is that `silting mutation' is always possible in the sense that
any summand of a silting object always can be replaced to get a new
silting object (Theorem \ref{thm:mutation is silting}).  
Hence it is natural to hope that `silting mutation' gives us sufficiently many
silting objects in triangulated categories. We pose the following question
(Question \ref{transitive question2}),
where the corresponding property for tilting objects is usually not satisfied.

\begin{question}\label{transitive question}
Let $A$ be a finite dimensional algebra over a field.
When does $A$ satisfy the following property (T) (respectively, (T$'$))?

(T) (respectively, (T$'$)) The action of iterated irreducible `silting mutation'
(respectively, iterated `silting mutation') on the set of basic silting objects in $\K^{\rm b}(\proj{A})$ is transitive.
\end{question}

We shall show the following partial answers in Corollary \ref{corollary indecomposable case} and Theorem \ref{hereditary case}.

\begin{theorem}
If $A$ is either a local algebra, a hereditary algebra or a canonical algebra, then (T) is satisfied.
\end{theorem}

It will be shown in \cite{A2} that (T) is satisfied also for representation-finite symmetric algebras.
On the other hand, there is a symmetric algebra such that (T) is not satisfied \cite{AGI}.
We do not know any algebra such that (T$'$) is not satisfied.

As a basic tool of the study of silting objects, we shall introduce
a partial order on silting objects (Theorem \ref{partial order}) and establish the relationship with
`silting mutation' (Theorem \ref{two quivers}) by generalizing the theory of Riedtmann-Schofield and
Happel-Unger for tilting modules.
In particular, Question \ref{transitive question} is equivalent to the following question.

\begin{question}
Let $A$ be a finite dimensional algebra over a field.
When is the Hasse quiver of the partially ordered set of basic silting objects in $\K^{\rm b}(\proj{A})$ connected?
\end{question}

%asking whether the Hasse quiver of the partially ordered set $\silt\TT$ is connected or not.


The comparison of three kinds of mutation is explained by the following table.
\[\begin{array}{|c||c|c|c|}
\hline
&\mbox{always possible}&\mbox{partial order exists}&\mbox{transitive for hereditary case}\\ \hline\hline
\mbox{cluster tilting mutation}&\mbox{Yes}&\mbox{No}&\mbox{Yes}\\ \hline
\mbox{silting mutation}&\mbox{Yes}&\mbox{Yes}&\mbox{Yes}\\
\cup&&&\\
\mbox{tilting mutation}&\mbox{No}&\mbox{Yes}&\mbox{No}\\ \hline
\mbox{exceptional mutation}&\mbox{Yes}&\mbox{No}&\mbox{Yes}\\ \hline
\end{array}\]

In our paper we study basic properties of silting objects/subcategories.
We show that non-isomorphic indecomposable objects in any silting subcategory form a basis of the Grothendieck group of the triangulated category
(Theorem \ref{grothendieck}).
In particular all basic silting objects have the same number of indecomposable summands (Corollary \ref{thm:number of indecomposable direct summands}).
We also introduce `silting reduction' (Theorem \ref{silting reduction}), which gives a bijection between silting subcategories containing a certain fixed subcategory
and silting subcategories in a certain quotient triangulated category.

In section \ref{The correspondence between silting subcategories and t-structures}
we study in detail the relationship between silting subcategories and t-structures under the assumption that the triangulated categories have arbitrary coproducts,
and improve some of pioneering results of Hoshino-Kato-Miyachi \cite{HKM}.
The advantage of this setting is that each set of compact objects gives rise to a torsion pair (Theorem \ref{torsion pair}),
which is not the case for the setting of section \ref{section: silting subcategories}.
This basic result seems to be new and to have independent interest.
As an application, we establish a one-to-one correspondence between silting subcategories
and t-structures satisfying certain conditions (Theorem \ref{correspondence}).
%Then we study the hearts of certain t-structures including those associated with silting subcategories.
%We show that the hearts contain generating subcategories consisting of compact projective objects,
%and in particular we realize hearts as certain functor categories (Theorem \ref{category P}, Corollary \ref{realization1}).

We notice that a different generalization of `tilting mutation' was given in \cite{IO,HX}.
Also some aspects of mutation were recently discussed in \cite{BRT,L,KY}.

Parts of results in this paper were presented in Trondheim (March 2009),
Nagoya (June 2009), Matsumoto (October 2009), Shizuoka (November 2009) and Tokyo (August 2010).

\medskip
\noindent
{\bf Notations }
Let $\TT$ be an additive category.
For morphisms $f:X\to Y$ and $g:Y\to Z$ in a category, we denote by $gf:X\to Z$ the composition.
We say that a morphism $f:X\to Y$ is \emph{right minimal} if any morphism $g:X\to X$
satisfying $fg=f$ is an isomorphism. Dually we define a \emph{left minimal} morphism.

For a collection $\XX$ of objects in $\TT$,
we denote by $\add\XX$ (respectively, $\Add\XX$) the smallest full subcategory of $\TT$
which is closed under finite (respectively, arbitrary) coproducts, summands and isomorphisms
and contains $\XX$.
We denote by $\summand\XX$ the smallest full subcategory of $\TT$ which is closed under summands and contains $\XX$.
When we say that $\XX$ is a \emph{subcategory} of $\TT$, we always assume that
$\XX$ is full and satisfies $\XX=\add\XX$.

Let $\XX$ be a subcategory of $\TT$.
We say that a morphism $f:X\to Y$ is a \emph{right $\XX$-approximation} of $Y$ 
if $X\in \XX$ and $\homo{\TT}{X}{f}$ is surjective for any $X\in \XX$. 
We say that $\XX$ is \emph{contravariantly finite} if any object in $\TT$ has a right $\XX$-approximation. 
Dually, we define a \emph{left $\XX$-approximation} and a \emph{covariantly finite subcategory}. 
We say that $\XX$ is \emph{functorially finite} if it is contravariantly and covariantly finite. 

When $\TT$ is a triangulated category, we denote by $\thick\XX$ the smallest thick subcategory of $\TT$ containing $\XX$.
For collections $\XX$ and $\YY$ of objects in $\TT$, we denote by
$\XX*\YY$ the collection of objects $Z\in\TT$ appearing in a triangle
$X\to Z\to Y\to X[1]$ with $X\in\XX$ and $Y\in\YY$.
We set
\[\begin{array}{ccccc}
\XX^\perp   &=& \XX^{\perp_{\TT}} &:=& \{T\in \TT\ |\ \Hom_{\TT}(\XX,T)=0\},\\
{}^\perp\XX &=& {}^{\perp_{\TT}}\XX &:=& \{T\in \TT\ |\ \Hom_{\TT}(T,\XX)=0\}
\end{array}\]

For an additive category $\AA$, we denote by $\K^{\rm b}(\AA)$ the homotopy
category of bounded complexes over $\AA$.
For an abelian category $\AA$, we denote by $\D^{\rm b}(\AA)$ the bounded
derived category of $\AA$.

For a ring $A$, we denote by $\Mod A$ (respectively, $\modu A$) the category
of all (respectively, finitely generated) right $A$-modules,
by $\proj A$ the category of finitely generated projective $A$-modules.
%and by $\fl A$ the category of finite length $A$-modules.
We denote by $J_A$ the Jacobson radical of $A$.
When $A$ is a finite dimensional algebra over a field $k$,
we denote by $D:=\Hom_k(-,k):\modu A\leftrightarrow\modu A^{\rm op}$ the $k$-duality.

\begin{definition}
Let $\TT$ be a triangulated category and $(\XX,\YY)$ a pair of subcategories of $\TT$.
\begin{itemize}
\item We say that $(\XX,\YY)$ is a \emph{torsion pair} 
if $\Hom_{\TT}(\XX,\YY)=0$ and $\XX*\YY=\TT$.
\item We say that $(\XX,\YY)$ is a \emph{t-structure} \cite{BBD}
if $(\XX[1],\YY)$ is a torsion pair and $\XX[1]\subset\XX$.
In this case $\XX\cap\YY$ is called the \emph{heart}.
\item We say that $(\XX,\YY)$ is a \emph{co-t-structure} \cite{P1} 
if $(\XX[-1],\YY)$ is a torsion pair and $\XX\subset\XX[1]$.
In this case $\XX\cap\YY$ is called the \emph{coheart}.
\item We say that a torsion pair $(\XX,\YY)$ is a \emph{stable t-structure} 
\cite{M} if $\XX=\XX[1]$.
\end{itemize}
\end{definition}


\medskip
\noindent
{\bf Acknowledgements }
The second author would like to thank Jiaqun Wei for stimulating discussion.
The authors would like to thank David Pauksztello, Apostolos Beligiannis and Sonia Trepode for useful information.

%%%%%%%%%%%%%%%%%%%%%%%%%%%%%%%%%%%%%%%%%%%%%%%%%%%%%%%%%%%%%%%%%%%%%%%%%%%%%%%%%%%%%%%%%%%%%%%%%%%%%%%%%%%
%%%%%%%%%%%%%%%%%%%%%%%%%%%%%%%%%%%%%%%%%%%%%%%%%%%%%%%%%%%%%%%%%%%%%%%%%%%%%%%%%%%%%%%%%%%%%%%%%%%%%%%%%%%

\section{Silting subcategories}\label{section: silting subcategories}

Let $\TT$ be a triangulated category. We do not assume anything else on $\TT$ unless otherwise stated.

Throughout this paper we write the vanishing condition $\Hom_{\TT}(X,Y[i])=0$ for any $i\in\Z$ by
\[\Hom_{\TT}(X,Y[\Z])=0\]
simply.  Similarly we often use
\[\Hom_{\TT}(X,Y[>0])=0,\ \Hom_{\TT}(X,Y[<0])=0,\ \Hom_{\TT}(X,Y[\neq0])=0\ \mbox{ and so on.}\]

\subsection{Definition and basic properties}

In this subsection we introduce silting subcategories/objects of triangulated categories
and study their basic properties.

\begin{definition}\label{definition of silting}
Let $\MM$ be a subcategory of a triangulated category $\TT$.
\begin{itemize}
\item[(a)] We say that $\MM$ is \emph{silting} if $\Hom_{\TT}(\MM, \MM[>0])=0$ and $\TT=\thick\MM$. 
We denote by $\silt\TT$ the collection of silting subcategories in $\TT$.
\item[(b)] We say that $\MM$ is \emph{tilting} if $\Hom_{\TT}(\MM, \MM[\neq0])=0$ and $\TT=\thick\MM$.
\item[(c)] We say that an object $M\in\TT$ is \emph{silting} (respectively, \emph{tilting}) if so is $\add M$.
\end{itemize}
\end{definition}

The following examples give typical classes of tilting/silting objects.

\begin{example}
\begin{itemize}
\item[(a)] Let $A$ be a ring.
Then $A$ regarded as a stalk complex is a tilting object in $\K^{\rm b}(\proj A)$.
More generally, any tilting complex of $A$ is a tilting object in $\K^{\rm b}(\proj A)$.
\item[(b)] Let $A$ be a differential graded ring and let $\D(A)$ be the derived category of $A$.
If $H^i(A)=0$ for any $i>0$, then the thick subcategory $\thick(A)$ of $\D(A)$ generated by $A$
has a silting object $A$.
\end{itemize}
\end{example}

The above example (a) is standard among triangulated categories with tilting objects in the following sense:

\begin{proposition}\cite{K}
Let $\TT$ be an algebraic triangulated category.
If $\TT$ has a tilting object $M$, then $\TT$ is triangle equivalent
to $\K^{\rm b}(\proj\End_{\TT}(M))$.
\end{proposition}

%No similar characterization of (algebraic) triangulated categories containing silting objects is known.

On the other hand, we have the following necessary conditions for existence of silting/tilting subcategories.

\begin{proposition}\label{silting vanishing}
Let $\TT$ be a triangulated category.
If $\TT$ contains a silting (respectively, tilting) subcategory,
then for any $X,Y\in\TT$ we have $\Hom_{\TT}(X,Y[i])=0$ for $i\gg0$ (respectively, $|i|\gg0$).
\end{proposition}

\begin{proof}
We only show the statement for silting subcategories.
Let $\UU:=\{X\in\TT\ |\ \Hom_{\TT}(X,M[\gg0])=0$ for any $M\in\MM\}$.
Then $\UU$ is a thick subcategory of $\TT$ containing $\MM$.
Thus we have $\UU=\TT$ since $\thick\MM=\TT$.
For any $X\in\TT$, let $\VV_X:=\{Y\in\TT\ |\ \Hom_{\TT}(X,Y[\gg0])=0\}$.
Then we have $\MM\subset\VV_X$ by the above argument.
Since $\VV_X$ is a thick subcategory of $\TT$ containing $\MM$, we have $\VV_X=\TT$ again.
Thus the assertion holds.
\end{proof}

Immediately we have the following observation.

\begin{example}
Let $A$ be a finite dimensional algebra over a field $k$.
\begin{itemize}
\item[(a)] $\D^{\rm b}(\modu A)$ contains silting subcategories if and only if ${\rm gl.dim} A<\infty$.
\item[(b)] Assume that $A$ is self-injective.
Then the stable module category $\sta{A}$ of $\modu{A}$ contains silting subcategories if and only if $A$ is semisimple.
\end{itemize}
\end{example}

\begin{proof}
We only have to show `only if' part.
Both for the cases (a) and (b), we have $\Ext_A^i(A/J_A,A/J_A)\neq0$ for any $i>0$.
Thus the assertions follow from Proposition \ref{silting vanishing}.
\end{proof}

It is often the case that triangulated categories have many silting subcategories which are not tilting.
But there are important triangulated categories such that all silting subcategories are tilting.

%the fact that $\K^{\rm{b}}(\proj{A})$ is a $0$-Calabi-Yau triangulated category in the following sense:

\begin{definition}
Let $\ell$ be an integer. We call a $k$-linear $\Hom$-finite triangulated category $\TT$ \emph{$\ell$-Calabi-Yau}
if there is a bifunctorial isomorphism $\homo{\TT}{-}{?}\simeq D\homo{\TT}{?}{-[\ell]}$. 
\end{definition}

%For example if $A$ is a symmetric algebra, then 
%and the stable module category $\sta{A}$ is $(-1)$-Calabi-Yau. 
One can easily check the following statements.

\begin{lemma}\label{prop:CY}
Let $\ell$ be an integer and $\TT$ an $\ell$-Calabi-Yau triangulated category.
If $\TT\neq0$, then the following assertions hold:
\begin{enumerate}[$(1)$]
\item If $\ell=0$, then any silting subcategory of $\TT$ is tilting; 
\item If $\ell>0$, then there exist no silting subcategories of $\TT$;
\item If $\ell<0$, then there exist no tilting subcategories of $\TT$.
\end{enumerate}
\end{lemma}

\begin{proof}
Since any non-zero object $X\in\TT$ satisfies $\homo{\TT}{X}{X[\ell]}\simeq D\homo{\TT}{X}{X}\neq0$,
the assertions (2) and (3) follow. We can prove (1) similarly.
\end{proof}

Immediately we have the following observation.

\begin{example}
Let $A$ be a finite dimensional symmetric algebra over a field.
Then any silting object in $\K^{\rm b}(\proj A)$ is tilting.
\end{example}

\begin{proof}
We know that $\K^{\rm b}(\proj A)$ is a 0-Calabi-Yau triangulated category by Auslander-Reiten duality
\cite{H}. Thus the assertion follows from Lemma \ref{prop:CY}.
\end{proof}

Notice that we do not know whether there exist $\ell$-Calabi-Yau triangulated categories
containing a silting subcategory and $\ell<0$.

We end this subsection with the following remark, where we say that a category
is \emph{skeletally small} if isomorphism classes of objects form a set.

\begin{remark}\label{T is small}
If a triangulated category $\TT$ has a skeletally small silting subcategory,
then $\TT$ is also skeletally small (e.g. Proposition \ref{generate2}).
\end{remark}


%%%%%%%%%%%%%%%%%%%%%%%%%%%%%%%%%%%%%%%%%%%%%%%%%%%%%%%%%%%%%%%%%%%%%%%%%%%%%%%%%%%%%%%%%%%%%%%%%%%%%%%%%%%
%%%%%%%%%%%%%%%%%%%%%%%%%%%%%%%%%%%%%%%%%%%%%%%%%%%%%%%%%%%%%%%%%%%%%%%%%%%%%%%%%%%%%%%%%%%%%%%%%%%%%%%%%%%


\subsection{Partial order on silting subcategories}

The aim of this subsection is to introduce a partial order on silting subcategories
as a generalization of the partial order on tilting modules introduced by
Riedtmann-Schofield \cite{RS} and Happel-Unger \cite{HU}.
Our main result in this subsection is Theorem \ref{partial order} below.

\begin{definition}
For $\MM,\NN\in\silt\TT$, we write $\MM\ge\NN$ if $\Hom_{\TT}(\MM, \NN[>0])=0$.
\end{definition}

\begin{theorem}\label{partial order}
$\ge$ gives a partial order on $\silt\TT$.
\end{theorem}

The following subcategory plays an important role in the proof of Theorem \ref{partial order}.

\begin{definition}
For any $\MM\in\silt\TT$, we define a subcategory of $\TT$ by
\[\TT_{\MM}^{\le0}:=\{X\in\TT\ |\ \Hom_{\TT}(\MM, X[>0])=0\}.\]
We put $\TT_{\MM}^{\le\ell}=\TT_{\MM}^{<\ell+1}:=\TT_{\MM}^{\le0}[-\ell]$ for $\ell\in\Z$.
\end{definition}

Clearly we have $\MM\subset\TT_{\MM}^{\le0}$ and $\TT_{\MM}^{<0}\subset\TT_{\MM}^{\le0}$.
For any non-zero object $X\in\TT$, we have $X[\gg0]\in\TT_{\MM}^{\le0}$ and $X[\ll0]\notin\TT_{\MM}^{\le0}$ 
by Proposition \ref{silting vanishing}.

Crucial results are the following.

\begin{proposition}\label{recover}
For any $\MM\in\silt\TT$, we have $\TT_{\MM}^{\le0}\cap{}^\perp(\TT_{\MM}^{<0})=\MM$.
\end{proposition}

\begin{proposition}\label{another condition}
Let $\MM,\NN\in\silt\TT$.
Then $\MM\ge\NN$ if and only if $\TT_{\MM}^{\le0}\supset\TT_{\NN}^{\le0}$.
\end{proposition}

Before proving these Propositions, we prove Theorem \ref{partial order} by using them.

(i) We shall show that $\MM\ge\NN$ and $\NN\ge\LL$ imply $\MM\ge\LL$.
By Proposition \ref{another condition}, we have $\TT_{\MM}^{\le0}\supset\TT_{\NN}^{\le0}\supset\TT_{\LL}^{\le0}$.
Thus $\TT_{\MM}^{\le0}\supset\TT_{\LL}^{\le0}$, and we have $\MM\ge\LL$ again by Proposition \ref{another condition}.

(ii) We shall show that $\MM\ge\NN$ and $\NN\ge\MM$ imply $\MM=\NN$.
By Proposition \ref{another condition}, we have $\TT_{\MM}^{\le0}=\TT_{\NN}^{\le0}$.
Thus we have ${}^\perp(\TT_{\MM}^{<0})={}^\perp(\TT_{\NN}^{<0})$.
By Proposition \ref{recover}, we have $\MM=\NN$.
\qed

\medskip
To prove Propositions \ref{recover} and \ref{another condition}, we need the following result.

\begin{lemma}\label{generate}
\begin{itemize}
\item[(a)] For any subcategory $\MM$ of $\TT$, we have
\[\thick\MM=\bigcup_{\ell>0,\ n_1,\cdots,n_\ell\in\Z}\summand
(\MM [n_1]*\cdots*\MM [n_{\ell}]).\]
\item[(b)] Moreover if $\Hom_{\TT}(\MM,\MM[>0])=0$, then we have
\[\thick\MM=\bigcup_{\ell\ge0}
\summand(\MM[-\ell]*\MM[1-\ell]*\cdots*\MM[\ell-1]*\MM[\ell]).\]
\end{itemize}
\end{lemma}

\begin{proof}
(a) Clearly the right hand side is contained in $\thick\MM$.
Since the right hand side clearly forms a thick subcategory of $\TT$, it contains $\thick\MM$.

(b) If $n\ge m$, then any triangle
\[M[n]\to X\to M'[m]\xrightarrow{f}M[n+1]\]
with $M,M'\in\MM$ satisfies $f=0$. This implies $X\simeq M[n]\oplus M'[m]$, so we have
$\MM[n]*\MM[m]\subset\MM[m]*\MM[n]$.
Thus we have reduced to the assertion from (a).
\end{proof}

The following observation is clear since we only consider subcategories $\MM$ satisfying $\MM=\add\MM$ by our notations.

\begin{lemma}\label{reduction}
Let $\MM$ and $\NN$ be subcategories of $\TT$ and $X\in\summand(\MM*\NN)$.
\begin{itemize}
\item[(a)] If $\Hom_{\TT}(\MM,X)=0$, then $X\in\NN$.
\item[(b)] If $\Hom_{\TT}(X,\NN)=0$, then $X\in\MM$.
\end{itemize}
\end{lemma}

Now we have the following description of $\TT_{\MM}^{\le0}$.

\begin{proposition}\label{generate2}
For any $\MM\in\silt\TT$, we have
\begin{eqnarray*}
\TT&=&\bigcup_{\ell\ge0}\summand(\MM[-\ell]*\MM[1-\ell]*\cdots*\MM[\ell-1]*\MM[\ell]),\\
\TT_{\MM}^{\le0}&=&\bigcup_{\ell\ge0}\summand(\MM*\MM[1]*\cdots*\MM[\ell]).
\end{eqnarray*}
%In particular we have $\TT=\bigcup_{i\ge0}\TT_{\MM}^{\le0}[-i]$.
\end{proposition}

\begin{proof}
The first equation is clear from Lemma \ref{generate}(b).

Fix any $X\in\TT_{\MM}^{\le0}$. We can take the smallest integer $n$ such that
\[X\in\summand(\MM[n]*\MM[n+1]\cdots*\MM[\ell])\]
for some $\ell\ge n$. By Lemma \ref{reduction}, the minimality of $n$ implies
$\Hom_{\TT}(\MM[n],X)\neq0$.
Since $X\in\TT_{\MM}^{\le0}$, we have $n\ge0$.
Thus we have $X\in\summand(\MM*\MM[1]*\cdots*\MM[\ell])$.
\end{proof}

Now we have the following important property of silting subcategories.

\begin{theorem}\label{no inclusion}
If $\MM,\NN\in\silt\TT$ satisfy $\MM\subset\NN$, then $\MM=\NN$.
\end{theorem}

\begin{proof}
Fix $X\in\NN$. Since $\NN$ is silting, we have $X\in\TT_{\MM}^{\le0}$.
By Proposition \ref{generate2}, we have
\[X\in\summand(\MM*\MM[1]*\cdots*\MM[\ell])\]
for some $\ell\ge0$.
Since $\Hom_{\TT}(X,\MM[1]*\cdots*\MM[\ell])=0$, we have $X\in\MM$
by Lemma \ref{reduction}.
\end{proof}

Now we give a proof of Proposition \ref{recover}.

Put $\NN:=\TT_{\MM}^{\le0}\cap{}^\perp(\TT_{\MM}^{<0})$.
We have $\MM\subset\TT_{\MM}^{\le0}$.
Since we have $\Hom_{\TT}(\MM,\TT_{\MM}^{<0})=0$, we have
$\MM\subset{}^\perp(\TT_{\MM}^{<0})$.
Thus we have $\MM\subset\NN$.

On the other hand, we have
\[\Hom_{\TT}(\NN,\NN[>0])\subset\Hom_{\TT}({}^\perp(\TT_{\MM}^{<0}),\TT_{\MM}^{\le0}[>0])=0.\]
Consequently, $\NN$ is also a silting subcategory of $\TT$.
By Theorem \ref{no inclusion}, we have $\MM=\NN$.
\qed

\medskip
Now we give a proof of Proposition \ref{another condition}.

First we assume $\TT_{\MM}^{\le0}\supset\TT_{\NN}^{\le0}$. Then we have $\NN\subset\TT_{\MM}^{\le0}$,
which implies $\Hom_{\TT}(\MM,\NN[>0])=0$.

Conversely we assume $\MM\ge\NN$. Then we have $\NN[\ge0]\subset\TT_{\MM}^{\le0}$.
By Proposition \ref{generate2}, we have
\[\TT_{\NN}^{\le0}=\bigcup_{\ell\ge0}\summand(\NN*\NN[1]*\cdots*\NN[\ell])
\subset\TT_{\MM}^{\le0}.\]
Thus we have completed the proof.
\qed

\medskip
Now we give the following property of the partial order.

\begin{proposition}\label{common summand}
If $\MM,\NN,\LL\in\silt\TT$ satisfy
$\MM\ge\LL\ge\NN$, then $\MM\cap\NN\subset\LL$.
\end{proposition}

\begin{proof}
By Proposition \ref{another condition}, we have $\TT_{\MM}^{\le0}\supset\TT_{\LL}^{\le0}\supset\TT_{\NN}^{\le0}$ and
${}^\perp(\TT_{\MM}^{<0})\subset{}^\perp(\TT_{\LL}^{<0})\subset{}^\perp(\TT_{\NN}^{<0})$.
By Proposition \ref{recover}, we have
\[\MM\cap\NN\subset{}^\perp(\TT_{\MM}^{\le0})\cap\TT_{\NN}^{\le0}
\subset{}^\perp(\TT_{\LL}^{\le0})\cap\TT_{\LL}^{\le0}=\LL.\]
\end{proof}

We end this subsection by the following observation.

\begin{proposition}\label{existence of additive generator}
If $\TT$ has a silting object, then any silting subcategory contains an additive generator.
\end{proposition}

\begin{proof}
Let $M$ be a silting object and $\NN$ be a silting subcategory.
By Lemma \ref{generate}, there exists $N_1,\ldots,N_\ell\in\NN$ and $n_1,\cdots,n_\ell\in\Z$ such that $M\in\summand(N_1[n_1]*\cdots*N_\ell[n_\ell])$.
Let $\NN':=\add(N_1\oplus\cdots\oplus N_\ell)$.
Since $M\in\thick\NN'$, we have $\TT=\thick\NN'$.
Thus $\NN'$ is a silting subcategory of $\TT$.
By Theorem \ref{no inclusion}, we have $\NN=\NN'$.
Thus $\NN$ has an additive generator.
\end{proof}

%%%%%%%%%%%%%%%%%%%%%%%%%%%%%%%%%%%%%%%%%%%%%%%%%%%%%%%%%%%%%%%%%%%%%%%%%%%%%%%%%%%%%%%%%%%%%%%%%%%%%%%%%%%
%%%%%%%%%%%%%%%%%%%%%%%%%%%%%%%%%%%%%%%%%%%%%%%%%%%%%%%%%%%%%%%%%%%%%%%%%%%%%%%%%%%%%%%%%%%%%%%%%%%%%%%%%%%

\subsection{Krull-Schmidt triangulated categories}

Let $\TT$ be a triangulated category.
In this subsection we always assume that $\TT$ is \emph{Krull-Schmidt}
in the sense that any object in $\TT$ is isomorphic to a finite coproduct
of objects whose endomorphism rings are local.
In this case such a coproduct is uniquely determined up to isomorphism.
We denote by $J_{\TT}$ the \emph{Jacobson radical} of $\TT$ \cite{ARS,ASS}.
For a subcategory $\MM$ of $\TT$, we denote by $\ind\MM$ the set
of isoclasses of indecomposable objects in $\MM$.

We say that an object $M\in\TT$ is \emph{basic} if $M$ is isomorphic to
a coproduct of indecomposable objects which are mutually non-isomorphic.
Since $\TT$ is Krull-Schmidt, we have a one-to-one correspondence
between the isomorphism classes of basic objects $M$ and subcategories $\MM$ of $\TT$
containing additive generators. It is given by $M\mapsto\MM=\add M$.

\begin{proposition}
Assume that $\TT$ has a silting object.
Then we can regard $\silt\TT$ as the set of isomorphism classes of basic silting objects in $\TT$.
\end{proposition}

\begin{proof}
$\silt\TT$ is a set by Remark \ref{T is small}.
By Proposition \ref{existence of additive generator}, any silting subcategory of $\TT$ is an additive closure of a silting object.
Thus we have the assertion.
\end{proof}

Thanks to Krull-Schmidt assumption, we have the following useful property (e.g. \cite[2.1, 2.3]{IY}),
where (b) and (c) is famous as a `Wakamatsu's Lemma'.

\begin{lemma}\label{direct summand}
Let $\MM$ be a subcategory of $\TT$.
Then the following statements hold.
\begin{itemize}
\item[(a)] If a subcategory $\NN$ of $\TT$ satisfies $\Hom_{\TT}(\MM,\NN)=0$, 
then $\MM*\NN$ is closed under summands.
\item[(b)] If $\MM$ is contravariantly finite and $\MM*\MM\subset\MM$, 
then $(\MM, \MM^{\perp})$ is a torsion pair.
\item[(c)] If $\MM$ is covariantly finite and $\MM*\MM\subset\MM$, 
then $({}^{\perp}\MM, \MM)$ is a torsion pair. 
\end{itemize}
\end{lemma}

Now we have the following equalities.

\begin{proposition}\label{generate3}
\begin{itemize}
\item[(a)] For any $\MM\in\silt\TT$, we have
\begin{eqnarray*}%\label{T is M}
\TT&=&\bigcup_{\ell\ge0}\MM[-\ell]*\MM[1-\ell]*\cdots*\MM[\ell-1]*\MM[\ell],\\ %\label{T< is M}
\TT_{\MM}^{\le0}&=&\bigcup_{\ell\ge0}\MM*\MM[1]*\cdots*\MM[\ell],\\ 
{}^{\perp}(\TT^{\leq 0}_{\MM})&=&\bigcup_{\ell>0}\MM[-\ell]*\MM[-\ell+1]*\cdots*\MM[-1].
\end{eqnarray*}
\item[(b)] $({}^{\perp}(\TT^{\leq 0}_{\MM}), \TT^{\leq 0}_{\MM})$ is a torsion pair, and so $({}^{\perp}(\TT^{<0}_{\MM}), \TT^{\leq 0}_{\MM})$ is a co-t-structure with the coheart $\MM$.
\end{itemize}
\end{proposition}

\begin{proof}
The first two equalities follow from Proposition \ref{generate2} and Lemma \ref{direct summand}(a).
By the second equality we have 
\[{}^{\perp}(\TT^{\leq 0}_{\MM})\supset\bigcup_{\ell>0}\MM[-\ell]*\MM[-\ell+1]*\cdots*\MM[-1].\]
Together with the first equality we have
\[(\bigcup_{\ell>0}\MM[-\ell]*\cdots*\MM[-1])*\TT^{\leq 0}_{\MM}=(\bigcup_{\ell>0}\MM[-\ell]*\cdots*\MM[-1])*(\bigcup_{\ell>0}\MM*\cdots\MM[\ell])=\TT.\]
Thus $(\bigcup_{\ell>0}\MM[-\ell]*\cdots*\MM[-1], \TT^{\leq 0}_{\MM})$ is a torsion pair, and hence we have
\[{}^{\perp}(\TT^{\leq 0}_{\MM})\subset\bigcup_{\ell>0}\MM[-\ell]*\MM[-\ell+1]*\cdots*\MM[-1].\]
Thus we get the third equality, and $({}^{\perp}(\TT^{\leq 0}_{\MM}), \TT^{\leq 0}_{\MM})$ is a torsion pair.
The coheart of the co-t-structure $({}^{\perp}(\TT^{<0}_{\MM}), \TT^{\leq 0}_{\MM})$ is $\MM$ by Proposition \ref{recover}.
\end{proof}

\begin{proposition}\label{resolution}
Let $\MM\in\silt\TT$. For any $N=N_0\in\TT_{\MM}^{\le0}$, we have triangles
\[\xymatrix@R=0.2cm{
N_1\ar[r]^{g_1}&M_0\ar[r]^{f_0}&N_0\ar[r]&N_1[1],\\
&\cdots,\\
N_\ell\ar[r]^{g_\ell}&M_{\ell-1}\ar[r]^{f_{\ell-1}}&N_{\ell-1}\ar[r]&N_\ell[1],\\
0\ar[r]^{g_{\ell+1}}&M_\ell\ar[r]^{f_\ell}&N_\ell\ar[r]&0,
}\]
for some $\ell\ge0$ such that $f_i$ is a minimal right $\MM$-approximation and
$g_{i+1}$ belongs to $J_{\TT}$ for any $0\le i\le \ell$.
\end{proposition}

\begin{proof}
Since $N_0\in\TT_{\MM}^{\le0}$, we have
\[N_0\in\MM*\MM[1]*\cdots*\MM[\ell]\]
for some $\ell\ge0$ by Proposition \ref{generate3}.
We can assume $\ell>0$. Then we have a triangle
\[N'_1\to M'_0\xrightarrow{f'_0}N_0\to N'_1[1]\]
with $M'_0\in\MM$ and $N'_1\in\MM*\MM[1]*\cdots*\MM[\ell-1]$.
Since $N'_1\in\TT_{\MM}^{\le0}$, we have that $f'_0$ is a right $\MM$-approximation.
Thus we can write
$f'_0=(f_0\ 0):M_0=M'_0\oplus M''_0\to N_0$ with a minimal right $\MM$-approximation $f_0$.
Then we have a triangle
\[N_1\xrightarrow{g_1}M_0\xrightarrow{f_0}N_0\to N_1[1]\]
such that $g_1$ belongs to $J_{\TT}$ and $N_1$ is a summand of $N'_1$.
By Lemma \ref{direct summand}, we have $N_1\in\MM*\MM[1]*\cdots*\MM[\ell-1]$.
Repeating similar construction, we obtain the desired triangles.
\end{proof}

The following observation is often useful.

\begin{lemma}\label{no common summand}
Let $\MM\in\silt\TT$ and $N_0,N_0'\in\TT_{\MM}^{\le0}$.
For $N_0$, we take $\ell\ge0$ and triangles in Proposition \ref{resolution}.
Also for $N_0'$, we take triangles 
\[\xymatrix@R=0.2cm{
N_1'\ar[r]^{g_1'}&M_0'\ar[r]^{f_0'}&N_0'\ar[r]&N_1'[1],\\
&\cdots,\\
N_{\ell'}'\ar[r]^{g_\ell'}&M_{\ell'-1}'\ar[r]^{f_{\ell'-1}'}&N_{\ell'-1}'\ar[r]&N_{\ell'}'[1],\\
0\ar[r]^{g_{\ell'+1}'}&M_{\ell'}'\ar[r]^{f_{\ell'}'}&N_{\ell'}'\ar[r]&0,
}\]
satisfying the same properties.
If $\Hom_{\TT}(N_0,N_0'[\ell])=0$ holds, then we have $(\add M_\ell)\cap(\add M_0')=0$.
\end{lemma}

\begin{proof}
For $\ell=0$, the assertion follows from right minimality of $f_0'$. So we assume $\ell>0$.
We only have to show that any morphism $a:M_\ell\to M_0'$ belongs
to $J_{\TT}$. Applying $\Hom_{\TT}(-,N_0')$ to triangles in Proposition \ref{resolution}, we have
\[\Hom_{\TT}(N_{\ell-1}[-1],N_0')\simeq\Hom_{\TT}(N_{\ell-2}[-2],N_0')
\simeq\cdots\simeq\Hom_{\TT}(N_0[-\ell],N_0')=0.\]
Thus we have the following commutative diagram of triangles.
\[\xymatrix{
N_{\ell-1}[-1] \ar[r] \ar[d] & M_\ell \ar[r]^{g_{\ell}} \ar[d]^{a} & M_{\ell-1} \ar[r]^{f_{\ell-1}} \ar[d]^{b} & N_{\ell-1}\ar[d] \\
N_1' \ar[r]^{g_{1}'} & M_0' \ar[r]^{f_{0}'} & N_0'\ar[r]&N_1'[1].
}\]
Since $f_0'$ is a right $\MM$-approximation,
there exists $b':M_{\ell-1}\to M_0'$ such that $b=f_0'b'$.
Since $f_0'(a-b'g_\ell)=0$, there exists $a':M_\ell\to N_1'$
such that $a=b'g_\ell+g_1'a'$. Since both $g_\ell$ and $g_1'$ belong to $J_{\TT}$,
we have $a\in J_{\TT}$.
Thus the proof is completed.
\end{proof}

As an application, we have the following result.

\begin{theorem}\label{indecomposable case}
If $\TT$ has an indecomposable silting object $M$, 
then we have $\silt{\TT}=\{M[i]\ |\ i\in \Bbb{Z}\}$.  
\end{theorem}

\begin{proof}
Let $N$ be a basic silting object in $\TT$. 
Take the smallest integer $k\in\Z$ such that $\homo{\TT}{M}{N[k]}\not=0$.
Replacing $N$ by $N[k]$, we can assume $N\in \TT^{\leq 0}_{M}$
and $\homo{\TT}{M}{N}\not=0$. 
We have triangles in Proposition \ref{resolution}. 
Then we have $M_0\neq0$ since $\homo{\TT}{M}{N}\neq0$,
and moreover we can assume $M_{\ell}\neq0$.
%\[\xymatrix@!R=1pt{
%N_{1} \ar[r] & M_{0} \ar[r]^{f_{0}} & N \ar[r] & N_{1}[1] \\
%\cdots ,\\
%N_{\ell} \ar[r] & M_{\ell-1} \ar[r]^{f_{\ell-1}} & N_{\ell-1} \ar[r] & N_{\ell}[1] \\
%0 \ar[r] & M_{\ell} \ar[r]^{f_{\ell}} & N_{\ell} \ar[r] & 0
%}\]
%for some $\ell\geq 0$ such that $f_{i}$ is a minimal right $\add{M}$-approximation. 

Assume $\ell>0$. Then we have $\Hom_{\TT}(N,N[\ell])=0$ and so
$\add{M_{\ell}}\cap \add{M_{0}}=0$ by Lemma \ref{no common summand}. 
Since both $M_0$ and $M_{\ell}$ are non-zero objects in $\add M$
where $M$ is indecomposable, this is a contradiction.
Thus we have $\ell=0$ and $N\simeq M_0\in \add M$.
Since $N$ is basic, we have $N\simeq M$.
\end{proof}

%%%%%%%%%%%%%%%%%%%%%%%%%%%%%%%%%%%%%%%%%%%%%%%%%%%%%%%%%%%%%%%%%%%%%%%%%%%%%%%%%%%%%%%%%%%%%%%%%%%%%%%%%%%

Next we shall show the following description of Grothendieck groups of triangulated categories with silting subcategories.

\begin{theorem}\label{grothendieck}
Let $\TT$ be a Krull-Schmidt triangulated category with a silting subcategory $\MM$.
Then the Grothendieck group $K_0(\TT)$ of $\TT$ is a free abelian group with a basis $\ind\MM$.
\end{theorem}
 
For an object $X\in\TT$, we denote by $\delta(X)$ the number of non-isomorphic indecomposable summands of $X$.
As an immediate consequence of Theorem \ref{grothendieck}, we have the following result.

\begin{corollary}\label{thm:number of indecomposable direct summands}
For any silting objects $M,N\in \TT$, we have $\delta (M)=\delta (N)$.
\end{corollary}

Let us start with proving Theorem \ref{grothendieck}.

Since $\thick\MM=\TT$, we have that $\ind\MM$ generate $K_0(\TT)$.
It is enough to show that they are linearly independent.
To prove this, we shall define a map
\[\gamma:{\rm ob}\TT\to\Z^{\ind\MM}.\]
\begin{itemize}
\item The map $\gamma$ naturally identifies ${\rm ob}\MM$ with $\Z_{\ge0}^{\ind\MM}$.
%\[\gamma(M^{(1)}{}^{a_1}\coprod\cdots\coprod M^{(n)}{}^{a_n}):=(a_1,\cdots,a_n).\]
\item For any $N\in\TT_{\MM}^{\le0}$, we take triangles in Proposition \ref{resolution} and put
\[\gamma(N):=\sum_{i=0}^\ell(-1)^i\gamma(M_i).\]
\item For general $N\in\TT$, take a sufficiently large $k$ such that $N[k]\in\TT_{\MM}^{\le0}$ and put
\[\gamma(N):=(-1)^k\gamma(N[k]).\]
\end{itemize}
In other words, for any $\ell\ge0$ and $M_i\in\MM$, we put
\[\gamma(M_{-\ell}[-\ell]*M_{1-\ell}[-\ell]*\cdots*M_{\ell-1}[\ell-1]*M_\ell[\ell]):=\sum_{i=-\ell}^\ell(-1)^i\gamma(M_i).\]
The crucial step is to prove the following observation.

\begin{lemma}\label{gamma is on grothendieck}
Let $X\to Y\to Z\to X[1]$ be a triangle in $\TT$. Then we have $\gamma(X)-\gamma(Y)+\gamma(Z)=0$.
\end{lemma}

\begin{proof}
Without loss of generality, we can assume $X,Y,Z\in\TT_{\MM}^{\le0}$.
To use induction, we consider the following assertions for $\ell\ge0$.

(i)$_\ell$ The statement is true if $X,Y,Z\in \MM*\MM[1]*\cdots*\MM[\ell]$.

(ii)$_\ell$ The statement is true if $X,Y\in \MM*\MM[1]*\cdots*\MM[\ell]$ and $Z\in\MM*\MM[1]*\cdots*\MM[\ell+1]$.

(iii)$_\ell$ The statement is true if $X\in \MM*\MM[1]*\cdots*\MM[\ell]$ and $Y,Z\in\MM*\MM[1]*\cdots*\MM[\ell+1]$.

The statement (i)$_0$ is true since any triangle $X\to Y\to Z\to X[1]$ with $X,Y,Z\in\MM$ splits,
so we have $Y\simeq X\coprod Z$ and $\gamma(X)-\gamma(Y)+\gamma(Z)=0$.

We shall show (i)$_\ell\Rightarrow$(ii)$_\ell$ for any $\ell\ge0$.

Let $X\to Y\to Z\to X[1]$ be a triangle with $X,Y\in \MM*\MM[1]*\cdots*\MM[\ell]$ and $Z\in\MM*\MM[1]*\cdots*\MM[\ell+1]$.
Take a triangle $Z_1\to M_0\to Z\to Z_1[1]$ such that $M_0\in\MM$, $Z_1\in\MM*\MM[1]*\cdots\MM[\ell]$ and
\begin{equation}\label{g1}
\gamma(Z)=\gamma(M_0)-\gamma(Z_1).
\end{equation}
By octahedral axiom, we have the following commutative diagram of triangles:
\[\xymatrix@R=.4cm{
&Z_1[1]\ar@{=}[r]&Z_1[1]\\
X\ar[r]&Y\ar[r]\ar[u]&Z\ar[r]\ar[u]&X[1]\\
X\ar[r]\ar@{=}[u]&W\ar[r]\ar[u]&M_0\ar[r]\ar[u]&X[1]\ar@{=}[u]\\
&Z_1\ar@{=}[r]\ar[u]&Z_1\ar[u].
}\]
Since $\Hom_{\TT}(M_0,X[1])=0$, the lower horizontal triangle splits and we have $W\simeq X\oplus M_0$. In particular, we have
\begin{equation}\label{g2}
\gamma(W)=\gamma(X)+\gamma(M_0).
\end{equation}
Since the left vertical triangle $Z_1\to W\to Y\to Z_1[1]$
satisfies $Z_1,W,Y\in\MM*\MM[1]*\cdots*\MM[\ell]$,
our assumption (i)$_\ell$ implies
\begin{equation}\label{g3}
\gamma(Z_1)-\gamma(W)+\gamma(Y)=0.
\end{equation}
Using \eqref{g1}, \eqref{g2} and \eqref{g3}, we have
\[\gamma(X)-\gamma(Y)+\gamma(Z)=
(\gamma(Z)-\gamma(M_0)+\gamma(Z_1))+(\gamma(X)-\gamma(W)+\gamma(M_0))-(\gamma(Z_1)-\gamma(W)+\gamma(Y))=0.\]
Thus (ii)$_\ell$ holds.

By a quite similar argument, one can show (ii)$_\ell\Rightarrow$(iii)$_\ell\Rightarrow$(i)$_{\ell+1}$ for any $\ell\ge0$.
Thus the assertion follows inductively.
\end{proof}

Now we are ready to prove Theorem \ref{grothendieck}.

By Lemma \ref{gamma is on grothendieck}, we have a homomorphism $\gamma:K_0(\TT)\to\Z^{\ind\MM}$.
Since $\gamma(\ind\MM)$ is a basis of $\Z^n$,
the set $\ind\MM$ must be linearly independent in $K_0(\TT)$.
Thus it forms a basis of $K_0(\TT)$.
\qed


%%%%%%%%%%%%%%%%%%%%%%%%%%%%%%%%%%%%%%%%%%%%%%%%%%%%%%%%%%%%%%%%%%%%%%%%%%%%%%%%%%%%%%%%%%%%%%%%%%%%%%%%%%%
%%%%%%%%%%%%%%%%%%%%%%%%%%%%%%%%%%%%%%%%%%%%%%%%%%%%%%%%%%%%%%%%%%%%%%%%%%%%%%%%%%%%%%%%%%%%%%%%%%%%%%%%%%%

\subsection{Silting mutation}

The aim of this subsection is to introduce silting mutation and give its basic properties.
Let $\TT$ be a triangulated category. We do not assume anything else on $\TT$ unless otherwise stated.

\begin{definition}\label{mutation}
Let $\MM\in\silt\TT$. For a covariantly finite subcategory $\DD$ of $\MM$, we define a subcategory $\mu^+(\MM;\DD)$ of $\TT$ as follows:
For any $M\in\MM$ we take a left $\DD$-approximation $f:M\to D$ and a triangle
\begin{equation}\label{exchange}
M\xrightarrow{f}D\xrightarrow{g}N_M\to M[1].
\end{equation}
We put
\begin{eqnarray*}
\mu^+(\MM;\DD)&:=&\add(\DD\cup\{N_M\ |\ M\in\MM\}).
\end{eqnarray*}
It is easily checked that $\mu^+(\MM;\DD)$ does not depend on a choice of left approximation $f$.
We call $\mu^+(\MM;\DD)$ a \emph{left mutation} of $\MM$.
Dually, we define a \emph{right mutation} $\mu^-(\MM;\DD)$ for a contravariantly finite subcategory $\DD$ of $\MM$.
\emph{(Silting) mutation} is a left or right mutation.
\end{definition}

\begin{theorem}\label{thm:mutation is silting}
Any mutation of a silting subcategory is again a silting subcategory.
\end{theorem}

\begin{proof}
The triangle \eqref{exchange} shows $\thick\mu^+(\MM;\DD)\supset\thick\MM=\TT$.

We only have to show $\Hom_{\TT}(\mu^+(\MM;\DD),\mu^+(\MM;\DD)[>0])=0$.

Applying $\Hom_{\TT}(\DD,-[>0])$ to \eqref{exchange}, we have an exact sequence
\begin{eqnarray*}%\label{apply (M,-)}
0=\Hom_{\TT}(\DD,D[>0])\to\Hom_{\TT}(\DD,N_M[>0])\to\Hom_{\TT}(\DD,M[>1])=0
\end{eqnarray*}
and so $\Hom_{\TT}(\DD,N_M[>0])=0$.
Applying $\Hom_{\TT}(-,\DD[>0])$ to \eqref{exchange}, we have an exact sequence
\begin{eqnarray*}%\label{apply (-,M[i])}
\Hom_{\TT}(D[1],\DD[>0])\xrightarrow{\cdot f[1]}
\Hom_{\TT}(M[1],\DD[>0])\to\Hom_{\TT}(N_M,\DD[>0])\to\Hom_{\TT}(D,\DD[>0])=0.
\end{eqnarray*}
Since the above $(\cdot f[1])$ is surjective, we have $\Hom_{\TT}(N_M,\DD[>0])=0$. 
Applying $\Hom_{\TT}(-,M[>1])$ to \eqref{exchange}, we have an exact sequence
\begin{equation*}%\label{apply (-,X[>1])}
0=\Hom_{\TT}(M[1],M[>1])\to\Hom_{\TT}(N_M,M[>1])\to\Hom_{\TT}(D,M[>1])=0
\end{equation*}
and so $\Hom_{\TT}(N_M,M[>1])=0$.
Applying $\Hom_{\TT}(N_M,-[>0])$ to \eqref{exchange}, we have an exact sequence
\begin{equation*}%\label{apply (Y_X,-)}
0=\Hom_{\TT}(N_M,D[>0])\to\Hom_{\TT}(N_M,N_M[>0])\to\Hom_{\TT}(N_M,M[>1])=0
\end{equation*}
and so we have $\Hom_{\TT}(N_M,N_M[>0])=0$.

Consequently we have $\Hom_{\TT}(\mu^+(\MM;\DD),\mu^+(\MM;\DD)[>0])=0$.
\end{proof}

In general silting mutation of a tilting subcategory is not necessarily a tilting subcategory. 
For this, we have the following criterion. 

\begin{theorem}\label{when silting is tilting}
Let $\MM$ be a tilting subcategory of $\TT$.
\begin{itemize}
\item[(a)] For a covariantly finite subcategory $\DD$ of $\MM$, the following conditions are equivalent.
\begin{itemize}
\item[(i)] $\mu^+(\MM;\DD)$ is tilting.
\item[(ii)] Any $M\in\MM$ has a left $\DD$-approximation $f$ such that
$\homo{\TT}{\DD}{f}$ is injective.
\end{itemize}
\item[(b)] For a contravariantly finite subcategory $\DD$ of $\MM$, the following conditions are equivalent.
\begin{itemize}
\item[(i)] $\mu^-(\MM;\DD)$ is tilting.
\item[(ii)] Any $M\in\MM$ has a right $\DD$-approximation $g$ such that
$\homo{\TT}{g}{\DD}$ is injective.
%$\homo{\TT}{g}{\MM_{\XX}}$ is a monomorphism where $g$ is a minimal right $\MM_{\XX}$-approximation of $X$ for any $X\in \XX$. 
\end{itemize}
\end{itemize}
In these cases mutation is called \emph{tilting mutation}.
%We say that mutation is \emph{tilting mutation} if both $\MM$ and $\mu_{\XX}^{-}(\MM)$ ($\mu_{\XX}^{+}(\MM)$) are tilting.  
\end{theorem}

\begin{proof}
We only prove the statement (a). 
The proof is parallel to that of Theorem \ref{thm:mutation is silting}.

Applying $\Hom_{\TT}(\DD,-[\neq0])$ to \eqref{exchange}, we have an exact sequence
\begin{eqnarray*}%\label{apply (M,-)}
0=\Hom_{\TT}(\DD,D[\neq0])\to\Hom_{\TT}(\DD,N_M[\neq0])\to\Hom_{\TT}(\DD,M[\neq1])\xrightarrow{f\cdot}\Hom_{\TT}(\DD,D[\neq1]).
\end{eqnarray*}
Thus $\Hom_{\TT}(\DD,N_M[\neq0])=0$ holds if and only if (ii) holds.

Applying $\Hom_{\TT}(-,\DD[\neq0])$ to \eqref{exchange}, we have an exact sequence
\begin{eqnarray*}%\label{apply (-,M[i])}
\Hom_{\TT}(D[1],\DD[\neq0])\xrightarrow{\cdot f[1]}
\Hom_{\TT}(M[1],\DD[\neq0])\to\Hom_{\TT}(N_M,\DD[\neq0])\to\Hom_{\TT}(D,\DD[\neq0])=0.
\end{eqnarray*}
Since the above $(\cdot f[1])$ is surjective, we have $\Hom_{\TT}(N_M,\DD[\neq0])=0$.

%The implication $(i)\Rightarrow (ii)$ can be shown easily. 
%We shall show (i)$\Rightarrow$(ii). 
%Let $X\in \XX$ and put $Y=Y_{X}$ in the triangle \eqref{exchange}. 
%By the exact sequence \eqref{apply (-,M[i])}, we have $\homo{\TT}{Y}{\MM_{\XX}[i]}=0$ for any $i\not=0$. 
%Since the exact sequence \eqref{apply (M,-)} and $\homo{\TT}{\MM_{\XX}}{f}$ is a monomorphism, 
%we have $\homo{\TT}{\MM_{\XX}}{Y[i]}=0$ for any $i\not=0$. 

Applying $\Hom_{\TT}(M,-[\neq-1,0])$ to \eqref{exchange}, we have an exact sequence 
\[0=\homo{\TT}{M}{D[\neq-1,0]}\to \homo{\TT}{M}{N_M[\neq-1,0]}\to \homo{\TT}{M}{M[\neq0,1]}=0\]
and so $\homo{\TT}{M}{N_M[\neq-1,0]}=0$. 
Applying $\Hom_{\TT}(-,N_M[\neq0,1])$ to \eqref{exchange}, we have an exact sequence 
\[0=\homo{\TT}{M}{N_M[\neq-1,0]}\to \homo{\TT}{N_M}{N_M[\neq0,1]}\to \homo{\TT}{D}{N_M[\neq0,1]}=0\]
and so $\homo{\TT}{N_M}{N_M[\neq0,1]}=0$. 
Thus we have $\homo{\TT}{N_M}{N_M[\neq0]}=0$ by Theorem \ref{thm:mutation is silting}.

Consequently we have $\homo{\TT}{\mu^+(\MM;\DD)}{\mu^+(\MM;\DD)[\neq0]}=0$
if and only if (ii) holds. 
\end{proof}

The following properties can be checked easily.

\begin{proposition}\label{mutation basic}
Let $\MM\in\silt\TT$.
\begin{itemize}
\item[(a)] For a covariantly finite subcategory $\DD$ of $\MM$,
we have $\mu^-(\mu^+(\MM;\DD);\DD)=\MM$ and $\MM\ge\mu^+(\MM;\DD)$, where the equality holds if and only if $\DD=\MM$.
\item[(b)] For a contravariantly finite subcategory $\DD$ of $\MM$,
we have $\mu^+(\mu^-(\MM;\DD);\DD)=\MM$ and $\mu^-(\MM;\DD)\ge\MM$, where the equality holds if and only if $\DD=\MM$.
\end{itemize}
\end{proposition}

\begin{proof}
(a) Applying $\Hom_{\TT}(\DD,-)$ to \eqref{exchange}, we have an exact sequence
\[\Hom_{\TT}(\DD,D)\stackrel{g\cdot}{\to}\Hom_{\TT}(\DD,N_M)\to\Hom_{\TT}(\DD,M[1])=0.\]
Thus $g$ is a right $\DD$-approximation.
Thus we have $\mu^-(\mu^+(\MM;\DD);\DD)=\MM$.

Applying $\Hom_{\TT}(\MM,-[>0])$ to \eqref{exchange}, we have an exact sequence
\[0=\Hom_{\TT}(\MM,D[>0])\to\Hom_{\TT}(\MM,N_M[>0])\to\Hom_{\TT}(\MM,M[>1])=0\]
and so $\Hom_{\TT}(\MM,N_M[>0])=0$.
Thus $\Hom_{\TT}(\MM,\mu^+(\MM;\DD)[>0])=0$ holds and we have $\MM\ge\mu^+(\MM;\DD)$.

If $\DD\neq\MM$, then take $M\in\MM\backslash\DD$.
Since $f$ is not a split monomorphism, we have $\Hom_{\TT}(N_M,M[1])\neq0$.
Thus $N_M\notin\MM$ holds, so we have $\MM\neq\mu^+(\MM;\DD)$.
\end{proof}

In the rest of this subsection, we assume that $\TT$ is a Krull-Schmidt triangulated category.
The following notation will be often used.

\begin{definition}
Let $\MM\in\silt\TT$ and $\XX$ a subcategory of $\MM$.
Define a subcategory $\MM_{\XX}$ of $\MM$ by $\ind\MM_{\XX}=\ind\MM\backslash\ind\XX$.
%If $\MM_{\XX}$ is a covariantly (respectively, contravariantly) finite subcategory of $\MM$
We put
\[\mu^{\pm}_{\XX}(\MM):=\mu^{\pm}(\MM;\MM_{\XX})\]
%\ \ \ (\mbox{respectively, }\ \mu^-_{\XX}(\MM):=\nu^-(\MM;\MM_{\XX}))}\]
We say that mutation is \emph{irreducible} if $\#\ind\XX=1$.
If $\XX=\add X$, we denote $\MM_{\XX}$ and $\mu_{\XX}^\pm$ by $\MM_X$ and $\mu_X^\pm$ respectively.
\end{definition}

Now we assume the following condition:
\begin{itemize}
\item[(F)] $\TT$ is Krull-Schmidt, and for any $\MM\in\silt\TT$ and $X\in\ind\MM$, the subcategory $\MM_X$ is functorially finite in $\MM$.
\end{itemize}
This is satisfied if $\TT$ is $k$-linear $\Hom$-finite and has a silting object (Proposition \ref{existence of additive generator}).

Under this condition, we shall show the following result.

\begin{theorem}\label{two quivers}
Let $\TT$ be a triangulated category satisfying the condition (F).
For any $\MM,\NN\in\silt\TT$, the following conditions are equivalent.
\begin{itemize}
\item[(a)] $\NN$ is an irreducible left mutation of $\MM$.
\item[(b)] $\MM$ is an irreducible right mutation of $\NN$.
\item[(c)] $\MM>\NN$ and there is no $\LL\in\silt\TT$ satisfying $\MM>\LL>\NN$.
\end{itemize}
\end{theorem}

The following property is crucial in our consideration.

\begin{proposition}\label{make closer}
If $\MM,\NN\in\silt\TT$ satisfy $\MM>\NN$,
then there exists an irreducible left mutation $\LL$ of $\MM$ such that $\MM>\LL\ge\NN$.
\end{proposition}

\begin{proof}
We take $N_0\in\NN$ which does not belong to $\MM$,
then consider $\ell\ge0$ and triangles in Proposition \ref{resolution}.
Then we have $\ell>0$. Take an indecomposable summand $X$ of $M_\ell$,
and let $\LL:=\mu_X^+(\MM)$ be an irreducible left mutation.
Then we have $\LL=\add(\MM_X\cup\{Y\})$, where $Y$ is given by the triangle
\[X\xrightarrow{f}D\xrightarrow{g}Y\to X[1]\]
with a left $\MM_X$-approximation $f$ of $X$.

We only have to prove $\LL\ge\NN$. We need to show $\Hom_{\TT}(Y,\NN[>0])=0$.
Since we have an exact sequence
\[0=\Hom_{\TT}(X,\NN[>0])\to\Hom_{\TT}(Y,\NN[>1])\to\Hom_{\TT}(D,\NN[>1])=0,\]
we have $\Hom_{\TT}(Y,\NN[>1])=0$.
Thus it remains to show $\Hom_{\TT}(Y,\NN[1])=0$.
Since we have an exact sequence
\[\Hom_{\TT}(D,\NN)\xrightarrow{\cdot f}\Hom_{\TT}(X,\NN)\to\Hom_{\TT}(Y,\NN[1])\to
\Hom_{\TT}(D,\NN[1])=0,\]
we only have to show that $\Hom_{\TT}(D,\NN)\xrightarrow{\cdot f}\Hom_{\TT}(X,\NN)$ is surjective.

Fix any $N_0'\in\NN$ and $a:X\to N_0'$. 
We consider the diagram
\[\xymatrix@R=.4cm{
&Y[-1] \ar[r]& X \ar[r]^{f} \ar[d]^{a} & D\ar[r]^g&Y\\
N_1'\ar[r]^{g_0'}&M_0' \ar[r]^{f_0'} & N_0' \ar[r] &  N_1'[1].
}\]
where the lower triangle is given in Lemma \ref{no common summand}.
Since $f_0'$ is a right $\MM$-approximation, there exists $b:X\to M_0'$ such that $a=f_0'b$.
Since $(\add M_\ell)\cap(\add M_0')=0$ holds by Lemma \ref{no common summand},
we have $X\notin\add M_0'$ and so $M_0'\in\MM_X$.
Since $f$ is a left $\MM_X$-approximation, there exists $c:D\to M_0'$
such that $b=cf$. Thus we have $a=(f_0'c)f$, and we have the assertion.
\end{proof}

Now we are ready to prove Theorem \ref{two quivers}.

(i) We shall show that (a) and (b) are equivalent.

We only have to show that (a) implies (b).
Assume that $\NN$ is an irreducible left mutation $\mu_X^+(\MM)$ of $\MM$
with respect to $X\in\ind\MM$. Take a triangle
\[X\xrightarrow{f}D\xrightarrow{g}Y\to X[1]\]
with a minimal left $\MM_X$-approximation $f$. By Proposition \ref{mutation basic},
we have $\MM=\mu_Y^-(\NN)$. Since $g$ is a right $\MM_X$-approximation, it is easy to check that $Y$ is indecomposable.
Thus $\MM$ is an irreducible right mutation of $\NN$.

(ii) We shall show that (a) and (c) are equivalent.

Assume that (c) is satisfied.
Since $\MM>\NN$, there exists an irreducible left mutation $\LL$ of $\MM$ such that $\MM>\LL\ge\NN$
by Proposition \ref{make closer}. By the condition (c), we have $\LL=\NN$,
and we have the assertion.

Assume that (a) is satisfied. Thus $\NN=\mu_X^+(\MM)$ for indecomposable $X$.
Assume that there exists $\LL$ such that $\MM>\LL>\NN$.
Again by Proposition \ref{make closer}, there exists a left mutation
$\KK$ of $\MM$ such that $\MM>\KK\ge\LL$.
By Proposition \ref{common summand}, we have $\MM_X=\MM\cap\NN\subset\KK$.
Since both of $\NN$ and $\KK$ are left mutation of $\MM$ with respect to $X$,
we have $\NN=\KK$, a contradiction. Thus (c) holds.
\qed


%%%%%%%%%%%%%%%%%%%%%%%%%%%%%%%%%%%%%%%%%%%%%%%%%%%%%%%%%%%%%%%%%%%%%%%%%%%%%%%%%%%%%%%%%%%%%%%%%%%%%%%%%%%
\subsection{Silting reduction}
In this subsection, we give a reduction theorem of silting subcategories,
which is an analogue of 2-Calabi-Yau reduction in cluster tilting theory \cite[4.9]{IY}.
The following result gives a bijection between silting subcategories of $\TT$
containing a functorially finite thick subcategory $\SSS$ and silting subcategories of the quotient triangulated category $\TT/\SSS$.

\begin{theorem}\label{silting reduction}
Let $\TT$ be a Krull-Schmidt triangulated category, $\SSS$ a thick subcategory of $\TT$
and $\UU:=\TT/\SSS$. Let $F:\TT\to\UU$ be the canonical functor.
\begin{itemize}
\item[(a)] If $\SSS$ is a contravariantly finite subcategory of $\TT$, then for any
$\DD\in\silt\SSS$ we have an injective map
\[\{\MM\in\silt\TT\ |\ \DD\subset\MM\}\to\silt\UU,\ \ \ \MM\mapsto F\MM.\]
\item[(b)] If $\SSS$ is a functorially finite subcategory of $\TT$, then the map in (a) is bijective.
\end{itemize}
\end{theorem}

The first step of the proof is to consider the subcategory $\SSS^{\perp_{\TT}}$ of $\TT$.
Since $\SSS$ is contravariantly finite, we have a stable t-structure $(\SSS,\SSS^{\perp_{\TT}})$ of $\TT$
by Lemma \ref{direct summand}.
Moreover, we can naturally identify $\UU$ with the subcategory $\SSS^{\perp_{\TT}}$ of $\TT$ \cite{M}.

Next we shall show the following observation.

\begin{lemma}\label{S in S^le0}
Let $\MM$ be as in Theorem \ref{silting reduction}(a).
For any $M\in\MM$, take a triangle
\begin{equation}\label{split of M}
S\stackrel{a}{\to}M\stackrel{b}{\to}FM\to S[1]
\end{equation}
with $S\in\SSS$ and $FM\in\UU$. 
Then $S\in \SSS^{\leq 0}_{\DD}$ holds, where
\[({}^{\perp_{\SSS}}(\SSS^{\leq 0}_{\DD}),\SSS^{\leq0}_{\DD})
=(\bigcup_{\ell>0}\DD[-\ell]*\cdots*\DD[-1],\bigcup_{\ell\ge0}\DD*\cdots*\DD[\ell])\]
is the torsion pair in $\SSS=\thick\DD$ given in Proposition \ref{generate3}.
\end{lemma}

\begin{proof}
Take a triangle
\[S^{>0}\stackrel{c}{\to}S\stackrel{d}{\to}S^{\le0}\to S^{>0}[1]\]
with $S^{>0}\in {}^{\perp_{\SSS}}(\SSS^{\leq 0}_{\DD})$ and $S^{\le0}\in\SSS^{\leq 0}_{\DD}$.
Since $\Hom_{\TT}({}^{\perp_{\SSS}}(\SSS^{\leq 0}_{\DD}),\MM)=0$, we have $ac=0$.
Thus there exists a morphism $e:S^{\le0}\to M$ such that $a=ed$.
Since $be=0$ by $\Hom_{\DD}(\SSS,\UU)=0$, there exists $f:S^{\le0}\to S$ such that $e=af$.
\[\xymatrix@R=.4cm{
&S^{>0}\ar[d]^c\\
FM[-1]\ar[r]&S\ar[r]^a\ar@<.2em>[d]^d&M\ar[r]^b&FM\\
&S^{\le0}\ar@<.2em>[u]^f\ar[ru]_e
}\]
Since $a(1_S-fd)=0$, we have that $1_S-fd:S\to S$ factors through $FM[-1]$.
Since $\Hom_{\TT}(\SSS,\UU)=0$ again, we have $1_S=fd$.
Thus $d$ is a split monomorphism, and we have $S\in \SSS^{\leq 0}_{\DD}$.
\end{proof}

Now we are ready to prove Theorem \ref{silting reduction}.

(a)(i) We shall show that $F\MM$ is a silting subcategory of $\UU$.

We have $\thick F\MM\supset F\thick\MM=F\TT=\UU$.
Thus we only have to show $\Hom_{\TT}(F\MM,F\MM[>0])=0$.

Applying $\Hom_{\TT}(M,-[>0])$ to \eqref{split of M}, we have an exact sequence
\[0\stackrel{\MM\in\silt\TT}{=}\Hom_{\TT}(M,M[>0])\to\Hom_{\TT}(M,FM[>0])\to\Hom_{\TT}(M,S[>1])\stackrel{\Hom_{\TT}(\MM,\SSS^{<0}_{\DD})=0}{=}0\]
and so $\Hom_{\TT}(M,FM[>0])=0$.
Applying $\Hom_{\TT}(-,FM[>0])$ to \eqref{split of M}, we have an exact sequence
\[0\stackrel{\Hom_{\TT}(\SSS,\UU)=0}{=}\Hom_{\TT}(S,FM[>-1])\to\Hom_{\TT}(FM,FM[>0])\to\Hom_{\TT}(M,FM[>0])=0\]
and so $\Hom_{\TT}(FM,FM[>0])=0$.

(ii) We shall show that the correspondence $\MM\mapsto F\MM$ is injective. 
By Theorem \ref{partial order}, it is enough to prove $\MM\geq \NN$ if $F\MM\geq F\NN$. 
Let $M\in \MM$ and $N\in \NN$. 
For any $f\in \homo{\TT}{M}{N[>0]}$, we have a commutative diagram 
\[\xymatrix@R=0.4cm{
S_M \ar[r]\ar[d] & M \ar[r]\ar[d]^{f} & FM \ar[r]\ar[d] & S_M[1] \ar[d] \\
S_N[>0] \ar[r]_{a_N[>0]} & N[>0] \ar[r]_{b_N[>0]} & FN[>0] \ar[r] & S_N[>1].
}\]
Since $b_N[>0]\cdot f=0$ by our assumption, $f$ factors through $a_N[>0]: S_N[>0]\to N[>0]$. 
This implies $f=0$ since $\homo{\TT}{M}{S_N[>0]}\subset \homo{\TT}{\MM}{\SSS^{<0}_{\DD}}=0$ by Lemma \ref{S in S^le0}. 
Thus we have $\MM\geq \NN$. 

(b) We shall show that the correspondence $\MM\mapsto F\MM$ is surjective.
Fix $\NN\in\silt\UU$.

Since $\SSS$ is covariantly finite in $\TT$ by our assumption and
$\SSS^{<0}_{\DD}$ is covariantly finite in $\SSS$ by Proposition \ref{generate3}, 
we have that $\SSS^{<0}_{\DD}$ is covariantly finite in $\TT$.
Thus we have a torsion pair $({}^{\perp_{\TT}}(\SSS^{<0}_{\DD}), \SSS^{<0}_{\DD})$ in $\TT$ by Lemma \ref{direct summand}(c).
For any $N\in\NN$, take a triangle
\begin{equation}\label{split of N}
S_N\to M_N\to N\to S_N[1]
\end{equation}
with $M_N\in {}^{\perp_{\TT}}(\SSS^{<0}_{\DD})$ and $S_N\in \SSS^{\leq 0}_{\DD}$.
Notice that $M_N$ is unique up to a summand in $\DD$
since $\SSS^{\leq 0}_{\DD}\cap{}^{\perp_{\TT}}(\SSS^{<0}_{\DD})=\SSS^{\leq 0}_{\DD}\cap{}^{\perp_{\SSS}}(\SSS^{<0}_{\DD})=\DD$ by Proposition \ref{recover}.

We shall show that $\MM:=\add(\DD\cup\{M_N\ |\ N\in\NN\})$ is a silting subcategory of $\TT$.
Since $\DD[>0]\subset\SSS^{<0}_{\DD}$, we have $\Hom_{\TT}(M_N,\DD[>0])=0$.
Applying $\Hom_{\TT}(\DD,-[>0])$ to \eqref{split of N}, we have an exact sequence
\[0\stackrel{\Hom_{\TT}(\DD,\SSS^{<0}_{\DD})=0}{=}\Hom_{\TT}(\DD,S_N[>0])\to\Hom_{\TT}(\DD,M_N[>0])\to\Hom_{\TT}(\DD,N[>0])
\stackrel{\Hom_{\TT}(\SSS,\UU)=0}{=}0\]
and so $\Hom_{\TT}(\DD,M_N[>0])=0$.

Applying $\Hom_{\TT}(-,N[>0])$ to \eqref{split of N}, we have an exact sequence
\[0\stackrel{\NN\in\silt\UU}{=}\Hom_{\TT}(N,N[>0])\to\Hom_{\TT}(M_N,N[>0])\to\Hom_{\TT}(S_N,N[>0])\stackrel{\Hom_{\TT}(\SSS,\UU)=0}{=}0\]
and so $\Hom_{\TT}(M_N,N[>0])=0$.
Applying $\Hom_{\TT}(M_N,-[>0])$ to \eqref{split of N}, we have an exact sequence
\[0\stackrel{\Hom_{\TT}({}^{\perp_{\TT}}(\SSS^{<0}_{\DD}),\SSS^{<0}_{\DD})=0}{=}
\Hom_{\TT}(M_N,S_N[>0])\to\Hom_{\TT}(M_N,M_N[>0])\to\Hom_{\TT}(M_N,N[>0])=0\]
and so $\Hom_{\TT}(M_N,M_N[>0])=0$.

Consequently we have $\Hom_{\TT}(\MM,\MM[>0])=0$.
By \eqref{split of N}, we have $\thick\MM\supset\NN$, so we have $\thick\MM\supset\thick(\DD\cup\NN)\supset\SSS*\UU=\TT$.
Thus $\MM$ is a silting subcategory of $\TT$.
%
%(ii) It remains to show that the correspondence $\MM\mapsto F\MM$ is injective.
%
%We only have to show that, for any $M\in\MM$, the object $M_{FM}$ given
%by the triangle \eqref{split of N} for $N:=FM$ is isomorphic to $M$ up to a summand in $\DD$.
%This is clear since the triangle \eqref{split of M} gives the triangle
%\eqref{split of N} since $M\in {}^{\perp_{\TT}}(\SSS^{<0}_{\DD})$ and
%$S\in\SSS^{\leq 0}_{\DD}$ hold by Lemma \ref{S in S^le0}.
\qed

\medskip
The functor $F:\MM\to F\MM$ has the following properties.

\begin{proposition}\label{M and FM}
In Theorem \ref{silting reduction}(a), the functor $F:\MM\to F\MM$ induces an equivalence
\[\MM/[\DD]\to F\MM,\]
where $[\DD]$ is the ideal of $\MM$ consisting of morphisms which factor through objects in $\DD$.
In particular, $F$ induces a bijection between $\ind\MM\backslash\ind\DD$ and $\ind F\MM$.
\end{proposition}

\begin{proof}
Since $F\DD=0$, we have the induced functor $F:\MM/[\DD]\to F\MM$.
We only have to show that this is fully faithful.
Let $M,M'\in\MM$. For each $f:M\to M'$, the morphism $Ff:FM\to FM'$ is given by the following commutative diagram.
\[\xymatrix@R=.4cm{
S\ar[r]\ar[d]&M\ar[r]\ar[d]^f&FM\ar[r]\ar[d]^{Ff}&S[1]\ar[d]\\
S'\ar[r]&M'\ar[r]&FM'\ar[r]&S'[1]
}\]

(i) We shall show that $F:\MM/[\DD]\to F\MM$ is faithful.

If $Ff=0$, then $f$ factors through $S'$. By Lemma \ref{S in S^le0}, we know
$S'\in\SSS^{\leq0}_{\DD}=\bigcup_{\ell\ge0}\DD*\cdots*\DD[\ell]$.
Since $\Hom_{\TT}(M,\DD[1]*\cdots*\DD[\ell])=0$, we have that $f$ factors through $\DD$.
Thus $f$ is zero in $\MM/[\DD]$.

(ii) We shall show that $F:\MM/[\DD]\to F\MM$ is full.

Fix $g\in\Hom_{\TT}(FM,FM')$. Since $\Hom_{\TT}(M,S'[1])=0$ again by Lemma \ref{S in S^le0}, there exists $f:M\to M'$ which makes the following diagram commutative.
\[\xymatrix@R=.4cm{
S\ar[r]\ar[d]&M\ar[r]\ar[d]^f&FM\ar[r]\ar[d]^{g}&S[1]\ar[d]\\
S'\ar[r]&M'\ar[r]&FM'\ar[r]&S'[1]
}\]
Then $f=Fg$ holds.
\end{proof}

Later we use the following compatibility of silting mutation and silting reduction.

\begin{lemma}\label{compatibility}
In Theorem \ref{silting reduction}(a), we have $F\mu^+_{\XX}(\MM)=\mu^+_{F\XX}(F\MM)$ for any covariantly finite subcategory $\XX$ of $\MM$ satisfying $\DD\cap\XX=0$.
\end{lemma}

\begin{proof}
By Proposition \ref{M and FM}, we have $(F\MM)_{F\XX}=F(\MM_{\XX})$.
For any $X\in\XX$, take a triangle
\[X\stackrel{f}{\to} D\to Y\to X[1]\]
where $f$ is a left $\MM_{\XX}$-approximation. 
We have to prove that $Ff$ is a left $F(\MM_{\XX})$-approximation.
It is enough to show $\homo{\TT}{FY}{F(\MM_{\XX})[1]}=0$.  

For any $M\in \MM_{\XX}$, there exists a triangle 
\begin{equation}%\label{prop:compat 1}
S \to M \to FM \to S[1]
\end{equation}
with $S\in \SSS^{\leq 0}_{\DD}$ by Lemma \ref{S in S^le0}.
Applying $\homo{\TT}{Y}{-}$, we have an exact sequence 
\[
0 \stackrel{{\rm Thm.}\ \ref{thm:mutation is silting}}{=}  \homo{\TT}{Y}{M[1]}\to \homo{\TT}{Y}{FM[1]} \to 
\homo{\TT}{Y}{S[2]}\stackrel{S\in \SSS^{\leq 0}_{\DD}}{=}0,
\]
so we have $\homo{\TT}{Y}{FM[1]}=0$.
Thus we have $\homo{\TT}{FY}{FM[1]}\simeq \homo{\TT}{Y}{FM[1]}=0$. 
\end{proof}

%%%%%%%%%%%%%%%%%%%%%%%%%%%%%%%%%%%%%%%%%%%%%%%%%%%%%%%%%%%%%%%%%%%%%%%%%%%%%%%%%%%%%%%%%%%%%%%%%%%%%%%%%%
%%%%%%%%%%%%%%%%%%%%%%%%%%%%%%%%%%%%%%%%%%%%%%%%%%%%%%%%%%%%%%%%%%%%%%%%%%%%%%%%%%%%%%%%%%%%%%%%%%%%%%%%%%%

\subsection{Silting quivers and examples}

Let $\TT$ be a Krull-Schmidt triangulated category.
The aim of this section is to introduce the silting quiver of $\TT$.

\begin{definition}
The \emph{silting quiver} of $\TT$ is defined as follows:
\begin{itemize}
\item The set of vertices is $\silt\TT$.
\item We draw an arrow $\MM\to\NN$ if $\NN$ is an irreducible left mutation of $\MM$.
\end{itemize}
If $\TT$ satisfies the condition (F), then the silting quiver is nothing but the Hasse quiver of the partially ordered set $\silt\TT$ by Theorem \ref{two quivers}.
\end{definition}

We pose the following question.

\begin{question}\label{transitive question2}
Let $A$ be a finite dimensional algebra $A$ over a field and $\TT:=\K^{\rm b}(\proj A)$.
When is the silting quiver of $\TT$ connected?
Equivalently, when is the action of iterated irreducible mutation on $\silt\TT$ transitive?
\end{question}

If $A$ is local, then we have a positive answer by the following observation.

\begin{corollary}\label{corollary indecomposable case}
If $\TT$ has an indecomposable silting object, 
then the action of iterated irreducible mutation on $\silt\TT$ is transitive.
\end{corollary}

\begin{proof}
This is immediate from Theorem \ref{indecomposable case}
since we have $\mu_{M[i]}^{\pm}(M[i])=M[i\pm1]$.
\end{proof}

Next we shall generalize Corollary \ref{corollary indecomposable case} by
introducing the following notion which generalizes `almost complete partial tilting modules'.

We call a subcategory $\DD$ of $\TT$ \emph{almost complete silting} if there exists a
silting subcategory $\MM$ of $\TT$ such that $\MM\supset\DD$ and $\#(\ind\MM\backslash\ind\DD)=1$.

Then we have the following analogue of \cite[2.1]{HU} for `tilting mutation'
and \cite[5.3]{IY} for `cluster tilting mutation'.

\begin{theorem}\label{almost transitivity}
Let $\TT$ be a Krull-Schmidt triangulated category and $\DD$ an almost complete silting subcategory.
If $\thick\DD$ is a functorially finite subcategory of $\TT$, then the set $\{\MM\in\silt\TT\ |\ \DD\subset\MM\}$ is transitive under iterated irreducible mutation.
\end{theorem}

\begin{proof}
Let $\UU:=\TT/\thick\DD$.
By Theorem \ref{silting reduction} we have a bijection
$\{\MM\in\silt\TT\ |\ \DD\subset\MM\}\to\silt\UU$.
This commutes with mutation by Lemma \ref{compatibility} and its dual.
By Corollary \ref{corollary indecomposable case}, we know that $\silt\UU$ is transitive under iterated irreducible mutation.
Thus the assertion follows.
\end{proof}

We give examples of silting quivers.

\begin{example}
Let $A$ be a path algebra of the quiver $\xymatrix{1\ar[r]&2}$. 
We have the AR-quiver of $\K^{\rm b}(\proj{A})$ as follows:
\[\xymatrix@R.3cm@C.3cm{
& X_{-2} \ar[dr] &                & X_{0} \ar[dr] &               & X_{2} \ar[dr] & &\cdots\\
\cdots\ar[ur]&                & X_{-1} \ar[ur] &               & X_{1} \ar[ur] &               & X_3\ar[ur]
}\]

Then the silting quiver of $A$ is the following (cf. Theorem \ref{hereditary case}):
\[
\xymatrix@R=0.5cm{
\cdots\ar[dd]\ar[rddd]&&\cdots\ar[lddd]\ar[rddd]&&\cdots\ar[lddd]\\
&X_{-3}\oplus X_1\ar[lddd]\ar[rddd]&&X_{-6}\oplus X_4\ar[lddd]\ar[rddd]&\\
X_{-1}\oplus X_0\ar[dd]\ar[rddd]&&X_{-4}\oplus X_3\ar[lddd]\ar[rddd]&&\cdots\ar[lddd]\\
&X_{-2}\oplus X_2\ar[lddd]\ar[rddd]&&X_{-5}\oplus X_5\ar[lddd]\ar[rddd]&\\
X_0\oplus X_1\ar[dd]\ar[rddd]&&X_{-3}\oplus X_4\ar[lddd]\ar[rddd]&&\cdots\ar[lddd]\\
&X_{-1}\oplus X_3\ar[lddd]\ar[rddd]&&X_{-4}\oplus X_6\ar[lddd]\ar[rddd]&\\
X_1\oplus X_2\ar[dd]\ar[rddd]&&X_{-2}\oplus X_5\ar[lddd]\ar[rddd]&&\cdots\ar[lddd]\\
&X_0\oplus X_4\ar[lddd]\ar[rddd]&&X_{-1}\oplus X_7\ar[lddd]\ar[rddd]&\\
X_2\oplus X_3\ar[dd]&&X_{-1}\oplus X_6&&\\
&X_1\oplus X_5&&X_{-2}\oplus X_8&\\
\cdots &&\cdots&&\cdots
}
\]
Identifying each silting object $T$ with $T[i]$ for any $i\in\Z$, we can simplify the quiver as follows,
where $\xymatrix{M\ar[r]|{[1]}&N}$ means $\xymatrix{M\ar[r]&N[1]}$.
\[
\xymatrix{
X_0\oplus X_1\ar[dd]\ar@<.3em>[r]&X_0\oplus X_4\ar@<.3em>[r]\ar@<.3em>[l]|{[1]}&X_0\oplus X_7\ar@<.3em>[r]\ar@<.3em>[l]|{[1]}
&X_0\oplus X_{10}\ar@<.3em>[r]\ar@<.3em>[l]|{[1]}&\cdots\ar@<.3em>[l]|(0.4){[1]}\\
&X_2\oplus X_3\ar[lu]|{[1]}\ar@<.3em>[r]&X_2\oplus X_6\ar@<.3em>[r]\ar@<.3em>[l]|{[1]}
&X_2\oplus X_9\ar@<.3em>[r]\ar@<.3em>[l]|{[1]}&X_2\oplus X_{12}\ar@<.3em>[r]\ar@<.3em>[l]|{[1]}&\cdots\ar@<.3em>[l]|(0.35){[1]}\\
X_1\oplus X_2\ar[ru]\ar@<.3em>[r]&X_1\oplus X_5\ar@<.3em>[r]\ar@<.3em>[l]|{[1]}&X_1\oplus X_8\ar@<.3em>[r]\ar@<.3em>[l]|{[1]}
&X_1\oplus X_{11}\ar@<.3em>[r]\ar@<.3em>[l]|{[1]}&\cdots\ar@<.3em>[l]|(0.4){[1]}
}
\]
\end{example}

\begin{example}
Let $A$ be a path algebra of the quiver 
$\xymatrix{
1\ar@<.5em>[r]^{}="a" & 2 \ar@{<-}@<.5em>[l]^{}="b"  
\ar@{.}"a";"b"
}$ 
with $\ell\ge2$ arrows. 
The AR-quiver of $\K^{\rm b}(\proj{A})$ contains the following connected component:
\[\xymatrix@R.5cm@C.5cm{
& X_{-2}\ar@<-0.7em>[dr]^{}="a1"\ar@<0.5em>[dr]_{}="a2" & & X_{0} \ar@<-0.7em>[dr]^{}="b1"\ar@<0.5em>[dr]_{}="b2" 
& & X_{2} \ar@<-0.7em>[dr]^{}="c1"\ar@<0.5em>[dr]_{}="c2" & &\cdots\\
\cdots\ar@<0.5em>[ur]_{}="g1"\ar@<-0.7em>[ur]^{}="g2" & & X_{-1} \ar@<0.5em>[ur]_{}="d1"\ar@<-0.7em>[ur]^{}="d2" & & X_{1} \ar@<0.5em>[ur]_{}="e1"\ar@<-0.7em>[ur]^{}="e2" & & 
X_3\ar@<0.5em>[ur]_{}="f1"\ar@<-0.7em>[ur]^{}="f2"
\ar@{.}"a1";"a2" \ar@{.}"b1";"b2" \ar@{.}"c1";"c2" \ar@{.}"d1";"d2" \ar@{.}"e1";"e2" \ar@{.}"f1";"f2" \ar@{.}"g1";"g2"
}\]
Then the silting quiver of $A$ is the following (cf. Theorem \ref{hereditary case}):
\[
\xymatrix{
\cdots\ar[d]\\
X_{-1}\oplus X_0\ar[d]\ar@<.3em>[r]&X_{-1}\oplus X_0[1]\ar@<.3em>[r]\ar@<.3em>[l]|{[1]}
&X_{-1}\oplus X_0[2]\ar@<.3em>[r]\ar@<.3em>[l]|{[1]}&X_{-1}\oplus X_0[3]\ar@<.3em>[r]\ar@<.3em>[l]|{[1]}
&\cdots\ar@<.3em>[l]|(0.3){[1]}\\
X_0\oplus X_1\ar[d]\ar@<.3em>[r]&X_0\oplus X_1[1]\ar@<.3em>[r]\ar@<.3em>[l]|{[1]}&X_0\oplus X_1[2]\ar@<.3em>[r]\ar@<.3em>[l]|{[1]}
&X_0\oplus X_1[3]\ar@<.3em>[r]\ar@<.3em>[l]|{[1]}&\cdots\ar@<.3em>[l]|(0.35){[1]}\\
X_1\oplus X_2\ar[d]\ar@<.3em>[r]&X_1\oplus X_2[1]\ar@<.3em>[r]\ar@<.3em>[l]|{[1]}&X_1\oplus X_2[2]\ar@<.3em>[r]\ar@<.3em>[l]|{[1]}
&X_1\oplus X_2[3]\ar@<.3em>[r]\ar@<.3em>[l]|{[1]}&\cdots\ar@<.3em>[l]|(0.35){[1]}\\
X_2\oplus X_3\ar[d]\ar@<.3em>[r]&X_2\oplus X_3[1]\ar@<.3em>[r]\ar@<.3em>[l]|{[1]}&X_2\oplus X_3[2]\ar@<.3em>[r]\ar@<.3em>[l]|{[1]}
&X_2\oplus X_3[3]\ar@<.3em>[r]\ar@<.3em>[l]|{[1]}&\cdots\ar@<.3em>[l]|(0.35){[1]}\\
\cdots\\
}\]
\end{example}

\begin{example}
Let $A$ be an algebra presented by a quiver $\xymatrix{1\ar^a@<.3em>[r]&\ar^b@<.3em>[l]2}$
with relations $ab=0=ba$.
Then the silting quiver of $A$ is the following, where $X:={\rm cone}(P_1\to P_2)$ and $Y:={\rm cone}(P_2\to P_1)$:
\[
\xymatrix{
&X\oplus P_2\ar@<.3em>[r]\ar[d]&X[1]\oplus P_2\ar@<.3em>[r]\ar@<.3em>[l]|{[1]}
&X[2]\oplus P_2\ar@<.3em>[r]\ar@<.3em>[l]|{[1]}&X[3]\oplus P_2\ar@<.3em>[r]\ar@<.3em>[l]|{[1]}&\cdots\ar@<.3em>[l]|(0.35){[1]}\\
&X\oplus P_1[1]\ar@<.3em>[r]\ar[ld]|{[1]}&X\oplus P_1[2]\ar@<.3em>[r]\ar@<.3em>[l]|{[1]}
&X\oplus P_1[3]\ar@<.3em>[r]\ar@<.3em>[l]|{[1]}&X\oplus P_1[4]\ar@<.3em>[r]\ar@<.3em>[l]|{[1]}&\cdots\ar@<.3em>[l]|(0.35){[1]}\\
P_1\oplus P_2\ar[ruu]\ar[rd]\\
&Y\oplus P_1\ar@<.3em>[r]\ar[d]&Y[1]\oplus P_1\ar@<.3em>[r]\ar@<.3em>[l]|{[1]}
&Y[2]\oplus P_1\ar@<.3em>[r]\ar@<.3em>[l]|{[1]}&Y[3]\oplus P_1\ar@<.3em>[r]\ar@<.3em>[l]|{[1]}&\cdots\ar@<.3em>[l]|(0.35){[1]}\\
&Y\oplus P_2[1]\ar@<.3em>[r]\ar[luu]|{[1]}&Y\oplus P_2[2]\ar@<.3em>[r]\ar@<.3em>[l]|{[1]}
&Y\oplus P_2[3]\ar@<.3em>[r]\ar@<.3em>[l]|{[1]}&Y\oplus P_2[4]\ar@<.3em>[r]\ar@<.3em>[l]|{[1]}&\cdots\ar@<.3em>[l]|(0.35){[1]}
}\]
\end{example}

\begin{example}
Let $A$ be an algebra presented by a quiver $\xymatrix{1\ar^a@<.3em>[r]&\ar^b@<.3em>[l]2}$
with relations $(ab)^{\ell}a=0=(ba)^{\ell}b$ ($\ell\ge1$).
The silting quiver of $A$ is connected by \cite{A2}, and it is the following:
\[
{\small{
\xymatrix@R=0.1pt @C=0.3cm{
&& & && && & && & & && & & & &&\\
&& & && && & && & & && P_1 \ar[r] & P_2 \ar[r] & P_2 \ar[r] & P_2 &&\\
&& & && && & && & & && 0 \ar[r] & 0 \ar[r] & 0 \ar[r] & P_2 &&\\
&& & && && & && P_1 \ar[r] & P_2 \ar[r] & P_2 && & & & &&\\
&& & && && & && 0 \ar[r] & 0 \ar[r] & P_2 && & & & &&\\   
&& & && && & && & & && P_1 \ar[r] & P_2 \ar[r] & P_2 \ar[r] & P_2 &&\\
&& & && && & && & & && P_1 \ar[r] & P_2 \ar[r] & P_2 \ar[r] & 0 &&\\
&& 0 \ar[r] & P_1 && && P_1 \ar[r] & P_2 && & & && & & & &&\\
&& P_2 \ar[r] & P_1 && && 0 \ar[r] & P_2 && & & && & & & &&\\
&& & && && & && & & && P_1 \ar[r] & P_1\oplus P_2 \ar[r] & P_2 \ar[r] & P_2 &&\\
&& & && && & && & & && 0 \ar[r] & P_1 \ar[r] & P_2 \ar[r] & P_2 &&\\
&& & && && & && P_1 \ar[r] & P_2 \ar[r] & P_2 && & & & &&\\
&& & && && & && P_1 \ar[r] & P_2 \ar[r] & 0 && & & & &&\\
&& & && && & && & & && P_1 \ar[r] & P_1\oplus P_2 \ar[r] & P_2 \ar[r] & P_2 &&\\
&& & && && & && & & && P_1 \ar[r] & P_2 \ar[r] & 0 \ar[r] & 0 &&\\
&& & && P_1\oplus P_2 && & && & & && & & & &&\\
&& & && && & && & & && P_1 \ar[r] & P_1 \ar[r] & P_1\oplus P_2 \ar[r] & P_2 &&\\
&& & && && & && & & && 0 \ar[r] & 0 \ar[r] & P_1 \ar[r] & P_2 &&\\
&& & && && & && P_1 \ar[r] & P_1 \ar[r] & P_2 && & & & &&\\
&& & && && & && 0 \ar[r] & P_1 \ar[r] & P_2 && & & & &&\\
&& & && && & && & & && P_1 \ar[r] & P_1 \ar[r] & P_1\oplus P_2 \ar[r] & P_2 &&\\
&& & && && & && & & && P_1 \ar[r] & P_1 \ar[r] & P_2 \ar[r] & 0 &&\\
&& P_2 \ar[r] & 0   && && P_1 \ar[r] & P_2 && & & && & & & &&\\
&& P_2 \ar[r] & P_1 && && P_1 \ar[r] & 0 && & & && & & & &&\\
&& & && && & && & & && P_1 \ar[r] & P_1 \ar[r] & P_1 \ar[r] & P_2 &&\\
&& & && && & && & & && 0 \ar[r] & P_1 \ar[r] & P_1 \ar[r] & P_2 &&\\
&& & && && & && P_1 \ar[r] & P_1 \ar[r] & P_2 && & & & &&\\
&& & && && & && P_1 \ar[r] & 0 \ar[r] & 0 && & & & &&\\
&& & && && & && & & && P_1 \ar[r] & P_1 \ar[r] & P_1 \ar[r] & P_2 &&\\
&& & && && & && & & && P_1 \ar[r] & 0 \ar[r] & 0 \ar[r] & 0 &&\\
&& & && && & && & & && & & & &&
\save "2,15"."3,18"*[F]\frm{}="a" \restore
\save "6,15"."7,18"*[F]\frm{}="b" \restore
\save "10,15"."11,18"*[F]\frm{}="c" \restore
\save "14,15"."15,18"*[F]\frm{}="d" \restore
\save "17,15"."18,18"*[F]\frm{}="e" \restore
\save "21,15"."22,18"*[F]\frm{}="f" \restore
\save "25,15"."26,18"*[F]\frm{}="g" \restore
\save "29,15"."30,18"*[F]\frm{}="h" \restore
\save "4,11"."5,13"*[F]\frm{}="i" \restore
\save "12,11"."13,13"*[F]\frm{}="j" \restore
\save "19,11"."20,13"*[F]\frm{}="k" \restore
\save "27,11"."28,13"*[F]\frm{}="l" \restore
\save "8,8"."9,9"*[F]\frm{}="m" \restore
\save "23,8"."24,9"*[F]\frm{}="n" \restore
\save "16,6"*[F]\frm{}="o" \restore
\save "8,3"."9,4"*[F]\frm{}="p" \restore
\save "23,3"."24,4"*[F]\frm{}="q" \restore
\ar@<6em> "a";"b"
\ar@<6em> "c";"d"
\ar@<6em> "e";"f"
\ar@<6em> "g";"h"
\ar@<2.5em> "i";"j"
\ar@<2.5em> "k";"l"
\ar@<1.3em> "m";"n"
\ar@<1.3em> "p";"q"
\ar "2,18";"1,20" 
\ar "6,18";"5,20" 
\ar "10,18";"9,20" 
\ar "14,18";"13,20" 
\ar "17,18";"17,20" 
\ar "21,18";"20,20" 
\ar "25,18";"24,20" 
\ar "29,18";"28,20" 
\ar "4,20";"3,18" |(0.4){[1]} 
\ar "8,20";"7,18" |(0.4){[1]} 
\ar "12,20";"11,18" |(0.4){[1]} 
\ar "15,20";"15,18" |(0.4){[1]} 
\ar "19,20";"18,18" |(0.4){[1]} 
\ar "23,20";"22,18" |(0.4){[1]} 
\ar "27,20";"26,18" |(0.4){[1]} 
\ar "31,20";"30,18" |(0.4){[1]} 
\ar "4,13";{"3,15"+<-0.9em, -0.1em>}
\ar "12,13";{"11,15"+<-0.9em, -0.1em>}
\ar "19,13";{"18,15"+<-0.9em, -0.1em>}
\ar "27,13";{"26,15"+<-0.9em, -0.1em>}
\ar "b";"i" |(0.2){[1]}
\ar "d";"j" |(0.2){[1]}
\ar "f";"k" |(0.2){[1]}
\ar "h";"l" |(0.2){[1]}
\ar "8,9";{"5,11"+<-0.9em, -0.4em>}
\ar "23,9";{"20,11"+<-0.9em, -0.4em>}
\ar "j";"m" |(0.35){[1]}
\ar "l";"n" |(0.35){[1]}
\ar "o";"m"
\ar "n";"o" |{[1]}
\ar "o";"p"
\ar {"q"+<1.7em,0.9em>};"o" |{[1]}
\ar "8,3";"6,1"
\ar "23,3";"21,1"
\ar "11,1";"9,3" |(0.4){[1]}
\ar "26,1";"24,3" |(0.4){[1]}
}}}
\]
\end{example}


%%%%%%%%%%%%%%%%%%%%%%%%%%%%%%%%%%%%%%%%%%%%%%%%%%%%%%%%%%%%%%%%%%%%%%%%%%%%%%%%%%%%%%%%%%%%%%%%%%%%%%%%%%%
%%%%%%%%%%%%%%%%%%%%%%%%%%%%%%%%%%%%%%%%%%%%%%%%%%%%%%%%%%%%%%%%%%%%%%%%%%%%%%%%%%%%%%%%%%%%%%%%%%%%%%%%%%%

\subsection{Okuyama-Rickard complexes and APR and BB tilting modules} \label{Examples of silting subcategories and objects}

Our silting mutation has two important origins in representation theory: 
\begin{itemize}
\item Okuyama-Rickard complexes and Okuyama's method in modular representation theory;
\item APR tilting modules and BB tilting modules;
%\item Riedtmann-Schofield and Happel-Unger theory on tilting mutation.
\end{itemize}
In this section we will explain the relationship between silting mutation and these notions.

Throughout this section, let $A$ be a finite dimensional algebra over a field $k$. 
We may assume that $A$ is basic and indecomposable as an $A$-$A$-bimodule.
For an $A$-module $X$, we denote by $P(X)$ a projective cover of $X$.
We denote by $\nu =D\homo{A}{-}{A}:\modu{A}\to\modu{A}$ the Nakayama functor.
We denote by $\tau$ and $\tau^{-1}$ the Auslander-Reiten translations \cite{ARS,ASS}. 

%\subsection{Okuyama-Rickard complexes and Okuyama's method}

\begin{definition}
For idempotent $e\in A$,
%Set $U$ as the $(-1)$-shift of a minimal projective presentation of $eA/eA(1-e)A$. 
the \emph{Okuyama-Rickard complex} with respect to $e$ is defined by
\begin{equation}\label{OR complex}
T:=\left\{\begin{array}{ccc}
\stackrel{0}{P(eA(1-e)A)}&\stackrel{p_e}{\longrightarrow}&\stackrel{1}{eA}\\
\oplus&&\\
(1-e)A&&
\end{array}\right.
\end{equation}
where $p_e$ gives a projective cover of the submodule $eA(1-e)A$ of $eA$.
%Now we define $T=U\oplus (1-e)A$ and call it the \emph{Okuyama-Rickard complex} with respect to $e$. 
\end{definition}

In \cite{O}, Okuyama constructed Okuyama-Rickard complexes and proved that it is tilting if $A$ is symmetric. 
The method of his construction is often called \emph{Okuyama's method}. 

Let us give basic properties of Okuyama-Rickard complexes in our context of silting mutation.

\begin{theorem}\label{prop:1_right_mu}
Let $e\in A$ be an idempotent and $T$ the Okuyama-Rickard complex with respect to $e$. 
\begin{itemize}
\item[(a)] $T$ is isomorphic to right mutation $\mu^{-}_{eA}(A)$ of $A$ with respect to $eA$.
\item[(b)] $T$ is a silting object in $\K^{\rm{b}}(\proj{A})$.
\item[(c)] The following conditions are equivalent.
\begin{itemize}
\item[(i)] $T$ is a tilting object in $\K^{\rm{b}}(\proj{A})$.
\item[(ii)] $\homo{A}{eA/eA(1-e)A}{(1-e)A}=0$.
\end{itemize}
\item[(d)] If $A$ is a self-injective algebra and $eA\simeq\nu (eA)$ holds, then $T$ is a tilting object in $\K^{\rm{b}}(\proj{A})$.
\end{itemize}
\end{theorem}

\begin{proof}
(a) Since $eA(1-e)A$ is minimal amongst submodules $X$ of $eA$ such that any composition factor of $eA/X$ belongs to $\add({\rm top}\ eA)$,
we have that $p_e$ in \eqref{OR complex} is a minimal right $\add{(1-e)A}$-approximation of $eA$. 
Thus the assertion follows.

(b)(c) These are immediate consequences of Theorems \ref{thm:mutation is silting} and \ref{when silting is tilting}.

(d) This is immediate from (c) since $\add({\rm top}\ eA)\cap\add({\rm soc}(1-e)A)=\add({\rm top}\ eA)\cap\add({\rm top} (1-e)A)=0$ holds by our assumption
and any composition factor of $eA/eA(1-e)A$ is isomorphic to ${\rm top}\ eA$.
\end{proof}

%\old{
%For Brauer tree algebras, we can describe the effect of tilting mutation
%by the following result.
%
%\begin{example}\label{thm:mutation of Brauer tree} \cite{A1}
%Let $A$ be a Brauer tree algebra and $e_1\in A$ a primitive idempotent corresponding to the edge $1$. 
%Suppose that $B$ is the endomorphism algebra of indecomposable right mutation of $A$ with respect to $P_1$. 
%Then $B:=\End_{\K^{\rm{b}}(\proj{A})}(\mu_{e_1A}^-(A))$ is a Brauer tree algebra of which Brauer tree is given by the following rule, 
%where the multiplicities of vertices do not change.
%\begin{itemize}
%\item[(a)] if the edge $1$ is not the end;
%\[\begin{xy}
%(-45,0)="A"*{\xymatrix{
%\circ \ar@{-}[rd]^{}_{}="a1" &       &    &   & \circ \ar@{-}[ld]_{3}^{}="b1" \\
%                  & \circ \ar@{-}[rr]^{1}& & \circ & \\
%\circ \ar@{-}[ru]_{2}^{}="a2"  &       &&       & \circ \ar@{-}[lu]^{}_{}="b2"
%\ar@{.}@/^/"a2";"a1" \ar@{.}@/^/"b1";"b2"
%}},
%"A"+(45,0)="B"*{\xymatrix{
%\circ \ar@{-}[rd]^{}_{}="a1" & && & \circ \ar@{-}[ld]^{3}="b1" \ar@{-}[ddllll]_{1} \\
% & \circ & & \circ &  \\
%\circ \ar@{-}[ru]^{2}="a2" & && & \circ \ar@{-}[lu]^{}_{}="b2"
%\ar@{.}@/^/"a2";"a1" \ar@{.}@/^/"b1";"b2"
%}},
%"A"+(0,-30)*{},
%"A"+(7,-12)="a"*{},
%"A"+(40,-12)="b"*{},
%\ar "a";"b"^{1}
%\end{xy}\]
%\item[(b)] if the edge $1$ is the end;
%\[\begin{xy}
%(-35,0)="A"*{\xymatrix{
%              &         & & \circ \ar@{-}[ld]^{}="b1"  \\
%        \circ \ar@{-}[rr]^{1} & & \circ & \\
%&&&\circ \ar@{-}[lu]_{}="b2" 
% \ar@{.}@/^/"b1";"b2"
%}},
%"A"+(45,0)="B"*{\xymatrix{
%              &         & & \circ \ar@{-}[ld]^{}="b1" \ar@{-}[dlll]_{1} \\
%        \circ & & \circ & \\
%&&&\circ \ar@{-}[lu]_{}="b2" 
% \ar@{.}@/^/"b1";"b2"
%}},
%"A"+(0,-30)*{},
%"A"+(7,-12)="a"*{},
%"A"+(42,-12)="b"*{},
%\ar "a";"b"^{1}
%\end{xy}\]
%\end{itemize}
%\end{example}
%}

%\begin{corollary}
%For any idempotents $e\in A$, the Okuyama-Rickard complex $T$ with respect to $e$ is always a silting complex 
%in $\K^{\rm{b}}(\proj{A})$. 
%\end{corollary}
%However, in general, it is not tilting. 

%\begin{remark}\label{remark:tilt_Rickard}
%In particular, if $A$ is self-injective with the Nakayama functor $\nu =D\homo{A}{-}{A}$ and $\add{eA}=\add{\nu (eA)}$, 
%then $T$ is tilting. 
%\end{remark}

The partial order $\le$ on $\silt{\K^{\rm{b}}(\proj{A})}$
have the following easy interpretation for Okuyama-Rickard complexes. 

\begin{proposition}\label{prop:relation of OR}
Let $e, e'\in A$ be idempotents and $T, T'$ the Okuyama-Rickard complexes with respect to $e, e'$ respectively. 
Then the following conditions are equivalent.
\begin{itemize}
\item[(a)] $\add{eA} \supseteq \add{e'A}$.
\item[(b)] $T \ge T'$.   
\end{itemize}
\end{proposition}

\begin{proof}
Assume $\add{eA} \supseteq \add{e'A}$. 
Clearly we have $\Hom_{\K^{\rm{b}}(\proj{A})}(T,T'[>1])=0$. 
Take any morphism
\[\xymatrix{
T\ar[d]&0\ar[r]\ar[d]&P(eA(1-e)A)\oplus(1-e)A\ar^(0.8){p_{e}}[r]\ar^{f}[d]&eA\ar[d]\\
T'[1]&P(e'A(1-e')A)\oplus(1-e')A\ar^(0.7){p_{e'}}[r]&e'A\ar[r]&0
}\]
Since $P(eA(1-e)A)\oplus(1-e)A\in\add{(1-e)A}\subset\add{(1-e')A}$
and $p_{e'}$ is a right $\add{(1-e')A}$-approximation, $f$ factors through $p_{e'}$.
Thus we have $\Hom_{\K^{\rm{b}}(\proj{A})}(T,T'[1])=0$, so $T\ge T'$. 

Assume $T\ge T'$. Then we have $\Hom_{\K^{\rm{b}}(\proj{A})}(T,T'[1])=0$.
In particular, for any morphism $f:P(eA(1-e)A)\oplus(1-e)A\to e'A$,
there exist $s$ and $t$ such that $f=p_{e'}s+tp_e$.
Since both $p_e$ and $p_{e'}$ belongs to $J_{\modu{A}}$, any morphism
in $\homo{A}{P(eA(1-e)A)\oplus(1-e)A}{e'A}$ belongs to $J_{\modu{A}}$.
In particular $(1-e)A$ and $e'A$ do not have a non-zero common summand,
and we have  $\add{eA}\supseteq \add{e'A}$. 
\end{proof}

%\subsection{APR tilting modules and BB tilting modules}

%Let us recall APR tilting modules and BB tilting modules. 
%We say that an $A$-module $T$ is a (classical) \emph{tilting module} if
%\begin{itemize}
%\item the projective dimension ${\rm{pd}}(T)$ of $T$  is at most $1$,
%\item ${\rm{Ext}}_{A}^{1}(T, T)=0$, 
%\item $\delta (T)=\delta (A)$. 
%\end{itemize}

We end this section by observing a connection between BB and APR tilting modules and silting mutation.

\begin{definition}\label{def:APR and BB}
Let $e\in A$ be a primitive idempotent and $S$ the corresponding simple $A$-module.
Assume ${\rm{Ext}}_{A}^{1}(S, S)=0$, ${\rm{pd}}(\tau^{-1}S)\leq 1$ and that $S$ is not injective.
%\old{that the following conditions are satisfied.
%\begin{itemize}
%\item ${\rm{Ext}}_{A}^{1}(S, S)=0$,
%\item ${\rm{pd}}(\tau^{-1}S)\leq 1$,
%\item $S$ is not injective.
%\end{itemize}}
We define a \emph{BB tilting module} \cite{BB} with respect to $e$ by
\[T:=\tau^{-1}S\oplus (1-e)A\] 
%where $V$ runs non-isomorphic simple $A$-modules except a simple $A$-module $S$.
%and $P(X)$ is a projective cover of $X$ for an $A$-module $X$. 
%Then $T$ is a tilting module by \cite{BB}.
A BB tilting module is called an \emph{APR tilting module} \cite{APR} if $S$ is projective. 
\end{definition}

%We shall show that they are special classes of the following silting complexes. 
%
%\begin{definition}
%For any $P\in \proj{A}$, we say that $\mu^{+}_{P}(A)$ is the \emph{APR silting complexes with respect to $P$}. 
%\end{definition}

%Note that APR silting complexes are a duality of Okuyama-Rickard complexes. 
%\old{Let us give basic properties of BB and APR tilting modules in our context of silting mutation.}

\begin{theorem}\label{prop:APR is APR}
Let $T$ be the BB tilting module with respect to $e$.
\begin{itemize}
\item[(a)] $T$ is isomorphic to left mutation $\mu_{eA}^+(A)$ of $A$ with respect to $eA$.
%the APR silting complex with respect to a projective cover $P(S)$ of $S$. 
\item[(b)] $T$ is a tilting $A$-module of projective dimension at most one.
\end{itemize}
\end{theorem}

%To prove this statement, we need the following property. 

%\begin{lemma}\label{lem:silt=tiltmod}
%Let $T$ be a silting complex in $\K^{\rm{b}}(\proj{A})$. We assume that $T^{i}=0$ unless $-1 \leq i \leq 0$, 
%and denote the $i$-th cohomological functor by $H^{i}$. 
%Then the following are equivalent:
%\begin{enumerate}[$(1)$]
%\item $H^{0}(T)$ is a tilting module;
%\item $H^{-1}(T)=0$. 
%\end{enumerate}
%\end{lemma}
%\begin{proof}
%Assume that $H^{0}(T)$ is a tilting module. Then we can easily show that $H^{-1}(T)$ is a projective module 
%satisfying $\Hom_{A}(H^{-1}(T), H^{0}(T))=0$. 
%Since $H^{0}(T)$ is a tilting module, we have $H^{-1}(T)=0$. 
%Conversely, we assume $H^{-1}(T)=0$. 
%we can easily show ${\rm{pd}}(H^{0}(T))\leq 1$, 
%and note that $T$ is isomorphic to $H^{0}(T)$ in $\D^{\rm{b}}(\modu{A})$. 
%Since $T$ is a silting complex, 
%\[\def\arraystretch{1.5}
%\begin{array}{rl}
%{\rm{Ext}}_{A}^{1}(H^{0}(T), H^{0}(T)) &\simeq \Hom_{\D^{\rm{b}}(\modu{A})}(H^{0}(T), H^{0}(T)[1]) \\
%                               &\simeq \Hom_{\D^{\rm{b}}(\modu{A})}(T, T[1]) \\
%                               &=0
%\end{array}\]
%Finally, $T$ is a tilting complex because $\Hom_{\K^{\rm{b}}(\proj{A})}(T, T[-1])\simeq \Hom_{A}(H^{0}(T), H^{-1}(T))=0$. 
%Hence the number of non-isomorphic indecomposable summands of $H^{0}(T)$ is equal to 
%the number of non-isomorphic simple $A$-modules. Thus we can obtain that $H^{0}(T)$ is a tilting module. 
%\end{proof}

\begin{proof}
(a) Take a minimal injective resolution of the $A$-module $S$:
\[0\longrightarrow S\longrightarrow D(Ae)\stackrel{f}{\longrightarrow}I\]
Since $\Ext^1_A(S,S)=0$, we have $I\in\add{D(A(1-e))}$.
Thus $f$ is a minimal left $\add{D(A(1-e))}$-approximation.
Applying $\nu^{-1}$, we have an exact sequence
\[eA\xrightarrow{\nu^{-1}f}\nu^{-1}I\longrightarrow\tau^{-1}S\longrightarrow0\]
with a minimal left $(1-e)A$-approximation $\nu^{-1}f$.
Since ${\rm{pd}}(\tau^{-1}S)\leq 1$ and $S$ is not injective, we have that $\nu^{-1}f$ is injective. 
Thus we have $T\simeq\nu_{eA}^+(A)$.

(b) Since $\nu^{-1}f$ is injective, so is $\Hom_{\K^{\rm{b}}(\proj{A})}((1-e)A,\nu^{-1}f)$.
Thus the assertion follows from Theorem \ref{when silting is tilting}.
\end{proof}

%\old{
%\begin{example}
%Let $A$ be the path algebra of a quiver
%$\xymatrix{
%1 \ar[r]  & 2 \ar[r] & 3 
%}$. Let us calculate left mutation of $A$. 
%\begin{itemize}
%\item[(a)] We have $\mu_{P_1}^+(A)=P_1[1]\oplus P_2\oplus P_3$, which is not a tilting complex. 
%\item[(b)] We have $\mu_{P_2}^+(A)=P_1\oplus S_1 \oplus P_3$, which is a BB tilting module with respect to $e_2$.
%\item[(c)] We have $\mu_{P_3}^+(A)=P_1\oplus P_2\oplus S_2$, which is an APR tilting module with respect to $e_3$.
%\end{itemize}
%\end{example}
%}

%%%%%%%%%%%%%%%%%%%%%%%%%%%%%%%%%%%%%%%%%%%%%%%%%%%%%%%%%%%%%%%%%%%%%%%%%%%%%%%%%%%%%%%%%%%%%%%%%%%%%%%%%%%
%%%%%%%%%%%%%%%%%%%%%%%%%%%%%%%%%%%%%%%%%%%%%%%%%%%%%%%%%%%%%%%%%%%%%%%%%%%%%%%%%%%%%%%%%%%%%%%%%%%%%%%%%%%


\section{Transitivity for piecewise hereditary algebras}

The aim of this section is to prove the following result.

\begin{theorem}\label{hereditary case}
Let $\TT=\K^{\rm b}(\proj{A})$ where $A$ is either a hereditary algebra or a canonical algebra over a field $k$.
Then the action of iterated irreducible silting mutation on $\silt\TT$ is transitive.
\end{theorem}

The idea of our proof of Theorem \ref{hereditary case} is to compare silting objects
with exceptional sequences.

\emph{In the rest of this section let $\TT$ be a triangulated category satisfying
the following property.}

\begin{assumption}\label{total Hom-finite}
For any $X,Y\in\TT$ we have $\Hom_{\TT}(X,Y[\ell])=0$ for any $|\ell|\gg0$.
\end{assumption}

\begin{definition}
Let $\TT$ be a triangulated category.
We say that an object $X\in\TT$ is \emph{exceptional} if
$\Hom_{\TT}(X,X[\neq0])=0$ and $\End_{\TT}(X)$ is a division algebra. 

We say that a sequence $(X_1,\cdots,X_n)$ of exceptional objects in $\TT$
is an \emph{exceptional sequence} if
\[\Hom_{\TT}(X_i,X_j[\Z])=0\ \mbox{ for any }\ 1\le j<i\le n.\]
We say that an exceptional sequence is \emph{full} if
$\thick(\coprod_{i=1}^nX_i)=\TT$.
We denote by $\exc\TT$ the set of isomorphism classes of full exceptional sequences in $\TT$.
\end{definition}

Clearly $\Z^n$ acts on $\exc\TT$ by
\[(\ell_1,\cdots,\ell_n)(X_1,\cdots,X_n):=(X_1[\ell_1],\cdots,X_n[\ell_n]).\]
Let $B_n$ be the \emph{braid group} generated by $\sigma_1,\cdots,\sigma_{n-1}$
with relations
\begin{eqnarray*}
\sigma_i\sigma_j=\sigma_j\sigma_i&|i-j|>1,\\
\sigma_i\sigma_{i+1}\sigma_i=\sigma_{i+1}\sigma_i\sigma_{i+1}.&
\end{eqnarray*} 
Then $B_n$ acts on $\exc\TT$ \cite{GR} as follows:
For an exceptional sequence ${\mathbf X}:=(X_1,\cdots,X_n)$ and $1\le i<n$,
define objects $L_{X_{i+1}}X_i$ and $R_{X_i}X_{i+1}$ in $\TT$ by
\begin{eqnarray*}
\xymatrix{
X_i\ar[r]&\coprod_{\ell\in\Z}D\Hom_{\TT}(X_i,X_{i+1}[\ell])\otimes_kX_{i+1}[\ell]
\ar[r]&L_{X_{i+1}}X_i[1]\ar[r]&X_i[1]
}\\
\xymatrix{
X_{i+1}[-1]\ar[r]& R_{X_i}X_{i+1}\ar[r]&\coprod_{\ell\in\Z}\Hom_{\TT}(X_{i}[\ell],X_{i+1})\otimes_kX_i[\ell]\ar[r]&X_{i+1}
}\end{eqnarray*}
and put
\begin{eqnarray*}
\sigma_i^{-1}{\mathbf X}&:=&(X_1,\cdots,X_{i-1},X_{i+1},L_{X_{i+1}}X_i,X_{i+2},\cdots,X_n),\\
\sigma_i{\mathbf X}&:=&(X_1,\cdots,X_{i-1},R_{X_i}X_{i+1},X_{i+1},X_{i+2},\cdots,X_n).
\end{eqnarray*}
The following transitivity result of exceptional sequences is well-known.

\begin{theorem}\cite{C,Rin,KM}\label{exceptional transitivity}
Let $\TT=\K^{\rm b}(\proj{A})$ where $A$ is either a hereditary algebra
or a canonical algebra over a field $k$. Then $B_n\times\Z^n$ acts on $\exc\TT$ transitively.
\end{theorem}

We shall deduce Theorem \ref{hereditary case} from Theorem \ref{exceptional transitivity}.

We have the following connection between silting objects and exceptional sequences,
asserting that any full exceptional sequence gives rise to a silting object.

\begin{proposition}\label{exceptional to silting}
Let ${\mathbf X}=(X_1,\cdots,X_n)$ be a full exceptional sequence in $\TT$.
Then there exists $a\in\Z$ such that $X_1[\ell_1]\oplus\cdots\oplus X_n[\ell_n]$
is a silting object for any integers $\ell_1\cdots,\ell_n\in\Z$
satisfying $\ell_i+a\le\ell_{i+1}$ for any $1\le i<n$.
\end{proposition}

\begin{proof}
By Assumption \ref{total Hom-finite}, there exists $a\ge0$
such that $\Hom_{\TT}(X_i,X_j[>a])=0$ and any $1\le i,j\le n$.
We shall show that $a$ satisfies the desired condition.

Let $\ell_1\cdots,\ell_n$ be integers satisfying $\ell_i+a\le\ell_{i+1}$ for any $i$.
Fix $1\le i<j\le n$. Since ${\mathbf X}\in\exc\TT$, we have
$\Hom_{\TT}(X_j[\ell_j],X_i[\ell_i+\ell])=0$ for any $\ell\in\Z$
and $\Hom_{\TT}(X_i[\ell_i],X_i[\ell_i+\ell])=0$ for any $\ell\neq0$.
On the other hand we have $\Hom_{\TT}(X_i[\ell_i],X_j[\ell_j+\ell])=0$
for any $\ell>0$ by $\ell_j+\ell-\ell_i>a$ and our choice of $a$.
Thus we have the assertion.
\end{proof}

We have the following easy observation.

\begin{lemma}\label{very silting}
Let ${\mathbf X}=(X_1,\cdots,X_n)$ be a full exceptional sequence in $\TT$
such that $M=X_1\oplus\cdots\oplus X_n$ is a silting object.
For any integers $\ell_1\le\cdots\le \ell_n$ we have the following.
\begin{itemize}
\item[(1)] $X_1[\ell_1]\oplus\cdots\oplus X_n[\ell_n]$ is a silting object.
\item[(2)] If $\ell_1\ge0$, then $X_1\oplus\cdots\oplus X_n$ and
$X_1[\ell_1]\oplus\cdots\oplus X_n[\ell_n]$ are transitive under
iterated irreducible silting mutation.
\end{itemize}
\end{lemma}

\begin{proof}
(1) Fix $i\le j$. Since ${\mathbf X}\in\exc\TT$, we have
$\Hom_{\TT}(X_j[\ell_j],X_i[\ell_i+\ell])=0$ for any $\ell\neq0$.
Since $X_1\oplus\cdots\oplus X_n$ is a silting object and $\ell_i\le\ell_j$,
we have $\Hom_{\TT}(X_i[\ell_i],X_j[\ell_j+\ell])=0$ for any $\ell>0$.
Thus we have the assertion.

(2) We use the induction on $\ell_1+\cdots+\ell_n$.
Let $\ell_0:=0$. Assume $\ell_{i-1}<\ell_i$ for some $1\le i\le n$.
Then $\Hom_{\TT}(X_j[\ell_j],X_i[\ell_i])=0$ for any $j>i$ since ${\mathbf X}\in\exc\TT$,
and also for any $j<i$ since $M\in\silt\TT$ and $\ell_i-\ell_j>0$.
Consequently we have $\Hom_{\TT}(X_j[\ell_j],X_i[\ell_i])=0$ for any $j\neq i$.
Thus we have
\[\mu_{X_i[\ell_i]}^-(X_1[\ell_1]\oplus\cdots\oplus X_n[\ell_n])=
X_1[\ell_1]\oplus\cdots X_i[\ell_i-1]\cdots\oplus X_n[\ell_n].\]
Thus the assertion follows inductively.
%
%(ii) Next we assume $\ell_1<0$.
%
%By (i), we know that $X_1\oplus\cdots\oplus X_n$ and $X_1\oplus X_2[\ell_2-\ell_1]\oplus\cdots\oplus X_n[\ell_n-\ell_1]$
%are transitive under iterated irreducible silting mutation,
%and that $X_1[\ell_1]\oplus\cdots\oplus X_n[\ell_n]$ and $X_1\oplus X_2[\ell_2-\ell_1]\oplus\cdots\oplus X_n[\ell_n-\ell_1]$ 
%are transitive under iterated irreducible silting mutation.
%Thus the assertion follows.
\end{proof}

We have the following transitivity result.

\begin{proposition}\label{first transitivity}
Let $(X_1,\cdots,X_n)$ be a full exceptional sequence in $\TT$
and $\ell_1,\cdots,\ell_n\in\Z$,
If $X_1\oplus\cdots\oplus X_n$ and
$X_1[\ell_1]\oplus\cdots\oplus X_n[\ell_n]$ are silting objects,
then they are transitive under iterated irreducible silting mutation.
\end{proposition}

\begin{proof}
Choose integers $0\le m_1\le\cdots\le m_n$ satisfying
$0\le m_1-\ell_1\le\cdots\le m_n-\ell_n$.
By Lemma \ref{very silting} we have that
$X_1[m_1]\oplus\cdots\oplus X_n[m_n]$ is an iterated irreducible silting mutation of $X_1\oplus\cdots\oplus X_n$
(respectively, $X_1[\ell_1]\oplus\cdots\oplus X_n[\ell_n]$).
Thus we have the assertion.
\end{proof}

\emph{In the rest of this section let $\TT=\K^{\rm b}(\proj{A})$
where $A$ is either a hereditary algebra or a canonical algebra over a field.}
We can identify $\TT$ with $\D^{\rm b}(\AA)$ for a hereditary abelian
category $\AA$.
Any indecomposable object in $\TT$ is isomorphic to $X[\ell]$ for some
$X\in\AA$ and $\ell\in\Z$.
Moreover $\TT$ satisfies Assumption \ref{total Hom-finite}.

Let us start with the following preliminary results.

\begin{lemma}\label{preliminaries}\cite[IV.1.2, IV.1.5]{H}
\begin{itemize}
\item[(1)] Let $X,Y\in\AA$ be indecomposable objects satisfying $\Ext^1_{\AA}(Y,X)=0$.
Then any non-zero morphism $f\in\Hom_{\AA}(X,Y)$ is either injective or surjective.
\item[(2)] Let $X,Y\in\AA$ be exceptional objects.
If $\Hom_{\AA}(X,Y)\neq0$ and $\Ext^1_{\AA}(X,Y)\neq0$, then we have $\Ext^1_{\AA}(Y,X)\neq0$.
\end{itemize}
\end{lemma}

We have the following result asserting that $\RHom_{\AA}(X,Y)$ is a stalk complex.

\begin{lemma}\label{only one extension}
Let $X,Y\in\TT$ be indecomposable objects. If $M:=X\oplus Y$ satisfies
$\Hom_{\TT}(M,M[>0])=0$,
then there exists at most one $\ell\in\Z$ such that $\Hom_{\TT}(X,Y[\ell])\neq0$.
\end{lemma}

\begin{proof}
Without loss of generality, we can assume that $X$ and $Z:=Y[a]$ belong to $\AA$.
Since $\AA$ is hereditary, we have $\Hom_{\TT}(X,Z[\neq0,1])=0$.
Assume that $\Hom_{\TT}(X,Z)\neq0$ and $\Hom_{\TT}(X,Z[1])\neq0$.
By Lemma \ref{preliminaries}(2) we have $\Hom_{\TT}(Z,X[1])\neq0$.

Since $\Hom_{\TT}(X,Z[1])\neq0$, we have $\Hom_{\TT}(M,M[1+a])\neq0$.
Thus we have $a<0$.
On the other hand, since $\Hom_{\TT}(Z,X[1])\neq0$, we have
$\Hom_{\TT}(M[a],M[1])\neq0$. Thus we have $a>0$, a contradiction.
\end{proof}

We have the following information about existence of $\ell$ in Lemma \ref{only one extension}.

\begin{lemma}\label{no cycle}
Let $X_{n+1}=X_1,\cdots,X_n\in\TT$ be pairwise non-isomorphic indecomposable objects.
Assume that the following conditions are satisfied.
\begin{itemize}
\item[(i)] $M:=X_1\oplus\cdots\oplus X_n$ satisfies $\Hom_{\TT}(M,M[>0])=0$.
\item[(ii)] There exist integers $\ell_1,\cdots,\ell_n\in\Z$ such that
$\Hom_{\TT}(X_i,X_{i+1}[\ell_i])\neq0$ for any $1\le i\le n$.
\end{itemize}
Then we have $n=1$.
\end{lemma}

\begin{proof}
By (i) we have $\ell_i\le0$ for any $i$.
We take $a_i\in\Z$ such that $X_i\in\AA[a_i]$.
Since $\Hom_{\TT}(X_i,X_{i+1}[\ell_i])\neq0$, we have
$\Hom_{\TT}(\AA[a_i],\AA[a_{i+1}+\ell_i])\neq0$.
Thus $a_{i+1}-a_i+\ell_i$ is either $0$ or $1$. We have
\[0\ge \sum_{i=1}^n\ell_i=\sum_{i=1}^n(a_{i+1}-a_i+\ell_i)
=\sum_{i=1}^n(0\ \mbox{ or }\ 1),\]
which implies $\ell_1=\cdots=\ell_n=0$ and $a_1=\cdots=a_n$.
Thus we can assume that each $X_i$ belongs to $\AA$.

Assume $n>1$ and take a non-zero morphism $f_i:X_i\to X_{i+1}$ for any $1\le i\le n$.
Let $f_{n+1}:=f_1$. Then $f_i$ is not an isomorphism and either injective or surjective
by Proposition \ref{preliminaries}(1).
Also each $f_{i+1}f_i$ is either injective or surjective again by Proposition \ref{preliminaries}(1).
So it is impossible that $f_i$ is surjective and $f_{i+1}$ is injective at the same time.
It is easy to conclude that either all $f_i$ are injective or all $f_i$ are surjective.
Thus $f_n\cdots f_1:X_1\to X_1$ is not an isomorphism and either injective or surjective.
This is a contradiction.
\end{proof}

Now we show a certain converse of Proposition \ref{exceptional to silting}
asserting that any silting object gives rise to a full exceptional sequence.
It is also possible to show this observation by applying \cite[Theorem A]{AST}.

\begin{proposition}\label{silting to exceptional}
Let $M=X_1\oplus\cdots\oplus X_n$ be a basic silting object in $\TT$.
Then we can change indices of $X_1,\cdots,X_n$ such that
$(X_1,\cdots,X_n)$ is a full exceptional sequence in $\TT$.
\end{proposition}

\begin{proof}
For each $1\le i,j\le n$, we write $X_i\le X_j$ if there exists a sequence
$X_i=X_{i_1},X_{i_2},\cdots,X_{i_m}=X_j$ such that
$\Hom_{\TT}(X_{i_a},X_{i_{a+1}}[\ell_a])\neq0$ for some $\ell_a\in\Z$.
Clearly we have $X_i\le X_i$ and that $X_i\le X_j\le X_k$ implies $X_i\le X_k$.
By Lemma \ref{no cycle} we have that $X_i\le X_j\le X_i$ implies $i=j$.
This means that $X_1,\cdots,X_n$ forms a partially ordered set.
Thus we can change indices of $X_1,\cdots,X_n$ such that $X_i\le X_j$ implies $i\le j$.
Then $(X_1,\cdots,X_n)$ forms an exceptional sequence.
\end{proof}

For an exceptional sequence ${\mathbf X}=(X_1,\cdots,X_n)$, we let
\[[[{\mathbf X}]]:=X_1\oplus\cdots\oplus X_n\in\TT.\]
The following is a main inductive step in the proof of Theorem \ref{hereditary case}.

\begin{lemma}\label{exceptional mutation by silting mutation}
Let ${\mathbf X}=(X_1,\cdots,X_n)$ be a full exceptional sequence in $\TT$
such that $[[{\mathbf X}]]$ is a silting object.
For any $1\le i<n$, there exists $\mbox{\boldmath $\ell$}\in\Z^n$ (respectively, $\mbox{\boldmath $m$}\in\Z^n$) such that
$[[(\sigma_i,\mbox{\boldmath $\ell $}){\mathbf X}]]$ (respectively, $[[(\sigma_i^{-1},\mbox{\boldmath $m$}){\mathbf X}]]$)
is a iterated irreducible silting mutation of $[[{\mathbf X}]]$.
\end{lemma}

\begin{proof}
By Lemma \ref{only one extension} there exists at most one
$a\in\Z$ such that $\Hom_{\TT}(X_i,X_{i+1}[a])\neq0$.
Then $a$ must be non-positive since $[[{\mathbf X}]]$ is silting.
If such $a$ does not exist, we let $a:=0$.

By Assumption \ref{total Hom-finite} there exists $b\le0$ such that
\begin{eqnarray}\label{vanishing 1}
\Hom_{\TT}(X_j[b],X_{i+1}[\ell])=0&\mbox{ for any }\ 1\le j<i\ \mbox{ and }\ a\le\ell\le0.
\end{eqnarray}
By Lemma \ref{very silting} we have an exceptional sequence
\[{\mathbf Y}:=(X_1[b],\cdots,X_{i-1}[b],X_i,X_{i+1},X_{i+2},\cdots,X_n)\]
such that $[[{\mathbf Y}]]$ is an iterated irreducible silting mutation of $[[{\mathbf X}]]$.
Let
\begin{eqnarray*}
{\mathbf Z}&:=&(X_1[b],\cdots,X_{i-1}[b],X_i,X_{i+1}[a],X_{i+2},\cdots,X_n),\\
\sigma_i{\mathbf Z}&=&(X_1[b],\cdots,X_{i-1}[b],R_{X_i}(X_{i+1}[a]),X_i,X_{i+2},\cdots,X_n).
\end{eqnarray*}
By our choice of $a$ and \eqref{vanishing 1} we have
\[\mu_{i+1}^{-(-a)}([[{\mathbf Y}]])=[[{\mathbf Z}]]\ \mbox{ and }\ \mu_{i+1}^{-(-a+1)}([[{\mathbf Y}]])=[[\sigma_i{\mathbf Z}]].\]
Consequently $[[\sigma_i{\mathbf Z}]]=[[(\sigma_i,\mbox{\boldmath $\ell$}){\mathbf X}]]$
is an iterated irreducible silting mutation of $[[{\mathbf X}]]$,
where we put $\mbox{\boldmath $\ell$}:=(b,\cdots,b,0,a,0,\cdots,0)$.
\end{proof}

Now we are ready to prove Theorem \ref{hereditary case}.

Let $T=X_1\oplus\cdots\oplus X_n$ and $U=Y_1\oplus\cdots\oplus Y_n$
be basic silting objects in $\TT$.
By Proposition \ref{silting to exceptional} we can assume that
${\mathbf X}=(X_1,\cdots,X_n)$ and ${\mathbf Y}=(Y_1,\cdots,Y_n)$ are
full exceptional sequences.
By Theorem \ref{exceptional transitivity} there exists
$(\sigma,\mbox{\boldmath $\ell$})\in B_n\times\Z^n$ such that ${\mathbf Y}=(\sigma,\mbox{\boldmath $\ell$}){\mathbf X}$.
Writing down $\sigma\in B_n$ as a product of $\sigma_1^{\pm1},\cdots,\sigma_{n-1}^{\pm1}$ and applying Lemma \ref{exceptional mutation by silting mutation} repeatedly,
we have $\mbox{\boldmath $m$}\in\Z^n$ such that $[[(\sigma,\mbox{\boldmath $m$}){\mathbf X}]]$ is an iterated irreducible silting mutation of $[[{\mathbf X}]]$.
By Proposition \ref{first transitivity}
$[[(\sigma,\mbox{\boldmath $m$}){\mathbf X}]]=[[(0,\mbox{\boldmath $m$}-\mbox{\boldmath $\ell $}){\mathbf Y}]]$
is an iterated irreducible silting mutation of $[[{\mathbf Y}$]].
Thus we have the assertion.
\qed

%%%%%%%%%%%%%%%%%%%%%%%%%%%%%%%%%%%%%%%%%%%%%%%%%%%%%%%%%%%%%%%%%%%%%%%%%%%%%%%%%%%%%%%%%%%%%%%%%%%%%%%%%%%
%%%%%%%%%%%%%%%%%%%%%%%%%%%%%%%%%%%%%%%%%%%%%%%%%%%%%%%%%%%%%%%%%%%%%%%%%%%%%%%%%%%%%%%%%%%%%%%%%%%%%%%%%%%

%\section{Preliminaries on t-structures in triangulated categories}
%\label{Preliminaries on t-structures in triangulated categories}

\section{Silting subcategories in triangulated categories with coproducts} 
\label{Silting subcategories in triangulated categories with coproducts}

The aim of this section is to study silting subcategories for triangulated categories which have arbitrary coproducts. 
We need to modify the definition of silting subcategories (Definition \ref{silting coproducts}) from that in section \ref{section: silting subcategories}.
The advantage of this setting is that each set of compact objects gives rise to a torsion pair (Theorem \ref{torsion pair}),
which is not the case for the setting of section \ref{section: silting subcategories}.
As an application, we deduce a result of Hoshino-Kato-Miyachi \cite{HKM}, which associates a t-structure for each silting subcategory (Corollary \ref{HKM2}).
Moreover we show that this gives a one-to-one correspondence between silting subcategories and certain t-structures (Theorem \ref{correspondence}).
%Then we study hearts of certain t-structures including those associated with silting subcategories,
%and in particular realize them as functor categories (Theorem \ref{category P}).
We also deduce a result of Pauksztello \cite{P1,P2} on co-t-structures.

\medskip
Throughout this section, let $\TT$ be a triangulated category with arbitrary coproducts.

We say that an object $X\in\TT$ is \emph{compact} 
if $\Hom_{\TT}(X,-)$ commutes with arbitrary coproducts.
We denote by $\TT^{\rm c}$ the full subcategory consisting of compact objects in $\TT$.
We say that a subcategory $\MM$ of $\TT$ is \emph{compact} if $\MM\subset\TT^{\rm c}$,
\emph{generating} if $\MM[\Z]^{\perp_{\TT}}=0$, and
\emph{skeletally small} if isomorphism classes of objects in $\MM$ form a set.

\begin{definition}\label{silting coproducts}
We say that a subcategory $\MM$ of $\TT$ is \emph{silting} if 
$\Hom_{\TT}(\MM,\MM[>0])=0$ and $\MM$ is a skeletally small compact and generating subcategory of $\TT$.
%(i.e. $\MM\subset\TT^{\rm c}$ and $\MM[\Z]^{\perp_{\TT}}=0$).
\end{definition}



In this case $\MM$ is a silting subcategory of $\TT^{\rm c}$ in the sense of Definition \ref{definition of silting}
by the following Neeman's result:
%\begin{itemize}
%\item[(ii)'] $\thick\MM=\TT^{\rm c}$.
%\end{itemize}

\begin{proposition}\cite{N}\label{Neeman}
For any compact generating subcategory $\MM$ of $\TT$, we have $\thick\MM=\TT^{\rm c}$.
\end{proposition}

\subsection{Torsion pairs induced by sets of compact objects}\label{Torsion pairs in triangulated categories induced by compact objects}

The following main result in this section asserts that each set of compact objects gives a t-structure.

\begin{theorem}\label{torsion pair}
Let $\TT$ be a triangulated category with arbitrary coproducts,
and let $\CC$ be a set of objects in $\TT^{\rm c}$.
Then $({}^\perp(\CC^\perp),\CC^\perp)$ is a torsion pair.
\end{theorem}

We need the following general observation of homotopy colimit \cite{N}.

\begin{proposition}\label{hocolim}
Let $\TT$ be a triangulated category with arbitrary coproducts.
For a sequence 
\[X_0\xrightarrow{a_0}X_1\xrightarrow{a_1}X_2\xrightarrow{a_2}\cdots\]
we put
\[1-{\rm shift}:=\left(\begin{smallmatrix}
1&0&0&0&\cdots\\
-a_0&1&0&0&\cdots\\
0&-a_1&1&0&\cdots\\
0&0&-a_2&1&\cdots\\
\vdots&\vdots&\vdots&\vdots&\ddots
\end{smallmatrix}\right):\coprod_{i\ge0}X_i\to\coprod_{i\ge0}X_i\]
and take a triangle
\[\coprod_{i\ge0}X_i\xrightarrow{1-{\rm shift}}\coprod_{i\ge0}X_i\to Y\to\coprod_{i\ge0}X_i[1],\]
where $Y$ is called the \emph{homotopy colimit}.
Then we have the following isomorphism on $\TT^{\rm c}$:
\[\Hom_{\TT}(-,Y)\simeq{\rm colim}(\Hom_{\TT}(-,X_0)\xrightarrow{a_0\cdot}
\Hom_{\TT}(-,X_1)\xrightarrow{a_1\cdot}\Hom_{\TT}(-,X_2)\xrightarrow{a_2\cdot}\cdots).\]
\end{proposition}

We are ready to prove our Theorem \ref{torsion pair}.

%\noindent{\it Proof of Theorem \ref{torsion pair}.}
For any $X\in\TT$, we only have to construct a triangle
$S_X\to X\to U_X\to S_X[1]$ in $\TT$ with $U_X\in\CC^\perp$ and $S_X\in{}^\perp(\CC^\perp)$.
This is given as follows:

\begin{proposition}\label{construction}
Fix $X=X_0\in\TT$. For each $i\ge0$ we take a triangle
\begin{equation}\label{approximation}
\xymatrix{
C_i\ar[r]^{b_i}&X_i\ar[r]^{a_i}&X_{i+1}\ar[r]&C_i[1]
}\end{equation}
such that $b_i$ is a right $\Add\CC$-approximation.
This is possible since $\CC$ forms a set.
Thus we have a sequence
\[\xymatrix@R=.4cm{
C_0\ar_{b_0}[d]&C_1\ar_{b_1}[d]&C_2\ar_{b_2}[d]\\
X_0\ar^{a_0}[r]&X_1\ar^{a_1}[r]&X_2\ar^{a_2}[r]&\cdots
}\]
of morphisms.
We denote by $\iota_i$ the inclusion $X_i\to\coprod_{i\ge0}X_i$.
We take a triangle
\begin{equation}\label{hocolim triangle}
\coprod_{i\ge0}X_i\xrightarrow{1-{\rm shift}}\coprod_{i\ge0}X_i\xrightarrow{c}U_X\to\coprod_{i\ge0}X_i[1]
\end{equation}
and for the composition $d:=c\iota_0:X\to U_X$ we take a triangle
\begin{equation}\label{torsion pair triangle}
\xymatrix{S_X\ar[r]&X\ar[r]^d&U_X\ar[r]&S_X[1].}
\end{equation}
Then we have $S_X\in{}^\perp(\CC^\perp)$ and $U_X\in\CC^\perp$.
\end{proposition}

\begin{proof}
(i) We shall show $U_X\in\CC^\perp$.

For any $i\ge0$ we have an exact sequence
\[\Hom_{\TT}(\CC,C_i)\xrightarrow{b_i\cdot}\Hom_{\TT}(\CC,X_i)\xrightarrow{a_i\cdot}\Hom_{\TT}(\CC,X_{i+1}).\]
Since $b_i$ is a right $\Add\CC$-approximation, the left map is surjective, and hence the right map is zero.
Since $\CC$ consists of compact objects, we have by Proposition \ref{hocolim}
\[\Hom_{\TT}(\CC,U_X)\simeq{\rm colim}(\Hom_{\TT}(\CC,X_0)\xrightarrow{0}
\Hom_{\TT}(\CC,X_1)\xrightarrow{0}\cdots)=0.\]
Thus we have $U_X\in\CC^\perp$.

(ii) We shall show that $d:X\to U_X$ is a left $\CC^\perp$-approximation.

Fix any $f_0:X=X_0\to Y$ with $Y\in\CC^\perp$.
Since $f_0b_0\in\Hom_{\TT}(C_0,Y)=0$, there exists $f_1:X_1\to Y$ such that $f_0=f_1a_0$
by the triangle \eqref{approximation}.
Similarly we have $f_i:X_i\to Y$ for any $i>0$ such that
\begin{equation}\label{f=af}
f_{i-1}=f_ia_{i-1}.
\end{equation}
\[\xymatrix@R=.4cm{
C_0\ar_{b_0}[d]&C_1\ar_{b_1}[d]&C_2\ar_{b_2}[d]\\
X_0\ar^{a_0}[r]\ar_{f_0}[d]&X_1\ar^{a_1}[r]\ar_{f_1}[ld]&X_2\ar^{a_2}[r]\ar^{f_2}[lld]&\cdots\\
Y
}\]
Now we consider a morphism
\[f:=(f_0\ f_1\ \cdots)\in\Hom_{\TT}(\coprod_{i\ge0}X_i,Y).\]
By \eqref{f=af} we have $f(1-{\rm shift})=0$.
By the triangle \eqref{hocolim triangle} there exists $g\in\Hom_{\TT}(U_X,Y)$
such that $f=gc$.
In particular we have $f_0=f\iota_0=gd$.

(iii) We shall show that
$\Hom_{\TT}(\coprod_{i\ge0}X_i,\CC^\perp)\xrightarrow{\cdot(1-{\rm shift})}
\Hom_{\TT}(\coprod_{i\ge0}X_i,\CC^\perp)$ is surjective.

Let $Y\in\CC^\perp$.
Fix any $f=(f_0\ f_1\ \cdots)\in\Hom_{\TT}(\coprod_{i\ge0}X_i,Y)$.
Since $f_0b_0=0$, there exists $g_1\in\Hom_{\TT}(X_1,Y)$ such that $f_0=-g_1a_0$
by the triangle \eqref{approximation}.
Since $(f_1-g_1)b_1=0$, there exists $g_2\in\Hom_{\TT}(X_2,Y)$ such that
$f_1-g_1=-g_2a_1$.
Similarly we have $g_i:X_i\to Y$ for any $i\ge0$ which makes the following diagram commutative.
\[\xymatrix@R=.4cm{
C_i\ar[r]^{b_i}&X_i\ar[r]^{a_i}\ar[d]_{f_i-g_i}&X_{i+1}\ar[r]\ar[dl]^{-g_{i+1}}&C_i[1]\\
&Y}\]
Now the morphism $g:=(0\ g_1\ g_2\ \cdots)$ satisfies $f=g(1-{\rm shift})$.

(iv) We shall show that
$\Hom_{\TT}(U_X,\CC^\perp[1])\xrightarrow{\cdot d}\Hom_{\TT}(X,\CC^\perp[1])$ is injective.

Let $Y\in\CC^\perp[1]$. Assume that $f\in\Hom_{\TT}(U_X,Y)$ satisfies $fd=0$.
Put $g=(g_0\ g_1\ \cdots):=fc\in\Hom_{\TT}(\coprod_{i\ge0}X_i,Y)$.
Then we have $g_0=fc\iota_0=fd=0$.
\[\xymatrix@R=.4cm{
&&Y\\
\coprod_{i\ge0}X_i\ar[r]^{1-{\rm shift}}&\coprod_{i\ge0}X_i\ar[r]^{c}\ar[ur]^g&U_X\ar[r]\ar[u]^f&\coprod_{i\ge0}X_i[1]\\
S_X\ar[r]\ar[u]&X\ar[r]^d\ar[u]^{\iota_0}&U_X\ar[r]\ar@{=}[u]&S_X[1]\ar[u]
}\]
Assume $g_i=0$ for some $i\ge0$. Since
\[-g_{i+1}a_i=g_i-g_{i+1}a_i=g(1-{\rm shift})\iota_i=fc(1-{\rm shift})\iota_i=0,\]
we have that $g_{i+1}$ factors through $C_i[1]$ by the triangle \eqref{approximation}.
Since $Y\in\CC^\perp[1]$, we have $g_{i+1}=0$.

Inductively we have $g=0$, so $f$ factors through $U_X\to\coprod_{i\ge0}X_i[1]$.
Now using (iii) we have $f=0$.

(v) We shall show $S_X\in{}^\perp(\CC^\perp)$.

We have an exact sequence
\[\Hom_{\TT}(U_X,\CC^\perp)\xrightarrow{\cdot d}\Hom_{\TT}(X,\CC^\perp)\xrightarrow{}\Hom_{\TT}(S_X,\CC^\perp)\xrightarrow{}
\Hom_{\TT}(U_X[-1],\CC^\perp)\xrightarrow{\cdot d[-1]}\Hom_{\TT}(X[-1],\CC^\perp).\]
By (ii) the left map $(\cdot d)$ is surjective, and
by (iv) the right map $(\cdot d[-1])$ is injective.
Thus we have $\Hom_{\CC}(S_X,\CC^\perp)=0$, and we have the assertion.
\end{proof}

As special cases of Theorem \ref{torsion pair}, we have the following results due to Beligiannis-Reiten and Pauksztello.

\begin{corollary}\label{t and co-t}
Let $\TT$ be a triangulated category with arbitrary coproducts,
and $\CC$ a set of objects in $\TT^{\rm c}$.
\begin{itemize}
\item[(a)] \cite{BR} If $\CC[1]\subset\CC$, then $({}^\perp(\CC^\perp),\CC^\perp [1])$ is a t-structure.
\item[(b)] \cite{P2} If $\CC\subset\CC[1]$, then $({}^\perp(\CC^\perp)[1],\CC^\perp)$ is a co-t-structure.
\end{itemize}
\end{corollary}

Let us apply our results to more special cases, where we can describe the category ${}^\perp(\CC^\perp)$ in a more direct way.

The first application is the following result (b) of Hoshino-Kato-Miyachi \cite{HKM}.

\begin{corollary}\label{HKM2}
Let $\TT$ be a triangulated category with arbitrary coproducts and $\MM$ a skeletally small compact subcategory satisfying $\Hom_{\TT}(\MM,\MM[>0])=0$.
\begin{itemize}
\item[(a)] We have ${}^\perp(\MM[>0]^\perp)\subseteq\MM[\le0]^\perp$.
\item[(b)] If $\MM$ is a silting subcategory, then equality in (a) holds and $(\MM[<0]^\perp,\MM[>0]^\perp)$ is a t-structure of $\TT$.
\end{itemize}
\end{corollary}

\begin{proof}
Let $\CC:=\MM[>0]$. For any $X\in\TT$, we apply Proposition \ref{construction} to get a triangle %\eqref{torsion pair triangle} 
$S_X\to X\stackrel{d}{\to} U_X\to S_X[1]$ with $S_X\in{}^\perp(\CC^\perp)$ and $U_X\in\CC^\perp$.

Applying $\Hom_{\TT}(\MM[\le0],-)$ to \eqref{approximation} for any $i\ge0$, we have an isomorphism
\[(a_i\cdot):\Hom_{\TT}(\MM[\le0],X_i)\xrightarrow{\sim}\Hom_{\TT}(\MM[\le0],X_{i+1}).\]
Since $\MM[\le0]$ is compact, Proposition \ref{hocolim} gives an isomorphism
\begin{equation}\label{key iso}
(d\cdot):\Hom_{\TT}(\MM[\le0],X)\xrightarrow{\sim}\Hom_{\TT}(\MM[\le0],U_X).
\end{equation}

(a) Let $X\in{}^\perp(\CC^\perp)$. Then the morphism $d:X\to U_X$ is zero since $U_X\in\CC^\perp$.
Since \eqref{key iso} is an isomorphism which is zero, we have $\Hom_{\TT}(\MM[\le0],X)=0$.

(b) By Theorem \ref{torsion pair} we only have to show ${}^\perp(\CC^\perp)=\MM[\le0]^\perp$.
By (a), we only have to show ${}^\perp(\CC^\perp)\supset\MM[\le0]^\perp$.
Let $X\in\MM[\le0]^\perp$.
By \eqref{key iso}, we have $\Hom_{\TT}(\MM[\le0],U_X)\simeq\Hom_{\TT}(\MM[\le0],X)=0$.
Thus $U_X\in\CC^\perp\cap\MM[\le0]^\perp=\MM[\Z]^\perp=0$, and we have $X\simeq S_X\in{}^\perp(\CC^\perp)$.
\end{proof}

The second application is Corollary \ref{Pauksztello} below, which is
a version of a result of Pauksztello \cite{P1} on co-t-structures, where the assumption $\Hom_{\TT}(\MM,\MM[1])=0$ in \cite{P1} is dropped.

We say that a subcategory $\MM$ of $\TT$ is \emph{cosilting} if 
$\Hom_{\TT}(\MM,\MM[<0])=0$ and $\MM$ is a skeletally small compact and generating subcategory of $\TT$.

We say that an additive category $\MM$ is \emph{semisimple}
if  $\End_{\MM}(X)$ is a semisimple ring for any $X\in\MM$.

\begin{corollary}\label{Pauksztello}
Let $\TT$ be a triangulated category with arbitrary coproducts and $\MM$ a skeletally small compact and semisimple subcategory satisfying $\Hom_{\TT}(\MM,\MM[<0])=0$.
\begin{itemize}
\item[(a)] We have ${}^\perp(\MM[<0]^\perp)\subseteq\MM[\ge0]^\perp$.
\item[(b)] If $\MM$ is a cosilting subcategory, then equality in (a) holds and $(\MM[>0]^\perp,\MM[<0]^\perp)$ is a co-t-structure of $\TT$.
\end{itemize}
\end{corollary}

\begin{proof}
Let $\CC:=\MM[<0]$. 
%\old{By Theorem \ref{torsion pair} we only have to show
%${}^\perp(\CC^\perp)=\MM[>0]^\perp[-1]=\MM[\ge0]^\perp$.}
For any $X\in\TT$, we apply Proposition \ref{construction} with an additional assumption
that $\Hom_{\TT}(\MM[-1],C_i)\xrightarrow{b_i\cdot}\Hom_{\TT}(\MM[-1],X_i)$ is an isomorphism for any $i\ge0$. This is possible since $\MM$ is semisimple.
Applying $\Hom_{\TT}(\MM[\ge0],-)$ to \eqref{approximation} for any $i\ge0$, we have an isomorphism
\[(a_i\cdot):\Hom_{\TT}(\MM[\ge0],X_i)\xrightarrow{\sim}\Hom_{\TT}(\MM[\ge0],X_{i+1}).\]
Since $\MM[\ge0]$ is compact, Proposition \ref{hocolim} implies that our triangle $S_X\to X\stackrel{d}{\to} U_X\to S_X[1]$ induces an isomorphism
\begin{equation}\label{key iso2}
(d\cdot):\Hom_{\TT}(\MM[\ge0],X)\xrightarrow{\sim}\Hom_{\TT}(\MM[\ge0],U_X).
\end{equation}

(a) Let $X\in{}^\perp(\CC^\perp)$. Then the morphism $d:X\to U_X$ is zero since $U_X\in\CC^\perp$.
Since \eqref{key iso2} is an isomorphism which is zero, we have $\Hom_{\TT}(\MM[\ge0],X)=0$.

(b) By Theorem \ref{torsion pair} we only have to show ${}^\perp(\CC^\perp)\supset\MM[\ge0]^\perp$.
Let $X\in\MM[\ge0]^\perp$. By \eqref{key iso2}, we have $\Hom_{\TT}(\MM[\ge0],U_X)\simeq\Hom_{\TT}(\MM[\ge0],X)=0$.
Thus $U_X\in\CC^\perp\cap\MM[\ge0]^\perp=\MM[\Z]^\perp=0$, and we have $X\simeq S_X\in{}^\perp(\CC^\perp)$.
\end{proof}


%%%%%%%%%%%%%%%%%%%%%%%%%%%%%%%%%%%%%%%%%%%%%%%%%%%%%%%%%%%%%%%%%%%%%%%%%%%%%%%%%%%%%%%%%%%%%%%%%%%%%%%%%%%
%%%%%%%%%%%%%%%%%%%%%%%%%%%%%%%%%%%%%%%%%%%%%%%%%%%%%%%%%%%%%%%%%%%%%%%%%%%%%%%%%%%%%%%%%%%%%%%%%%%%%%%%%%%

\subsection{The correspondence between silting subcategories and t-structures}\label{The correspondence between silting subcategories and t-structures}

Let $\TT$ be a triangulated category with arbitrary coproducts.
For a silting subcategory $\MM$ of $\TT$, we observed in Corollary \ref{HKM2} that we have a t-structure 
$(\TT_{\MM}^{\le0},\TT_{\MM}^{\ge0})$ given by
\begin{eqnarray*}
\TT_{\MM}^{\le0}:={}^\perp(\MM[\ge0]^\perp)=\MM[<0]^\perp\ \mbox{ and }\ \TT_{\MM}^{\ge0}:=\MM[>0]^\perp.
\end{eqnarray*}
%If $\MM$ is a silting subcategory of $\TT$, we have $\MM\subset\TT_{\MM}^{\le0}$.
%Moreover if $\MM$ is a tilting subcategory of $\TT$, we also have $\MM\subset\TT_{\MM}^{\ge0}$.
The aim of this subsection is to characterize these t-structures obtained from silting subcategories.
The following definition is suggested by Proposition \ref{recover}.

\begin{definition}
Let $(\TT^{\le0},\TT^{\ge0})$ be a t-structure in $\TT$ with the heart $\TT^0$.
Define a subcategory of $\TT$ by
\[\MM=\MM(\TT^{\le0},\TT^{\ge0}):=\TT^{\le0}\cap{}^\perp(\TT^{<0})\cap\TT^{\rm c}.\]
\begin{itemize}
\item[(a)] We say that the t-structure $(\TT^{\le0},\TT^{\ge0})$ is \emph{silting} if $\MM$ is skeletally small and generating.
\item[(b)] We say that a silting t-structure is \emph{tilting} if $\MM\subset\TT^0$.
\end{itemize}
\end{definition}

The names of these t-structures are explained by the following our main results in this subsection.

\begin{theorem}\label{correspondence}
\begin{itemize}
\item[(a)] We have mutually inverse bijections 
\[\MM\mapsto(\TT_{\MM}^{\le0},\TT_{\MM}^{\ge0})\ \mbox{ and }\ (\TT^{\le0},\TT^{\ge0})\mapsto\MM(\TT^{\le0},\TT^{\ge0})\]
between silting subcategories of $\TT$ and silting t-structures of $\TT$.
\item[(b)] These induce bijections between tilting subcategories of $\TT$ and tilting t-structures of $\TT$.
\end{itemize}
\end{theorem}

We need the following observation.

\begin{lemma}
Let $(\TT^{\le0},\TT^{\ge0})$ be an arbitrary t-structure in $\TT$ and $\MM:=\MM(\TT^{\le0},\TT^{\ge0})$. Then
\begin{eqnarray}\label{from M to T 1}
\Hom_{\TT}(\MM,\TT^{\le0}[>0])=0,\\ \label{from M to T 2}
\Hom_{\TT}(\MM,\TT^{\ge0}[<0])=0,\\ \label{vanishing of M}
\Hom_{\TT}(\MM,\MM[>0])=0.
\end{eqnarray}
\end{lemma}

\begin{proof}
By $\MM\subset{}^\perp(\TT^{<0}$), we have \eqref{from M to T 1}.
By $\MM\subset\TT^{\le0}$, we have \eqref{from M to T 2}.
By $\MM\subset\TT^{\le0}$ and \eqref{from M to T 1}, we have \eqref{vanishing of M}.
\end{proof}

%\begin{lemma}\label{M recovers T 1}
%If $(\TT^{\le0},\TT^{\ge0})$ is a silting t-structure in $\TT$ and $\MM:=\MM(\TT^{\le0},\TT^{\ge0})$, then $\TT^{\ge0}=\TT_{\MM}^{\ge0}$.
%\end{lemma}
%
%\begin{proof}
%By \eqref{from M to T 2} we have $\TT^{\ge0}\subset\TT_{\MM}^{\ge0}$.
%We shall show the converse.
%For any $X\in\TT_{\MM}^{\ge0}$, take a triangle
%\[X^{<0}\to X\to X^{\ge0}\to X^{<0}[1]\]
%with $X^{<0}\in\TT^{<0}$ and $X^{\ge0}\in\TT^{\ge0}$.
%By \eqref{from M to T 1} we have $\Hom_{\TT}(\MM,X^{<0}[\ge0])=0$.
%On the other hand, for any $i<0$ we have an exact sequence
%\[0\stackrel{\eqref{from M to T 2}}{=}\Hom_{\TT}(\MM,X^{\ge0}[<-1])\to\Hom_{\TT}(\MM,X^{<0}[<0])\to\Hom_{\TT}(\MM,X[<0])\stackrel{X\in\TT_{\MM}^{\ge0}}{=}0.\]
%Thus we have $\Hom_{\TT}(\MM,X^{<0}[<0])=0$.
%Consequently $\Hom_{\TT}(\MM,X^{<0}[\Z])=0$ holds, so we have $X^{<0}=0$ and $X\simeq X^{\ge0}\in\TT^{\ge0}$.
%\end{proof}
%
%\begin{lemma}\label{M recovers T 2}
%If $(\TT^{\le0},\TT^{\ge0})$ is a silting t-structure in $\TT$ and $\MM:=\MM(\TT^{\le0},\TT^{\ge0})$, then $\TT^{\le0}=\TT_{\MM}^{\le0}$.
%\end{lemma}
%
%\begin{proof}
%By \eqref{from M to T 1} we have $\TT^{\le0}\subset\TT_{\MM}^{\le0}$.
%We shall the converse.
%or any $X\in\TT_{\MM}^{\le0}$, take a triangle
%\[X^{\le0}\to X\to X^{>0}\to X^{\le0}[1]\]
%with $X^{\le0}\in\TT^{\le0}$ and $X^{>0}\in\TT^{>0}$.
%By \eqref{from M to T 2} we have $\Hom_{\TT}(\MM,X^{>0}[i])=0$ for any $i\le0$.
%On the other hand, for any $i>0$ we have an exact sequence
%\[[0\stackrel{X\in\TT_{\MM}^{\le0}}{=}\Hom_{\TT}(\MM,X[i])]\to\Hom_{\TT}(\MM,X^{>0}[i])\to[\Hom_{\TT}(\MM,X^{\le0}[i+1])\stackrel{\eqref{from M to T 1}}{=}0].\]
%Thus we have $\Hom_{\TT}(\MM,X^{>0}[i])=0$ for any $i>0$.
%Consequently we have $\Hom_{\TT}(\MM,X^{>0}[i])=0$ for any $i\in\Z$.
%This implies $X^{>0}=0$ and $X\simeq X^{\le0}\in\TT^{\le0}$.
%\end{proof}

%\begin{lemma}\label{M and N}
%Let $\MM$ and $\NN$ be silting subcategories of $\TT$ satisfying $\MM\subset\NN$.
%Then we have $\MM=\NN$.
%\end{lemma}

Now we are ready to prove Theorem \ref{correspondence}.

(a)(i) Let $(\TT^{\le0},\TT^{\ge0})$ be a silting t-structure and $\MM:=\MM(\TT^{\le0},\TT^{\ge0})$.
Then $\MM$ is skeletally small compact and generating by definition. Thus $\MM$ is a silting subcategory by \eqref{vanishing of M}. 

We have $\TT^{\ge0}\subset\TT_{\MM}^{\ge0}$ by \eqref{from M to T 2}, and we have $\TT^{\le0}\subset\TT_{\MM}^{\le0}$ by \eqref{from M to T 1}.
Moreover we have
\begin{eqnarray*}
\TT^{\le0}={}^\perp(\TT^{>0})\supset{}^\perp(\TT_{\MM}^{>0})=\TT_{\MM}^{\le0},\\
\TT^{\ge0}=(\TT^{<0})^\perp\supset(\TT_{\MM}^{<0})^\perp=\TT_{\MM}^{\ge0}
\end{eqnarray*}
by Corollary \ref{HKM2}, so $(\TT^{\le0},\TT^{\ge0})=(\TT_{\MM}^{\le0},\TT_{\MM}^{\ge0})$ holds.

(ii) Let $\MM$ be a silting subcategory of $\TT$.
Then $(\TT_{\MM}^{\le0},\TT_{\MM}^{\ge0})$ is a t-structure of $\TT$ by Corollary \ref{HKM2}.
Let $\NN:=\MM(\TT_{\MM}^{\le0},\TT_{\MM}^{\ge0})$.
Then $\NN$ is skeletally small since $\TT^{\rm c}$ is skeletally small by Remark \ref{T is small}.
We have
\begin{eqnarray}\label{N 1}
\Hom_{\TT}(\NN,\NN[>0])=0
\end{eqnarray}
by \eqref{vanishing of M}.
Since $\MM$ is silting, we have $\MM\subset\MM[<0]^\perp=\TT_{\MM}^{\le0}$.
Since $\TT_{\MM}^{<0}=\MM[\le0]^\perp$, we have
$\Hom_{\TT}(\MM,\TT_{\MM}^{<0})=0$ and so $\MM\subset{}^\perp(\TT_{\MM}^{<0})$.
Consequently we have
\begin{equation}\label{M is in N}
\MM\subset\TT_{\MM}^{\le0}\cap{}^\perp(\TT_{\MM}^{<0})\cap\TT^{\rm c}=\NN.
\end{equation}
%If $X\in\TT$ satisfies $\Hom_{\TT}(\NN,X[\Z])=0$, then $\Hom_{\TT}(\MM,X[\Z])=0$. Thus we have $X=0$.
Since $\MM$ is generating, so is $\NN$.
By \eqref{N 1}, we have that $\NN$ is a silting subcategory of $\TT$.
In particular, $(\TT_{\MM}^{\le0},\TT_{\MM}^{\ge0})$ is a silting t-structure.

Now $\MM$ and $\NN$ are silting subcategories of $\TT^{\rm c}$ in the sense of Definition \ref{definition of silting} by Proposition \ref{Neeman}.
By \eqref{M is in N} and Theorem \ref{no inclusion} we have $\MM=\NN$.
Thus our correspondences are mutually inverse.

(b) Let $\MM$ be a tilting subcategory of $\TT$.
Then $\MM$ is contained in both $\TT_{\MM}^{\le0}$ and $\TT_{\MM}^{\ge0}$,
so we have $\MM\subset\TT_{\MM}^0$.
Since $\MM(\TT_{\MM}^{\le0},\TT_{\MM}^{\ge0})$ equals to $\MM$,
we have that $(\TT_{\MM}^{\le0},\TT_{\MM}^{\ge0})$ is a tilting t-structure.

Let $(\TT^{\le0},\TT^{\ge0})$ be a tilting t-structure and $\MM:=\MM(\TT^{\le0},\TT^{\ge0})$. Since $\MM\subset\TT^0$, we have
\[\Hom_{\TT}(\MM,\MM[<0])\subset\Hom_{\TT}(\TT^{\le0},\TT^{>0})=0.\]
Thus $\MM$ is a tilting subcategory of $\TT$.
\qed

%%%%%%%%%%%%%%%%%%%%%%%%%%%%%%%%%%%%%%%%%%%%%%%%%%%%%%%%%%%%%%%%%%%%%%%%%%%%%%%%%%%%%%%%%%%%%%%%%%%%%%%%%%%
%%%%%%%%%%%%%%%%%%%%%%%%%%%%%%%%%%%%%%%%%%%%%%%%%%%%%%%%%%%%%%%%%%%%%%%%%%%%%%%%%%%%%%%%%%%%%%%%%%%%%%%%%%%

\begin{thebibliography}{15}
%\bibitem[AH]{AH} H. Abe, M. Hoshino, \emph{On derived equivalences for selfinjective algebras}. 
%Comm. in algebra {\bf{34}} (2006), no. 12, 4441--4452.
%\bibitem[AH2]{AH2} H. Abe, M. Hoshino, Original manuscript of tilting mutation. 
%\bibitem[A1]{A1} T. Aihara, \emph{Mutating Brauer trees}, arXiv:1009.3210.
\bibitem[A]{A2} T. Aihara, \emph{Silting mutation for self-injective algebras}, arXiv:1012.3265.
\bibitem[AGI]{AGI} T. Aihara, J. Grant, O. Iyama, in preparation.
%\bibitem[ARi]{ARi} S. Al-nofayee, J. Rickard, \emph{Rigidity of tilting complexes and derived equivalence 
%for self-injective algebras}. preprint. 
\bibitem[AST]{AST} I. Assem, M. Souto Salorio, S. Trepode,
\emph{Ext-projectives in suspended subcategories},
J. Pure Appl. Algebra 212 (2008), no. 2, 423--434.
\bibitem[ASS]{ASS} I. Assem, D. Simson, A. Skowronski,
\emph{Elements of the representation theory of associative algebras. Vol. 1. Techniques of representation theory},
London Mathematical Society Student Texts, 65. Cambridge University Press, Cambridge, 2006.
\bibitem[APR]{APR} M. Auslander, M. I. Platzeck, I. Reiten, \emph{Coxeter functors without diagrams}, 
Trans. Amer. Math. Soc. 250 (1979), 1--46.
\bibitem[ARS]{ARS} M. Auslander, I. Reiten, S. O. Smalo, \emph{Representation theory of Artin algebras},
Cambridge Studies in Advanced Mathematics, 36. Cambridge University Press, Cambridge, 1997.
\bibitem[BR]{BR} A. Beligiannis, I. Reiten, \emph{Homological and homotopical aspects of torsion theories},
Mem. Amer. Math. Soc. 188 (2007), no. 883,
\bibitem[BBD]{BBD} A.A. Beilinson, J. Bernstein, P. Deligne, \emph{Faisceaux pervers}, Analysis and topology on singular spaces, I (Luminy, 1981), 5--171, Asterisque, 100, Soc. Math. France, Paris, 1982.
\bibitem[BGP]{BGP} I. N. Bernstein, I. M. Gelfand, V. A. Ponomarev, \emph{Coxeter functors, and Gabriel's theorem}.
Uspehi Mat. Nauk 28 (1973), no. 2(170), 19--33.
%\bibitem[B]{B} K. Bongartz, Tilted algebras. \emph{Representations of algebras (Puebla, 1980)}, 
%Lecture Notes in Math. Vol. 903. Springer, Berlin-New York, 1981, 26--38.
\bibitem[BB]{BB} S. Brenner, M. C. R. Butler, \emph{Generalizations of the Bern\v{s}te\v{i}n Gel'fand Ponomarev reflection functors}. 
Representation theory, II , pp. 103--169, Lecture Notes in Math., 832, Springer, Berlin-New York, 1980.
\bibitem[BIRS]{BIRS} A. B. Buan, O. Iyama, I. Reiten, J. Scott,
\emph{Cluster structures for 2-Calabi-Yau categories and unipotent groups},
Compos. Math. 145 (2009), no. 4, 1035--1079. 
\bibitem[BMRRT]{BMRRT} A. B. Buan, R. Marsh, M. Reineke, I. Reiten, G. Todorov,
\emph{Tilting theory and cluster combinatorics}, Adv. Math. 204 (2006), no. 2, 572--618.
\bibitem[BRT]{BRT} A. B. Buan, I. Reiten, H. Thomas, \emph{Three kinds of mutation}, arXiv:1005.0276.
\bibitem[BIKR]{BIKR} I. Burban, O. Iyama, B. Keller, I. Reiten,
\emph{Cluster tilting for one-dimensional hypersurface singularities},
Adv. Math. 217 (2008), no. 6, 2443--2484.
\bibitem[C]{C} W. Crawley-Boevey, \emph{Exceptional sequences of representations of quivers}, 
Representations of algebras (Ottawa, ON, 1992), 117--124, CMS Conf. Proc., 14, Amer. Math. Soc., Providence, RI, 1993. 
\bibitem[GR]{GR} A. L. Gorodentsev, A. N. Rudakov, \emph{Exceptional vector bundles on projective spaces}, 
Duke Math. J. 54 (1987), no. 1, 115--130.
\bibitem[H]{H} D. Happel, \emph{Triangulated categories in the representation theory of finite-dimensional algebras}, 
London Mathematical Society Lecture Note Series, 119. Cambridge University Press, Cambridge, 1988.
\bibitem[HU]{HU} D. Happel, L. Unger, \emph{On a partial order of tilting modules}. 
Algebr. Represent. Theory 8 (2005), no. 2, 147--156.
%\bibitem[Ho]{Ho} M. Hoshino, \emph{Tilting modules and torsion theories}. 
%Bull. London Math Soc., {\bf{14}} (1982) No. 4, 334--336. 
\bibitem[HK]{HK} M. Hoshino, Y. Kato, \emph{Tilting complexes defined by idempotents}, Comm. Algebra 30 (2002), no. 1, 83--100.
\bibitem[HKM]{HKM} M. Hoshino, Y. Kato, J. Miyachi, \emph{On $t$-structures and torsion theories induced by compact objects}, 
J. Pure Appl. Algebra 167 (2002), no. 1, 15--35. 
\bibitem[HX]{HX} W. Hu, C, Xi,
\emph{$D$-split sequences and derived equivalences}, Adv. Math. 227 (2011), no. 1, 292--318.
\bibitem[IO]{IO} O. Iyama, S. Oppermann,
\emph{n-representation-finite algebras and n-APR tilting},
arXiv:0909.0593, to appear in Trans. Amer. Math. Soc.
\bibitem[IR]{IR} O. Iyama, I. Reiten, 
\emph{Fomin-Zelevinsky mutation and tilting modules over Calabi-Yau algebras},
Amer. J. Math. 130 (2008), no. 4, 1087--1149.
\bibitem[IW]{IW} O. Iyama, M. Wemyss, \emph{Auslander-Reiten duality and maximal modifications for non-isolated singularities},
arXiv:1007.1296.
\bibitem[IY]{IY} O. Iyama, Y. Yoshino, \emph{Mutation in triangulated categories and rigid Cohen-Macaulay modules}, 
Invent. Math. 172 (2008), no. 1, 117--168.
\bibitem[K]{K} B. Keller, \emph{Deriving DG categories},
Ann. Sci. Ecole Norm. Sup. (4) 27 (1994), no. 1, 63--102. 
\bibitem[KV]{KV} B. Keller, D. Vossieck, \emph{Aisles in derived categories},
Deuxieme Contact Franco-Belge en Algebre (Faulx-les-Tombes, 1987).
Bull. Soc. Math. Belg. Ser. A 40 (1988), no. 2, 239--253.
\bibitem[KY]{KY}, B. Keller, D. Yang, \emph{Derived equivalences from mutations of quivers with potential},
Adv. Math. 226 (2011), no. 3, 2118--2168.
\bibitem[KM]{KM} D. Kussin, H. Meltzer, \emph{The braid group action for exceptional curves}, 
Arch. Math. (Basel) 79 (2002), no. 5, 335--344. 
\bibitem[L]{L} S. Ladkani, \emph{Perverse equivalences, BB-tilting, mutations and applications}, arXiv:1001.4765.
\bibitem[M]{M} J. Miyachi, \emph{Localization of triangulated categories and derived categories}, J. Algebra 141 (1991), no. 2, 463--483.
\bibitem[N]{N} A. Neeman, \emph{The connection between the $K$-theory localization theorem of Thomason, 
Trobaugh and Yao and the smashing subcategories of Bousfield and Ravenel},
Ann. Sci. Ecole Norm. Sup. (4) 25 (1992), no. 5, 547--566.
\bibitem[O]{O} T. Okuyama, \emph{Some examples of derived equivalent blocks of finite groups}. 
preprint, 1998. 
%\bibitem[Ok2]{Ok2} Original manuscript of tilting mutation.
\bibitem[P1]{P1} D. Pauksztello, \emph{Compact corigid objects in triangulated categories and co-$t$-structures},  
Cent. Eur. J. Math.  6  (2008),  no. 1, 25--42. 
\bibitem[P2]{P2} D. Pauksztello, \emph{A note on compactly generated co-t-structures}, arXiv:1006.5347.
\bibitem[Ric1]{Ric1} J. Rickard, \emph{Morita theory for derived categories}, J. London Math. Soc. (2) 39 (1989), no. 3, 436--456. 
\bibitem[Ric2]{Ric2} J. Rickard, \emph{Derived categories and stable equivalence}. 
J. Pure. Appl. 61 (1989), 303--317.
%\bibitem[Ri2]{Ri2} Original manuscript of tilting mutation.
\bibitem[RS]{RS} C. Riedtmann, A. Schofield,
\emph{On a simplicial complex associated with tilting modules},
Comment. Math. Helv. 66 (1991), no. 1, 70--78. 
\bibitem[Rin]{Rin} C. M. Ringel, \emph{The braid group action on the set of exceptional sequences of a hereditary Artin algebra}, 
Abelian group theory and related topics (Oberwolfach, 1993), 339--352, Contemp. Math., 171, Amer. Math. Soc., Providence, RI, 1994.
\bibitem[SY]{SY} Y. Sekiya, K. Yamaura, \emph{Tilting theoretical approach to moduli spaces over preprojective algebras}, arXiv:1008.4451.
%\bibitem[S]{S} S. O. Smal$\phi $, \emph{Torsion theories and tilting modules}. 
%Bull. London Math. Soc., {\bf{16}} (1984) No. 5, 518--522. 
\bibitem[W]{W} J. Wei, \emph{Semi-tilting complexes}, preprint.
\end{thebibliography}
\end{document}






%%%%%%%%%%%%%%%%%%%%%%%%%%%%%%%%%%%%%%%%%%%%%%%%%%%%%%%%%%%%%%%%%%%%%%%%%%%%%%%%%%%%%%%%%%%%%%%%%%%%%%%%%%%
%%%%%%%%%%%%%%%%%%%%%%%%%%%%%%%%%%%%%%%%%%%%%%%%%%%%%%%%%%%%%%%%%%%%%%%%%%%%%%%%%%%%%%%%%%%%%%%%%%%%%%%%%%%

\subsection{The number of indecomposable summands of silting objects}

In this section, let $\TT$ be a triangulated category with arbitrary coproducts, 
and we assume that $\TT$ is $k$-linear for a field $k$, 
$\TT^{\rm{c}}$ is Krull-Schmidt and $\TT^{\rm{c}}/J_{\TT^{\rm{c}}}$ is Hom-finite. 

For any $X\in \TT^{\rm{c}}$, we denote by $\delta (X)$ the number of non-isomorphic indecomposable direct summands of $X$. 
The aim of this section is to show the following result. 

\begin{theorem}\label{thm:number of indecomposable direct summands}
For any silting objects $M,N\in \TT$, we have $\delta (M)=\delta (N)$. 
\end{theorem}

To prove this statement, the following full subcategory is useful. 

\begin{definition}
We define a full subcategory of $\TT$ by 
\[\TT^{\rm{b}}=\{X\in \TT \ |\ \sum_{i\in \Bbb{Z}}{\rm{dim}}_{k}\homo{\TT}{Y}{X[i]}<\infty \ \mbox{for any } Y\in \TT^{\rm{c}}\}.\]
Obviously, $\TT^{\rm{b}}$ is a thick subcategory of $\TT$. 
\end{definition}

We recall the Grothendieck group of $\TT^{\rm{b}}$. 

\begin{definition}
We define $\Bbb{Z}(\TT^{\rm{b}})$ to be the free abelian group on isomorphism classes of objects of $\TT^{\rm{b}}$. 
An object $X\in \TT^{\rm{b}}$, viewed as an element in $\Bbb{Z}(\TT^{\rm{b}})$, is denoted $[X]$.  
We define a subgroup $R(\TT^{\rm{b}})$ of $\Bbb{Z}(\TT^{\rm{b}})$ to be the subgroup generated by sums $[X]-[Y]+[Z]$ 
where $X\to Y\to Z\to X[1]$ is a triangle in $\TT^{\rm{b}}$. 
Then we define a group, namely the \emph{Grothendieck group} of $\TT^{\rm{b}}$, 
by $K_{0}(\TT^{\rm{b}})=\Bbb{Z}(\TT^{\rm{b}})/R(\TT^{\rm{b}})$. 
\end{definition}

Theorem \ref{thm:number of indecomposable direct summands} is an immediate consequence of the following result.

\begin{theorem}\label{rank and delta}
Assume that $\TT$ has a silting object $M$. Then $K_{0}(\TT^{\rm{b}})$ is a free abelian group and the rank is equal to $\delta(M)$.
\end{theorem}

Let us start with the following easy observation. 

\begin{lemma}\label{prop:uniqueness of bounded}
For any silting subcategory $\MM$ of $\TT$, we have 
\[\TT^{\rm{b}}=\{X\in \TT \ |\ \sum_{i\in \Bbb{Z}}{\rm{dim}}_{k}\homo{\TT}{Y}{X[i]}<\infty \ \mbox{for any } Y\in\MM\}.\]
\end{lemma}

\begin{proof}
Let $X\in \TT$ satisfy $\sum_{i\in \Bbb{Z}}{\rm{dim}}_{k}\homo{\TT}{Y}{X[i]}<\infty$ for any $Y\in\MM$. 
Let 
\[\mathcal{U}_X:=\{Y\in \TT^{\rm{c}}\ |\ \sum_{i\in \Bbb{Z}}{\rm{dim}}_{k}\homo{\TT}{Y}{X[i]}<\infty\}.\]
Then $\mathcal{U}_X$ is a thick subcategory of $\TT^{\rm{c}}$ containing $\MM$. 
Thus $\mathcal{U}_X\supset\thick\MM=\TT^{\rm{c}}$ holds, so we have $X\in \TT^{\rm{b}}$.
\end{proof}

Recall that we have a t-structure $(\TT^{\leq 0}_{M}, \TT^{\geq 0}_{M})$ of $\TT$ by Corollary \ref{HKM2}. Let
\[(\TT^{\rm{b}})^{\leq 0}_{M}:=\{X^{\leq 0}\ |\ X\in\TT^{\rm{b}}\}
\ \mbox{ and } (\TT^{\rm{b}})^{\geq 0}_{M}:=\{X^{\geq 0}\ |\ X\in\TT^{\rm{b}}\}.\]
We have the following crucial result. 

\begin{proposition}\label{prop:t-structure on bounded}
Let $M$ be a silting object in $\TT$. Then the following assertions hold.
\begin{itemize}
\item[(a)] $((\TT^{\rm{b}})^{\leq 0}_{M}, (\TT^{\rm{b}})^{\geq 0}_{M})$ is a t-structure on $\TT^{\rm{b}}$.
\item[(b)] We have an equivalence
\[\homo{\TT}{M}{-}:\ (\TT^{\rm{b}})^{0}_{T}\to \fl{\ehomo{\TT}{M}}\] 
\end{itemize}
\end{proposition}

\begin{proof}
(a) By Corollary \ref{HKM2}, for any $X\in \TT^{\rm{b}}$ we have a triangle
\[\xymatrix{
X^{\leq 0} \ar[r] & X \ar[r] & X^{>0} \ar[r] & X^{\leq 0}[1]
}\] 
with $X^{\leq 0}\in \TT_M^{\leq 0}$ and $X^{>0}\in \TT_M^{>0}$. 
Applying $\Hom_{\TT}(M,-)$, we have isomorphisms
\[\homo{\TT}{M}{X[i]}\simeq 
\begin{cases} 
 \ \homo{\TT}{M}{X^{\leq 0}[i]} & \mbox{for any } i\leq 0 \\
 \ \homo{\TT}{M}{X^{>0}[i]} & \mbox{for any } i>0
\end{cases}\]
By Lemma \ref{prop:uniqueness of bounded}, these imply $X^{\leq 0}, X^{>0}\in \TT^{\rm{b}}$. 
Hence $((\TT^{\rm{b}})^{\leq 0}_{M}, (\TT^{\rm{b}})^{\geq 0}_{M})$ is a t-structure on $\TT^{\rm{b}}$. 

(b) By Theorem \ref{realization1}, 
the functor $\homo{\TT}{M}{-}:\ (\TT^{\rm{b}})^{0}_{M}\to \fl{\ehomo{\TT}{M}}$ is fully faithful. 
Let $U\in \fl{\ehomo{\TT}{T}}$. 
By Theorem \ref{realization1}, there exists an object $X\in \TT^{0}_{T}$ such that $U\simeq \homo{\TT}{M}{X}$. 
Since $\ehomo{\TT}{M}/J_{\ehomo{\TT}{M}}$ is a finite dimensional $k$-algebra and $U$ has a finite length, 
we have ${\rm{dim}}_{k}\homo{\TT}{M}{X}<\infty$. 
By Lemma \ref{prop:uniqueness of bounded}, we obtain $X\in (\TT^{\rm{b}})^{0}_{M}$. 
Thus the functor $\homo{\TT}{M}{-}:\ (\TT^{\rm{b}})^{0}_{M}\to \fl{\ehomo{\TT}{M}}$ is dense. 
\end{proof}

Note that $\TT^{\rm{b}}$ is skeletally small by Proposition \ref{prop:t-structure on bounded}. 

\begin{definition}
Let $M$ be a silting object in $\TT$. 
For any non-zero $X\in \TT^{\rm{b}}$, we put 
\[a_{X}^{M}=\max\{i\in \Bbb{Z}\ |\ \homo{\TT}{M}{X[i]}\not=0\},\]
\[b_{X}^{M}=\min\{i\in \Bbb{Z}\ |\ \homo{\TT}{M}{X[i]}\not=0\}.\]
We define a positive integer, namely the \emph{length of $X$ with respect to $M$}, by $\ell_{X}^{M}=a_{X}^{M}-b_{X}^{M}+1$.
\end{definition}

We shall show the following statement. 

\begin{lemma}\label{lemma:Grothendieck group}
For any silting objects $M\in \TT$, we have $K_{0}(\TT^{\rm{b}})=K_{0}((\TT^{\rm{b}})^{0}_{M})$. 
\end{lemma}

\begin{proof}
It is clear that $K_{0}((\TT^{\rm{b}})^{0}_{M})$ is a subgroup of $K_{0}(\TT^{\rm{b}})$. 
Let $X\in \TT^{\rm{b}}$. We will prove by induction on the length of $X$ with respect to $M$. 
We can assume $a_{X}^{M}=0$. 
By Proposition \ref{prop:t-structure on bounded}, we have a triangle 
\[\xymatrix{
X^{<0} \ar[r] & X \ar[r] & X^{\geq 0} \ar[r] & X^{<0}[1] 
}\]
with $X^{<0}\in (\TT^{\rm{b}})^{<0}_{M}$ and $X^{\geq 0}\in (\TT^{\rm{b}})^{\geq 0}_{M}$. 
We can easily show $X^{\geq 0}\in (\TT^{\rm{b}})^{0}_{M}$ and $\homo{\TT}{M}{X^{<0}[i]}=0$ for $i<b_{X}^{M}, 0\leq i$. 
By induction hypothesis, we have $[X^{<0}]\in K_{0}((\TT^{\rm{b}})^{0}_{M})$. 
Hence we obtain $[X]=[X^{<0}]+[X^{\geq 0}]\in K_{0}((\TT^{\rm{b}})^{0}_{M})$. 
\end{proof}

Now we are ready to prove Theorem \ref{rank and delta}. 

Let $M$ be a silting object in $\TT$. 
By Proposition \ref{prop:t-structure on bounded} and Lemma \ref{lemma:Grothendieck group}, 
we have $K_0(\TT^{\rm{b}})\simeq K_0(\fl{\ehomo{\TT}{M}})$.
Since $\TT$ is Krull-Schmidt and $\ehomo{\TT}{M}/J_{\ehomo{\TT}{M}}$ is a finite dimensional $k$-algebra,
we have that the rank of $K_0(\fl{\ehomo{\TT}{M}})$ is equal to $\delta(M)$.
Thus the assertion follows.
\qed



%%%%%%%%%%%%%%%%%%%%%%%%%%%%%%%%%%%%%%%%%%%%%%%%%%%%%%%%%%%%%%%%%%%%%%%%%%%%%%%%%%%%%%%%%%%%%%%%%%%%%%%%%%%
%%%%%%%%%%%%%%%%%%%%%%%%%%%%%%%%%%%%%%%%%%%%%%%%%%%%%%%%%%%%%%%%%%%%%%%%%%%%%%%%%%%%%%%%%%%%%%%%%%%%%%%%%%%

\subsection{Properties of the hearts of t-structures}\label{Realization of hearts as functor categories}

The aim of this subsection is to study the hearts of certain t-structures including silting t-structures.
We shall show that the hearts have compact and generating subcategories in the sense below. Moreover we shall realize the hearts as functor categories defined below.

Let $\TT$ be a triangulated category with arbitrary coproducts. 
Let $\MM$ be a skeletally small subcategory of $\TT^{\rm c}$ satisfying $\Hom_{\TT}(\MM,\MM[>0])=0$ 
(We do not assume that $\MM$ is generating).
By Corollary \ref{t and co-t}(a), we have a t-structure $(\TT^{\le0},\TT^{\ge0})$ of $\TT$ given by
\[\TT^{\le0}:={}^\perp(\MM[\ge0]^\perp)\ \mbox{ and }\ \TT^{\ge0}:=\MM[>0]^\perp.\]
Thus for any $X\in\TT$, we have a triangle
\begin{equation}\label{t-structure triangle}
\xymatrix{X^{<0}\ar[r]^{b_X}&X\ar[r]^{a_X}&X^{\ge0}\ar[r]&X^{<0}[1]}
\end{equation}
with $X^{<0}\in\TT^{<0}$ and $X^{\ge0}\in\TT^{\ge0}$. We denote by
\[\AA:=\TT^{\le0}\cap\TT^{\ge0}\]
the heart of the t-structure.
Then $\AA$ is an abelian subcategory of $\TT$ closed under arbitrary coproducts.

Let $\AA$ be an abelian category with arbitrary coproducts in general.
We say that an object $X\in\AA$ is \emph{compact} if $\Hom_{\AA}(X,-)$ commutes with arbitrary coproducts.
We say that a subcategory $\PP$ of $\AA$ is \emph{generating} if for any object $A\in\AA$
there exists a surjection $\coprod P_i\to A$ with $P_i\in\PP$.

For a skeletally small category $\PP$, the functor category $\Mod\PP$ is defined as follows:
\begin{itemize}
\item An object is a contravariant additive functor from $\PP$ to the category of abelian groups.
\item A morphism is a natural transformation of functors.
\end{itemize}
The category $\Mod\PP$ is an abelian category with arbitrary products and coproducts.
In fact, kernels (respectively, cokernels, products, coproducts) are given by `objectwise' kernels (respectively, cokernels, products, coproducts).

Now we are ready to state our main result in this subsection.

\begin{theorem}\label{category P}
Under the above setting, the subcategory
\[\PP:=\{X^{\ge0}\ |\ X\in\MM\}\]
of $\TT$ satisfies the following properties.
\begin{itemize}
\item[(a)] $\PP$ is a skeletally small generating subcategory of $\AA$ consisting of compact projective objects in $\AA$.
\item[(b)] We have an equivalence $\AA\to\Mod\PP$ given by $A\mapsto\Hom_{\TT}(-,A)$.
\end{itemize}
\end{theorem}

As an application, we shall deduce the result below from Theorem \ref{category P}.
This is a generalization of a result of Hoshino-Kato-Miyachi \cite{HKM}, and their assumed that $\MM$ is generating is dropped.

\begin{corollary}\label{realization1}
Under the above setting, we have an equivalence $\AA\to\Mod\MM$ given by $A\mapsto\Hom_{\TT}(-,A)$.
\end{corollary}

The first step of the proof of Theorem \ref{category P} is the following.

\begin{lemma}
$\PP$ is contained in $\AA$.
\end{lemma}

\begin{proof}
For any $X\in\MM$, we shall show $X^{\ge0}\in\AA$.
We only have to show $X^{\ge0}\in\TT^{\le0}$. This follows from
\[X^{\ge0}\stackrel{\eqref{t-structure triangle}}{\in}X*X^{<0}[1]\subset\TT^{\le0}*\TT^{<-1}\subset\TT^{\le0}.\]
\end{proof}

\begin{lemma}\label{iso on A}
For any $X\in\TT$, we have an isomorphism $\Hom_{\TT}(X^{\ge0},\AA)\xrightarrow{\cdot a_X}\Hom_{\TT}(X,\AA)$.
\end{lemma}

\begin{proof}
Applying $\Hom_{\TT}(-,\AA)$ to \eqref{t-structure triangle}, we have an exact sequence
\[\Hom_{\TT}(X^{<0}[1],\AA)\to\Hom_{\TT}(X^{\ge0},\AA)\xrightarrow{\cdot a_X}
\Hom_{\TT}(X,\AA)\to\Hom_{\TT}(X^{<0},\AA).\]
Since $\Hom_{\TT}(X^{<0}[\ge0],\AA)\subset\Hom_{\TT}(\TT^{<0},\TT^{\ge0})=0$ holds, we have the desired isomorphism.
\end{proof}

\begin{lemma}\label{compact}
$\PP$ consists of compact objects in $\AA$.
\end{lemma}

\begin{proof}
The assertion follows from
\[\Hom_{\TT}(X^{\ge0},\coprod-)\stackrel{{\rm Lem.}\ \ref{iso on A}}{\simeq}
\Hom_{\TT}(X,\coprod-)\stackrel{X\in\TT^{\rm c}}{\simeq}
\coprod\Hom_{\TT}(X,-)\stackrel{{\rm Lem.}\ \ref{iso on A}}{\simeq}
\coprod\Hom_{\TT}(X^{\ge0},-).\]
\end{proof}

%\begin{lemma}\label{crucial}
%We have $\Hom_{\TT}(\MM,\TT^{<0})=0$.
%\end{lemma}
%
%\begin{proof}
%This was already shown.
%\end{proof}

\begin{lemma}\label{projective}
$\PP$ consists of projective objects in $\AA$.
\end{lemma}

\begin{proof}
Applying $\Hom_{\TT}(-,\AA[1])$ to \eqref{t-structure triangle}, we have an exact sequence
\[0\stackrel{\Hom_{\TT}(\TT^{<-1},\TT^{\ge-1})=0}{=}\Hom_{\TT}(X^{<0}[1],\AA[1])\to\Hom_{\TT}(X^{\ge0},\AA[1])\to\Hom_{\TT}(X,\AA[1]).\]
We have $\Hom_{\TT}(X,\AA[1])\subset\Hom_{\TT}(\MM,\TT^{<0})=0$ by Corollary \ref{HKM2}(a).
Thus $\Hom_{\TT}(X^{\ge0},\AA[1])=0$ holds, and $X^{\ge0}$ is projective in $\AA$.
\end{proof}

\begin{lemma}\label{shift}
We have $\TT^{\le0}\cap\MM^\perp\subset\TT^{<0}$.
\end{lemma}

\begin{proof}
Let $X\in\TT^{\le0}\cap\MM^\perp$.
% in the left hand side, we take a triangle
%\[X^{<0}\xrightarrow{g} X\xrightarrow{f} X^{\ge0}\to X^{<0}[1]\]
%with $X^{<0}\in\TT^{<0}$ and $X^{\ge0}\in\TT^{\ge0}$.
Applying $\Hom_{\TT}(\MM,-)$ to \eqref{t-structure triangle}, we have an exact sequence
\[0\stackrel{X\in\MM^\perp}{=}\Hom_{\TT}(\MM,X)\to\Hom_{\TT}(\MM,X^{\ge0})\to\Hom_{\TT}(\MM,X^{<0}[1])\]
with $\Hom_{\TT}(\MM,X^{<0}[1])\subset\Hom_{\TT}(\MM,\TT^{<0})=0$ by Corollary \ref{HKM2}(a).
Thus $\Hom_{\TT}(\MM,X^{\ge0})=0$ holds.
Thus we have $X^{\ge0}\in\MM[\ge0]^\perp=\TT^{>0}$, and so $a_X\in\Hom_{\TT}(X,X^{\ge0})\subset\Hom_{\TT}(\TT^{\le0},\TT^{>0})=0$.
Thus $b_X$ is a split epimorphism, and we have $X\in\TT^{<0}$.
\end{proof}

\begin{lemma}\label{generating}
$\PP$ is a generating subcategory of $\AA$.
\end{lemma}

\begin{proof}
Fix $A\in\AA$. Since $\MM$ is skeletally small, we can take a right $\Add\MM$-approximation $f:\coprod M_i\to A$ with $M_i\in\MM$.
Since $\Hom_{\TT}(M_i^{<0},A)=0$, our morphism $f$ factors through
$g:\coprod M_i^{\ge0}\to A$.
We shall show that $g$ is an epimorphism in $\AA$. Take a triangle in $\TT$:
\begin{equation}\label{surjection}
\xymatrix{B\ar[r]&\coprod M_i^{\ge0}\ar[r]^(0.65)g&A\ar[r]&B[1].}
\end{equation}
We have $B\in A[-1]*\coprod M_i^{\ge0}\subset\AA[-1]*\AA\subset\TT^{\le1}$.
Applying $\Hom_{\TT}(\MM,-)$ to \eqref{surjection}, we have an exact sequence
\[\Hom_{\TT}(\MM,\coprod M_i^{\ge0})\xrightarrow{g\cdot}\Hom_{\TT}(\MM,A)\to\Hom_{\TT}(\MM,B[1])\to\Hom_{\TT}(\MM,\coprod M_i^{\ge0}[1]).\]
The left map $(g\cdot)$ is surjective by our construction, and $\Hom_{\TT}(\MM,\coprod M_i^{\ge0}[1])\in\Hom_{\TT}(\MM,\TT^{<0})=0$ holds by Corollary \ref{HKM2}(a).
Thus $\Hom_{\TT}(\MM,B[1])=0$ holds, and we have $B\in\TT^{\le1}\cap\MM[-1]^\perp\subset\TT^{\le0}$ by Lemma \ref{shift}.
Applying $\Hom_{\TT}(-,\AA)$ to \eqref{surjection}, we have an exact sequence
\[0\stackrel{\Hom_{\TT}(\TT^{<0},\TT^{\ge0})=0}{=}\Hom_{\TT}(B[1],\AA)\to\Hom_{\TT}(A,\AA)\xrightarrow{g\cdot}\Hom_{\TT}(\coprod M_i^{\ge0},\AA).\]
Thus $g$ is an epimorphism in $\AA$.
\end{proof}

We have shown Theorem \ref{category P}(a).
Now Theorem \ref{category P}(b) is an immediate consequence of the following standard result in Morita theory for abelian categories.

\begin{proposition}\label{morita}
Let $\AA$ be an abelian category with arbitrary coproducts,
and $\PP$ a skeletally small generating subcategory of $\AA$ consisting of compact projective objects.
Then we have an equivalence $\AA\to\Mod\PP$ given by $A\mapsto\Hom_{\AA}(-,A)$.
\end{proposition}

\begin{proof}
For convenience of the reader, we give a complete proof here.
For simplicity reason, for each $A\in\AA$ we put
\[H_A:=\Hom_{\AA}(-,A):\PP\to\Mod\Z.\]
For any set $\{A_i\}$ of objects in $\AA$, a coproduct of $H_{A_i}$ is given by
\begin{equation}\label{coproduct compact}
\coprod_iH_{A_i}=H_{\coprod_iA_i}
\end{equation}
from our assumption that $\PP$ consists of compact objects in $\AA$ and the general description of coproducts in $\Mod\PP$.

(i) We shall show that our functor is fully faithful.

Fix two objects $A,B\in\AA$.
Since $\PP$ is generating, we can take an exact sequence
\begin{equation}\label{projective resolution of A}
\coprod_jQ_j\xrightarrow{}\coprod_iP_i\xrightarrow{}A\xrightarrow{}0
\end{equation}
with $P_i,Q_j\in\PP$.
Applying $\Hom_{\AA}(-,B)$, we have an exact sequence
\begin{equation}\label{Hom(-,B)}
0\xrightarrow{}\Hom_{\AA}(A,B)\xrightarrow{}\prod_i\Hom_{\AA}(P_i,B)
\xrightarrow{}\prod_j\Hom_{\AA}(Q_j,B).
\end{equation}
Applying $\Hom_{\AA}(\PP,-)$ to \eqref{projective resolution of A} and using \eqref{coproduct compact}, we have an exact sequence
\[\coprod_jH_{Q_j}\xrightarrow{}\coprod_iH_{P_i}\xrightarrow{}H_A\xrightarrow{}0\]
since $\PP$ consists of projective objects.
Applying $\Hom_{\Mod\PP}(-,H_B)$, we have an exact sequence
\begin{equation}\label{Hom((P,-),(P,B))}
0\xrightarrow{}\Hom_{\Mod\PP}(H_A,H_B)\xrightarrow{}\prod_i\Hom_{\Mod\PP}(H_{P_i},H_B)\xrightarrow{}\prod_j\Hom_{\Mod\PP}(H_{Q_j},H_B).
\end{equation}
Now \eqref{Hom(-,B)} and \eqref{Hom((P,-),(P,B))} form a commutative diagram
\[\xymatrix{
0\ar[r]&\Hom_{\AA}(A,B)\ar[r]\ar^{H_{A,B}}[d]&\prod_i\Hom_{\AA}(P_i,B)\ar[r]\ar^{\prod_iH_{P_i,B}}[d]&\prod_j\Hom_{\AA}(Q_j,B)\ar^{\prod_jH_{Q_j,B}}[d]\\
0\ar[r]&\Hom_{\Mod\PP}(H_A,H_B)\ar[r]&\prod_i\Hom_{\Mod\PP}(H_{P_i},H_B)\ar[r]&\prod_j\Hom_{\Mod\PP}(H_{Q_j},H_B).
}\]
By Yoneda's lemma, the right and the middle morphisms are isomorphisms. Thus the left morphism is also an isomorphism, and we have the assertion. 

(ii) We shall show that our functor is dense.

Since $\PP$ is skeletally small, we have that $\{H_P\ |\ P\in\PP\}$ is a generating subcategory of $\Mod\PP$.
For any $F\in\Mod\PP$, take an exact sequence
\begin{equation}\label{Hom(P,-)}
\coprod_jH_{Q_j}\xrightarrow{f}\coprod_iH_{P_i}\xrightarrow{}F\xrightarrow{}0.
\end{equation}
%Since we have
%\begin{eqnarray*}
%&&\Hom_{\Mod\PP}(\coprod_jH_{Q_j},\coprod_iH_{P_i})\simeq\prod_j\Hom_{\Mod\PP}(H_{Q_j},\coprod_iH_{P_i})\\
%&\stackrel{\mbox{Yoneda}}{\simeq}&\prod_j(\coprod_iH_{P_i})(Q_j)\stackrel{\eqref{coproduct compact}}{\simeq}\prod_j\Hom_{\AA}(Q_j,\coprod_iP_i)
%\simeq\Hom_{\AA}(\coprod_jQ_j,\coprod_iP_i),
%\end{eqnarray*}
Using Yoneda's Lemma, we have a morphism $g\in\Hom_{\AA}(\coprod_jQ_j,\coprod_iP_i)$ such that $f=H_g$.
Let $A\in\AA$ be a cokernel of $g$, so we have an exact sequence
\[\coprod_jQ_j\xrightarrow{g}\coprod_iP_i\xrightarrow{}A\xrightarrow{}0.\]
Applying $\Hom_{\AA}(\PP,-)$ and using \eqref{coproduct compact}, we have an exact sequence
\begin{equation}\label{Hom(P,-) 2}
\coprod_jH_{Q_j}\xrightarrow{f}\coprod_iH_{P_i}\xrightarrow{}H_A\xrightarrow{}0
\end{equation}
since $\PP$ consists of compact objects. Comparing \eqref{Hom(P,-)} and \eqref{Hom(P,-) 2}, we have $F\simeq H_A$.
Thus our functor is dense.
\end{proof}

In the rest we shall show Corollary \ref{realization1}.

\begin{lemma}\label{iso on M}
For any $X\in\TT$, we have an isomorphism $\Hom_{\TT}(\MM,X)\xrightarrow{\cdot a_X}\Hom_{\TT}(\MM,X^{\ge0})$.
\end{lemma}

\begin{proof}
Applying $\Hom_{\TT}(\MM,-)$ to \eqref{t-structure triangle}, we have an exact sequence
\[\Hom_{\TT}(\MM,X^{<0})\to\Hom_{\TT}(\MM,X)\xrightarrow{\cdot a_X}\Hom_{\TT}(\MM,X^{\ge0})\to\Hom_{\TT}(\MM,X^{<0}[1]).\]
Since $\Hom_{\TT}(\MM,X^{<0}[\ge0])\subset\Hom_{\TT}(\MM,\TT^{<0}[\ge0])=0$ holds by Corollary \ref{HKM2}(a), we have the assertion.
\end{proof}

\begin{lemma}\label{equivalence of M and P}
The functor $(-)^{\ge0}:\TT\to\TT^{\ge0}$ induces an equivalence $\MM\to\PP$.
\end{lemma}

\begin{proof}
We only have to show fully faithfulness.
For any $X,Y\in\MM$, we have isomorphisms
\[\xymatrix@R=.1cm@C=2cm{
\Hom_{\TT}(X,Y)\ar[r]_\sim^{{\rm Lem.}\ref{iso on M}}&\Hom_{\TT}(X,Y^{\ge0})&\Hom_{\TT}(X^{\ge0},Y^{\ge0})\ar[l]^\sim_{{\rm Lem.}\ref{iso on A}}\\
f\ar@{|->}[r]&a_Yf=f^{\ge0}a_X&f^{\ge0}.\ar@{|->}[l]
}\]
Thus the assertion follows.
\end{proof}

Now we are ready to prove Corollary \ref{realization1}.

By Lemma \ref{equivalence of M and P}, we have an equivalence $\Mod\PP\to\Mod\MM$, $F\mapsto F\circ(-)^{\ge0}$.
By Lemma \ref{iso on A}, we have a commutative diagram of functors
\[\xymatrix@C=.2cm@R=.1cm{
A\ar@{|->}[rrrrrr]&&&&&&\Hom_{\TT}(-,A)\\
\AA\ar^\sim[rrrrrr]\ar[ddrrrrrr]&&&&&&\Mod\PP\ar^\sim[dd]&F\ar@{|->}[dd]\\
A\ar@{|->}[ddrrrrrr]\\
&&&&&&\Mod\MM&F\circ(-)^{\ge0}\\
&&&&&&\Hom_{\TT}(-,A)}\]
where the horizontal functor is an equivalence by Theorem \ref{category P}.
Thus our functor $\AA\to\Mod\MM$ is also an equivalence.
\qed

